\Needspace{7\baselineskip}
\subsection{Row 1: 3.1 Root filter, coefficient action}
\label{row:1}
\textbf{Exact row scope.} Exact formal coefficient selector.\par\smallskip
\textbf{Manuscript labels.} \code{thm:root_filter}\par\smallskip
{\small\code{Komorebi.PaperIV.cyclicProjector}\par}
{\footnotesize\color{muted}definition; \code{KomorebiPaperIV/CyclicFilters.lean}, line 25.\par}\smallskip
{\small\code{Komorebi.PaperIV.coeff_cyclicProjector}\par}
{\footnotesize\color{muted}theorem; \code{KomorebiPaperIV/CyclicFilters.lean}, line 29.\par}\smallskip
{\small\code{Komorebi.PaperIV.cyclicProjector_eq_zero_of_residue_ne}\par}
{\footnotesize\color{muted}theorem; \code{KomorebiPaperIV/CyclicFilters.lean}, line 47.\par}\smallskip
{\small\code{Komorebi.PaperIV.cyclicProjector_eq_of_mod_eq}\par}
{\footnotesize\color{muted}theorem; \code{KomorebiPaperIV/CyclicFilters.lean}, line 52.\par}\smallskip
\Needspace{7\baselineskip}
\subsection{Row 2: 3.1 Root filter, projector resolution}
\label{row:2}
\textbf{Exact row scope.} Exact over every commutative semiring; completeness assumes \code{N > 0}.\par\smallskip
\textbf{Manuscript labels.} \code{thm:root_filter}\par\smallskip
{\small\code{Komorebi.PaperIV.cyclicProjector_idempotent}\par}
{\footnotesize\color{muted}theorem; \code{KomorebiPaperIV/CyclicFilters.lean}, line 34.\par}\smallskip
{\small\code{Komorebi.PaperIV.cyclicProjector_orthogonal}\par}
{\footnotesize\color{muted}theorem; \code{KomorebiPaperIV/CyclicFilters.lean}, line 39.\par}\smallskip
{\small\code{Komorebi.PaperIV.sum_cyclicProjector}\par}
{\footnotesize\color{muted}theorem; \code{KomorebiPaperIV/CyclicFilters.lean}, line 74.\par}\smallskip
\Needspace{7\baselineskip}
\subsection{Row 3: 3.3 Cyclic product law}
\label{row:3}
\textbf{Exact row scope.} Exact pure-packet label addition, off-label vanishing, and full finite double-sum product formula.\par\smallskip
\textbf{Manuscript labels.} \code{thm:cyclic_product}\par\smallskip
{\small\code{Komorebi.PaperIV.cyclicProjector_add}\par}
{\footnotesize\color{muted}theorem; \code{KomorebiPaperIV/CyclicFilters.lean}, line 58.\par}\smallskip
{\small\code{Komorebi.PaperIV.cyclicProjector_sum}\par}
{\footnotesize\color{muted}theorem; \code{KomorebiPaperIV/CyclicFilters.lean}, line 63.\par}\smallskip
{\small\code{Komorebi.PaperIV.cyclicProjector_product_packet}\par}
{\footnotesize\color{muted}theorem; \code{KomorebiPaperIV/CyclicFilters.lean}, line 88.\par}\smallskip
{\small\code{Komorebi.PaperIV.cyclicProjector_product_other}\par}
{\footnotesize\color{muted}theorem; \code{KomorebiPaperIV/CyclicFilters.lean}, line 114.\par}\smallskip
{\small\code{Komorebi.PaperIV.cyclic_product_law}\par}
{\footnotesize\color{muted}theorem; \code{KomorebiPaperIV/CyclicFilters.lean}, line 122.\par}\smallskip
\Needspace{7\baselineskip}
\subsection{Row 4: 3.1 Root filter, explicit Fourier average}
\label{row:4}
\textbf{Exact row scope.} Exact primitive-root finite Fourier average in \code{ℂ⟦X⟧}, proved equal to the coefficient projector.\par\smallskip
\textbf{Manuscript labels.} \code{thm:root_filter}\par\smallskip
{\small\code{Komorebi.PaperIV.finiteRootAverage}\par}
{\footnotesize\color{muted}definition; \code{KomorebiPaperIV/FiniteAnalytic.lean}, line 30.\par}\smallskip
{\small\code{Komorebi.PaperIV.coeff_finiteRootAverage}\par}
{\footnotesize\color{muted}theorem; \code{KomorebiPaperIV/FiniteAnalytic.lean}, line 69.\par}\smallskip
{\small\code{Komorebi.PaperIV.finiteRootAverage_eq_cyclicProjector}\par}
{\footnotesize\color{muted}theorem; \code{KomorebiPaperIV/FiniteAnalytic.lean}, line 83.\par}\smallskip
\Needspace{7\baselineskip}
\subsection{Row 5: 4.1 Transfer exponential basis}
\label{row:5}
\textbf{Exact row scope.} Exact formal realization of the lacunary exponential series.\par\smallskip
\textbf{Manuscript labels.} \code{def:transfer_exponential}\par\smallskip
{\small\code{Komorebi.PaperIV.exponentialGeneratingSeries}\par}
{\footnotesize\color{muted}definition; \code{KomorebiPaperIV/SectionCycles.lean}, line 26.\par}\smallskip
{\small\code{Komorebi.PaperIV.coeff_exponentialGeneratingSeries}\par}
{\footnotesize\color{muted}theorem; \code{KomorebiPaperIV/SectionCycles.lean}, line 31.\par}\smallskip
{\small\code{Komorebi.PaperIV.transferSectionCoeff}\par}
{\footnotesize\color{muted}definition; \code{KomorebiPaperIV/SectionCycles.lean}, line 68.\par}\smallskip
{\small\code{Komorebi.PaperIV.transferSection}\par}
{\footnotesize\color{muted}definition; \code{KomorebiPaperIV/SectionCycles.lean}, line 72.\par}\smallskip
{\small\code{Komorebi.PaperIV.coeff_transferSection}\par}
{\footnotesize\color{muted}theorem; \code{KomorebiPaperIV/SectionCycles.lean}, line 76.\par}\smallskip
\Needspace{7\baselineskip}
\subsection{Row 6: 4.1 Transfer exponential, root-average representation}
\label{row:6}
\textbf{Exact row scope.} Exact finite root-of-unity representation in \code{ℂ⟦X⟧}, including rescaling by the spectral parameter.\par\smallskip
\textbf{Manuscript labels.} \code{def:transfer_exponential}\par\smallskip
{\small\code{Komorebi.PaperIV.finiteRootAverage_exp_eq_transferSection}\par}
{\footnotesize\color{muted}theorem; \code{KomorebiPaperIV/FiniteAnalytic.lean}, line 90.\par}\smallskip
{\small\code{Komorebi.PaperIV.finiteRootAverage_rescale_exp_eq_transferSection}\par}
{\footnotesize\color{muted}theorem; \code{KomorebiPaperIV/FiniteAnalytic.lean}, line 104.\par}\smallskip
\Needspace{7\baselineskip}
\subsection{Row 7: 4.2 Differential cycle and initial data}
\label{row:7}
\textbf{Exact row scope.} Exact formal derivative cycle, full-period equation, and initial data. Analytic solution-space dimension and the Wronskian claim are outside this row.\par\smallskip
\textbf{Manuscript labels.} \code{thm:exponential_cycle}\par\smallskip
{\small\code{Komorebi.PaperIV.derivative_exponentialGeneratingSeries}\par}
{\footnotesize\color{muted}theorem; \code{KomorebiPaperIV/SectionCycles.lean}, line 37.\par}\smallskip
{\small\code{Komorebi.PaperIV.iterate_derivative_exponentialGeneratingSeries}\par}
{\footnotesize\color{muted}theorem; \code{KomorebiPaperIV/SectionCycles.lean}, line 55.\par}\smallskip
{\small\code{Komorebi.PaperIV.transferSectionCoeff_succ_of_pos}\par}
{\footnotesize\color{muted}theorem; \code{KomorebiPaperIV/SectionCycles.lean}, line 84.\par}\smallskip
{\small\code{Komorebi.PaperIV.transferSectionCoeff_succ_zero}\par}
{\footnotesize\color{muted}theorem; \code{KomorebiPaperIV/SectionCycles.lean}, line 113.\par}\smallskip
{\small\code{Komorebi.PaperIV.derivative_transferSection_of_pos}\par}
{\footnotesize\color{muted}theorem; \code{KomorebiPaperIV/SectionCycles.lean}, line 225.\par}\smallskip
{\small\code{Komorebi.PaperIV.derivative_transferSection_zero}\par}
{\footnotesize\color{muted}theorem; \code{KomorebiPaperIV/SectionCycles.lean}, line 237.\par}\smallskip
{\small\code{Komorebi.PaperIV.transferSectionCoeff_add_period}\par}
{\footnotesize\color{muted}theorem; \code{KomorebiPaperIV/SectionCycles.lean}, line 269.\par}\smallskip
{\small\code{Komorebi.PaperIV.iterate_derivative_transferSection}\par}
{\footnotesize\color{muted}theorem; \code{KomorebiPaperIV/SectionCycles.lean}, line 298.\par}\smallskip
{\small\code{Komorebi.PaperIV.constantCoeff_iterate_derivative_transferSection}\par}
{\footnotesize\color{muted}theorem; \code{KomorebiPaperIV/SectionCycles.lean}, line 320.\par}\smallskip
\Needspace{7\baselineskip}
\subsection{Row 8: 5.1 Signed cyclic sections}
\label{row:8}
\textbf{Exact row scope.} Exact formal realization of the signed lacunary series.\par\smallskip
\textbf{Manuscript labels.} \code{def:signed_sections}\par\smallskip
{\small\code{Komorebi.PaperIV.signedSectionCoeff}\par}
{\footnotesize\color{muted}definition; \code{KomorebiPaperIV/SectionCycles.lean}, line 137.\par}\smallskip
{\small\code{Komorebi.PaperIV.signedSection}\par}
{\footnotesize\color{muted}definition; \code{KomorebiPaperIV/SectionCycles.lean}, line 141.\par}\smallskip
{\small\code{Komorebi.PaperIV.coeff_signedSection}\par}
{\footnotesize\color{muted}theorem; \code{KomorebiPaperIV/SectionCycles.lean}, line 145.\par}\smallskip
\Needspace{7\baselineskip}
\subsection{Row 9: 5.2 Relation with the transfer exponential}
\label{row:9}
\textbf{Exact row scope.} Exact formal identity for a chosen \code{rho} with \code{rho ^ N = -1}, including the factor \code{(rho ^ j)⁻¹}.\par\smallskip
\textbf{Manuscript labels.} \code{prop:signed_exponential_relation}\par\smallskip
{\small\code{Komorebi.PaperIV.signedSection_eq_inv_pow_mul_transferSection}\par}
{\footnotesize\color{muted}theorem; \code{KomorebiPaperIV/FiniteAnalytic.lean}, line 143.\par}\smallskip
\Needspace{7\baselineskip}
\subsection{Row 10: 5.3 Signed derivative cycle}
\label{row:10}
\textbf{Exact row scope.} Exact formal derivative cycle, negative wraparound, full-period equation, and initial data.\par\smallskip
\textbf{Manuscript labels.} \code{thm:signed_cycle}\par\smallskip
{\small\code{Komorebi.PaperIV.signedSectionCoeff_succ_of_pos}\par}
{\footnotesize\color{muted}theorem; \code{KomorebiPaperIV/SectionCycles.lean}, line 170.\par}\smallskip
{\small\code{Komorebi.PaperIV.signedSectionCoeff_succ_zero}\par}
{\footnotesize\color{muted}theorem; \code{KomorebiPaperIV/SectionCycles.lean}, line 200.\par}\smallskip
{\small\code{Komorebi.PaperIV.derivative_signedSection_of_pos}\par}
{\footnotesize\color{muted}theorem; \code{KomorebiPaperIV/SectionCycles.lean}, line 248.\par}\smallskip
{\small\code{Komorebi.PaperIV.derivative_signedSection_zero}\par}
{\footnotesize\color{muted}theorem; \code{KomorebiPaperIV/SectionCycles.lean}, line 258.\par}\smallskip
{\small\code{Komorebi.PaperIV.signedSectionCoeff_add_period}\par}
{\footnotesize\color{muted}theorem; \code{KomorebiPaperIV/SectionCycles.lean}, line 284.\par}\smallskip
{\small\code{Komorebi.PaperIV.iterate_derivative_signedSection}\par}
{\footnotesize\color{muted}theorem; \code{KomorebiPaperIV/SectionCycles.lean}, line 310.\par}\smallskip
{\small\code{Komorebi.PaperIV.constantCoeff_iterate_derivative_signedSection}\par}
{\footnotesize\color{muted}theorem; \code{KomorebiPaperIV/SectionCycles.lean}, line 337.\par}\smallskip
\Needspace{7\baselineskip}
\subsection{Row 11: 5.6 Finite exponential representation}
\label{row:11}
\textbf{Exact row scope.} Exact finite representation in \code{ℂ⟦X⟧} for every signed section and the distinguished denominator. Pointwise evaluation is supplied by row 47.\par\smallskip
\textbf{Manuscript labels.} \code{thm:finite_exponential_representation}\par\smallskip
{\small\code{Komorebi.PaperIV.signedSection_eq_finiteRootAverage_rescale_exp}\par}
{\footnotesize\color{muted}theorem; \code{KomorebiPaperIV/FiniteAnalytic.lean}, line 153.\par}\smallskip
{\small\code{Komorebi.PaperIV.signedSection_zero_eq_finite_exponential_average}\par}
{\footnotesize\color{muted}theorem; \code{KomorebiPaperIV/FiniteAnalytic.lean}, line 166.\par}\smallskip
\Needspace{7\baselineskip}
\subsection{Row 12: 5.7 and 7.1 Reduced denominator and reciprocal family}
\label{row:12}
\textbf{Exact row scope.} Exact named formal specialization \code{Delta_N(w)^(-r) = exp(-r log Delta_N(w))}, with parameters \code{-1}, \code{0}, and \code{1} linked to the denominator, unit, and inverse. No analytic logarithm branch is asserted.\par\smallskip
\textbf{Manuscript labels.} \code{def:reciprocal_family}, \code{def:reduced_denominator}\par\smallskip
{\small\code{Komorebi.PaperIV.formalExpFamily}\par}
{\footnotesize\color{muted}definition; \code{KomorebiPaperIV/FormalExponential.lean}, line 27.\par}\smallskip
{\small\code{Komorebi.PaperIV.formalExpFamily_zero}\par}
{\footnotesize\color{muted}theorem; \code{KomorebiPaperIV/FormalExponential.lean}, line 37.\par}\smallskip
{\small\code{Komorebi.PaperIV.formalExpFamily_one}\par}
{\footnotesize\color{muted}theorem; \code{KomorebiPaperIV/FormalExponential.lean}, line 45.\par}\smallskip
{\small\code{Komorebi.PaperIV.reducedDenominator}\par}
{\footnotesize\color{muted}definition; \code{KomorebiPaperIV/ReciprocalCore.lean}, line 28.\par}\smallskip
{\small\code{Komorebi.PaperIV.coeff_reducedDenominator}\par}
{\footnotesize\color{muted}theorem; \code{KomorebiPaperIV/ReciprocalCore.lean}, line 33.\par}\smallskip
{\small\code{Komorebi.PaperIV.constantCoeff_reducedDenominator}\par}
{\footnotesize\color{muted}theorem; \code{KomorebiPaperIV/ReciprocalCore.lean}, line 39.\par}\smallskip
{\small\code{Komorebi.PaperIV.signedSection_zero_eq_expand_reducedDenominator}\par}
{\footnotesize\color{muted}theorem; \code{KomorebiPaperIV/ReciprocalCore.lean}, line 46.\par}\smallskip
{\small\code{Komorebi.PaperIV.negativeLogDenominator}\par}
{\footnotesize\color{muted}definition; \code{KomorebiPaperIV/ReciprocalCore.lean}, line 63.\par}\smallskip
{\small\code{Komorebi.PaperIV.constantCoeff_negativeLogDenominator}\par}
{\footnotesize\color{muted}theorem; \code{KomorebiPaperIV/ReciprocalCore.lean}, line 67.\par}\smallskip
{\small\code{Komorebi.PaperIV.exp_subst_logOf}\par}
{\footnotesize\color{muted}theorem; \code{KomorebiPaperIV/ReciprocalCore.lean}, line 74.\par}\smallskip
{\small\code{Komorebi.PaperIV.mul_derivative_logOf}\par}
{\footnotesize\color{muted}theorem; \code{KomorebiPaperIV/ReciprocalCore.lean}, line 88.\par}\smallskip
{\small\code{Komorebi.PaperIV.formalExpFamily_eq_exp_subst}\par}
{\footnotesize\color{muted}theorem; \code{KomorebiPaperIV/ReciprocalCore.lean}, line 101.\par}\smallskip
{\small\code{Komorebi.PaperIV.reciprocalSeries}\par}
{\footnotesize\color{muted}definition; \code{KomorebiPaperIV/ReciprocalCore.lean}, line 112.\par}\smallskip
{\small\code{Komorebi.PaperIV.constantCoeff_reciprocalSeries}\par}
{\footnotesize\color{muted}theorem; \code{KomorebiPaperIV/ReciprocalCore.lean}, line 116.\par}\smallskip
{\small\code{Komorebi.PaperIV.reciprocalSeries_zero}\par}
{\footnotesize\color{muted}theorem; \code{KomorebiPaperIV/ReciprocalCore.lean}, line 128.\par}\smallskip
{\small\code{Komorebi.PaperIV.reciprocalSeries_neg_one}\par}
{\footnotesize\color{muted}theorem; \code{KomorebiPaperIV/ReciprocalCore.lean}, line 133.\par}\smallskip
{\small\code{Komorebi.PaperIV.reciprocalSeries_one_eq_inverse}\par}
{\footnotesize\color{muted}theorem; \code{KomorebiPaperIV/ReciprocalCore.lean}, line 151.\par}\smallskip
{\small\code{Komorebi.PaperIV.reciprocalCoeff}\par}
{\footnotesize\color{muted}definition; \code{KomorebiPaperIV/ReciprocalCore.lean}, line 161.\par}\smallskip
{\small\code{Komorebi.PaperIV.reciprocalOrdCoeff}\par}
{\footnotesize\color{muted}definition; \code{KomorebiPaperIV/ReciprocalCore.lean}, line 165.\par}\smallskip
{\small\code{Komorebi.PaperIV.reciprocalCoeff_zero_parameter}\par}
{\footnotesize\color{muted}theorem; \code{KomorebiPaperIV/ReciprocalCore.lean}, line 169.\par}\smallskip
{\small\code{Komorebi.PaperIV.reciprocalCoeff_zero_parameter_of_pos}\par}
{\footnotesize\color{muted}theorem; \code{KomorebiPaperIV/ReciprocalCore.lean}, line 173.\par}\smallskip
\Needspace{7\baselineskip}
\subsection{Row 13: 7.4 Additive parameter law for reciprocal powers}
\label{row:13}
\textbf{Exact row scope.} Exact formal parameter law and exact \code{(Nn choose Nq)} convolution for the named \code{P_{N,n}} specialization.\par\smallskip
\textbf{Manuscript labels.} \code{thm:reciprocal_parameter_addition}\par\smallskip
{\small\code{Komorebi.PaperIV.formalExpFamily_semigroup}\par}
{\footnotesize\color{muted}theorem; \code{KomorebiPaperIV/FormalExponential.lean}, line 30.\par}\smallskip
{\small\code{Komorebi.PaperIV.blockCoeff}\par}
{\footnotesize\color{muted}definition; \code{KomorebiPaperIV/FormalExponential.lean}, line 129.\par}\smallskip
{\small\code{Komorebi.PaperIV.blockCoeff_mul}\par}
{\footnotesize\color{muted}theorem; \code{KomorebiPaperIV/FormalExponential.lean}, line 134.\par}\smallskip
{\small\code{Komorebi.PaperIV.reciprocal_block_convolution}\par}
{\footnotesize\color{muted}theorem; \code{KomorebiPaperIV/FormalExponential.lean}, line 163.\par}\smallskip
{\small\code{Komorebi.PaperIV.reciprocalSeries_add}\par}
{\footnotesize\color{muted}theorem; \code{KomorebiPaperIV/ReciprocalCore.lean}, line 142.\par}\smallskip
{\small\code{Komorebi.PaperIV.reciprocalCoeff_addition}\par}
{\footnotesize\color{muted}theorem; \code{KomorebiPaperIV/ReciprocalCore.lean}, line 178.\par}\smallskip
\Needspace{7\baselineskip}
\subsection{Row 14: 7.5 Binomial-type normalization}
\label{row:14}
\textbf{Exact row scope.} Exact exponential-generating normalization and binomial convolution.\par\smallskip
\textbf{Manuscript labels.} \code{cor:binomial_type}\par\smallskip
{\small\code{Komorebi.PaperIV.egfCoeff}\par}
{\footnotesize\color{muted}definition; \code{KomorebiPaperIV/FormalExponential.lean}, line 101.\par}\smallskip
{\small\code{Komorebi.PaperIV.egfCoeff_mul}\par}
{\footnotesize\color{muted}theorem; \code{KomorebiPaperIV/FormalExponential.lean}, line 106.\par}\smallskip
{\small\code{Komorebi.PaperIV.reciprocal_binomial_type}\par}
{\footnotesize\color{muted}theorem; \code{KomorebiPaperIV/FormalExponential.lean}, line 174.\par}\smallskip
\Needspace{7\baselineskip}
\subsection{Row 15: 7.6, 8.2, 8.3, and 8.4 Newton, tangent, and reciprocal derivative identities}
\label{row:15}
\textbf{Exact row scope.} Exact formal logarithmic derivative and finite coefficient recurrences in the reduced variable. The displayed \code{z}-variable tangent scaling is not separately named.\par\smallskip
\textbf{Manuscript labels.} \code{cor:reciprocal_derivative}, \code{prop:tangent_kappa}, \code{thm:logarithmic_newton_recurrence}, \code{thm:tangent_recurrence}\par\smallskip
{\small\code{Komorebi.PaperIV.derivative_rescale_exp}\par}
{\footnotesize\color{muted}theorem; \code{KomorebiPaperIV/FormalExponential.lean}, line 50.\par}\smallskip
{\small\code{Komorebi.PaperIV.derivative_formalExpFamily}\par}
{\footnotesize\color{muted}theorem; \code{KomorebiPaperIV/FormalExponential.lean}, line 75.\par}\smallskip
{\small\code{Komorebi.PaperIV.formalExpFamily_coefficient_recurrence}\par}
{\footnotesize\color{muted}theorem; \code{KomorebiPaperIV/FormalExponential.lean}, line 84.\par}\smallskip
{\small\code{Komorebi.PaperIV.kappaCoeff}\par}
{\footnotesize\color{muted}definition; \code{KomorebiPaperIV/ReciprocalCore.lean}, line 191.\par}\smallskip
{\small\code{Komorebi.PaperIV.reducedTangentSeries}\par}
{\footnotesize\color{muted}definition; \code{KomorebiPaperIV/ReciprocalCore.lean}, line 195.\par}\smallskip
{\small\code{Komorebi.PaperIV.coeff_reducedTangentSeries}\par}
{\footnotesize\color{muted}theorem; \code{KomorebiPaperIV/ReciprocalCore.lean}, line 199.\par}\smallskip
{\small\code{Komorebi.PaperIV.reducedDenominator_mul_reducedTangentSeries}\par}
{\footnotesize\color{muted}theorem; \code{KomorebiPaperIV/ReciprocalCore.lean}, line 207.\par}\smallskip
{\small\code{Komorebi.PaperIV.reducedTangent_coefficient_convolution}\par}
{\footnotesize\color{muted}theorem; \code{KomorebiPaperIV/ReciprocalCore.lean}, line 217.\par}\smallskip
{\small\code{Komorebi.PaperIV.derivative_reciprocalSeries}\par}
{\footnotesize\color{muted}theorem; \code{KomorebiPaperIV/ReciprocalCore.lean}, line 228.\par}\smallskip
{\small\code{Komorebi.PaperIV.reciprocalSeries_coefficient_recurrence}\par}
{\footnotesize\color{muted}theorem; \code{KomorebiPaperIV/ReciprocalCore.lean}, line 235.\par}\smallskip
{\small\code{Komorebi.PaperIV.reciprocalOrdCoeff_newton}\par}
{\footnotesize\color{muted}theorem; \code{KomorebiPaperIV/ReciprocalCore.lean}, line 245.\par}\smallskip
\Needspace{7\baselineskip}
\subsection{Row 16: 9.2 Generic inverse support}
\label{row:16}
\textbf{Exact row scope.} Exact inverse convolution for every unit-constant series over a field.\par\smallskip
\textbf{Manuscript labels.} \code{thm:triangular_inverse}\par\smallskip
{\small\code{Komorebi.PaperIV.inverse_coefficient_convolution}\par}
{\footnotesize\color{muted}theorem; \code{KomorebiPaperIV/FormalExponential.lean}, line 188.\par}\smallskip
{\small\code{Komorebi.PaperIV.constantCoeff_inverse_of_one}\par}
{\footnotesize\color{muted}theorem; \code{KomorebiPaperIV/FormalExponential.lean}, line 198.\par}\smallskip
\Needspace{7\baselineskip}
\subsection{Row 17: 9.1 and 9.2 Normalized residue section and triangular inverse}
\label{row:17}
\textbf{Exact row scope.} Exact manuscript normalization, factorial-rescaled triangular identity, and solved recurrence over a characteristic-zero field.\par\smallskip
\textbf{Manuscript labels.} \code{def:normalized_section}, \code{thm:triangular_inverse}\par\smallskip
{\small\code{Komorebi.PaperIV.normalizedResidueSection}\par}
{\footnotesize\color{muted}definition; \code{KomorebiPaperIV/AlgebraicCompletion.lean}, line 35.\par}\smallskip
{\small\code{Komorebi.PaperIV.coeff_normalizedResidueSection}\par}
{\footnotesize\color{muted}theorem; \code{KomorebiPaperIV/AlgebraicCompletion.lean}, line 41.\par}\smallskip
{\small\code{Komorebi.PaperIV.constantCoeff_normalizedResidueSection}\par}
{\footnotesize\color{muted}theorem; \code{KomorebiPaperIV/AlgebraicCompletion.lean}, line 47.\par}\smallskip
{\small\code{Komorebi.PaperIV.residueInverseCoeff}\par}
{\footnotesize\color{muted}definition; \code{KomorebiPaperIV/AlgebraicCompletion.lean}, line 56.\par}\smallskip
{\small\code{Komorebi.PaperIV.residueInverseCoeff_triangular}\par}
{\footnotesize\color{muted}theorem; \code{KomorebiPaperIV/AlgebraicCompletion.lean}, line 61.\par}\smallskip
{\small\code{Komorebi.PaperIV.residueInverseCoeff_zero}\par}
{\footnotesize\color{muted}theorem; \code{KomorebiPaperIV/AlgebraicCompletion.lean}, line 126.\par}\smallskip
{\small\code{Komorebi.PaperIV.residueInverseCoeff_recurrence}\par}
{\footnotesize\color{muted}theorem; \code{KomorebiPaperIV/AlgebraicCompletion.lean}, line 131.\par}\smallskip
\Needspace{7\baselineskip}
\subsection{Row 18: 9.5, 9.6, and 9.7 Appell realization and exact sparse polynomial}
\label{row:18}
\textbf{Exact row scope.} Exact finite Appell realization, lowering and translation laws, and the manuscript's sparse sum with only degrees \code{Nn}. The residue specialization here is \code{r = 1}.\par\smallskip
\textbf{Manuscript labels.} \code{def:block_appell}, \code{prop:explicit_appell}, \code{thm:appell_laws}\par\smallskip
{\small\code{Komorebi.PaperIV.appellSeries}\par}
{\footnotesize\color{muted}definition; \code{KomorebiPaperIV/Appell.lean}, line 26.\par}\smallskip
{\small\code{Komorebi.PaperIV.appellCoeff}\par}
{\footnotesize\color{muted}definition; \code{KomorebiPaperIV/Appell.lean}, line 30.\par}\smallskip
{\small\code{Komorebi.PaperIV.egfCoeff_rescale_exp}\par}
{\footnotesize\color{muted}theorem; \code{KomorebiPaperIV/Appell.lean}, line 58.\par}\smallskip
{\small\code{Komorebi.PaperIV.appellCoeff_zero_parameter}\par}
{\footnotesize\color{muted}theorem; \code{KomorebiPaperIV/Appell.lean}, line 73.\par}\smallskip
{\small\code{Komorebi.PaperIV.appell_addition}\par}
{\footnotesize\color{muted}theorem; \code{KomorebiPaperIV/Appell.lean}, line 80.\par}\smallskip
{\small\code{Komorebi.PaperIV.appellCoeff_at_zero}\par}
{\footnotesize\color{muted}theorem; \code{KomorebiPaperIV/Appell.lean}, line 89.\par}\smallskip
{\small\code{Komorebi.PaperIV.appellPolynomial}\par}
{\footnotesize\color{muted}definition; \code{KomorebiPaperIV/Appell.lean}, line 97.\par}\smallskip
{\small\code{Komorebi.PaperIV.derivative_appellPolynomial}\par}
{\footnotesize\color{muted}theorem; \code{KomorebiPaperIV/Appell.lean}, line 103.\par}\smallskip
{\small\code{Komorebi.PaperIV.formalAppellPolynomial}\par}
{\footnotesize\color{muted}definition; \code{KomorebiPaperIV/Appell.lean}, line 133.\par}\smallskip
{\small\code{Komorebi.PaperIV.derivative_formalAppellPolynomial}\par}
{\footnotesize\color{muted}theorem; \code{KomorebiPaperIV/Appell.lean}, line 137.\par}\smallskip
{\small\code{Komorebi.PaperIV.eval_formalAppellPolynomial}\par}
{\footnotesize\color{muted}theorem; \code{KomorebiPaperIV/Appell.lean}, line 142.\par}\smallskip
{\small\code{Komorebi.PaperIV.blockCoefficientSequence}\par}
{\footnotesize\color{muted}definition; \code{KomorebiPaperIV/AlgebraicCompletion.lean}, line 155.\par}\smallskip
{\small\code{Komorebi.PaperIV.blockCoefficientSequence_mul}\par}
{\footnotesize\color{muted}theorem; \code{KomorebiPaperIV/AlgebraicCompletion.lean}, line 159.\par}\smallskip
{\small\code{Komorebi.PaperIV.appellPolynomial_eq_sparse_sum}\par}
{\footnotesize\color{muted}theorem; \code{KomorebiPaperIV/AlgebraicCompletion.lean}, line 166.\par}\smallskip
{\small\code{Komorebi.PaperIV.appellPolynomial_eq_block_sum}\par}
{\footnotesize\color{muted}theorem; \code{KomorebiPaperIV/AlgebraicCompletion.lean}, line 185.\par}\smallskip
{\small\code{Komorebi.PaperIV.residueAppellPolynomial_explicit}\par}
{\footnotesize\color{muted}theorem; \code{KomorebiPaperIV/AlgebraicCompletion.lean}, line 282.\par}\smallskip
\Needspace{7\baselineskip}
\subsection{Row 19: 9.8 Joint Appell convolution}
\label{row:19}
\textbf{Exact row scope.} Exact formal product and coefficient convolution in space and order parameters for the generic zero-constant logarithm family.\par\smallskip
\textbf{Manuscript labels.} \code{thm:appell_two_parameter_convolution}\par\smallskip
{\small\code{Komorebi.PaperIV.appellSeries_two_parameter}\par}
{\footnotesize\color{muted}theorem; \code{KomorebiPaperIV/Appell.lean}, line 34.\par}\smallskip
{\small\code{Komorebi.PaperIV.appell_two_parameter_convolution}\par}
{\footnotesize\color{muted}theorem; \code{KomorebiPaperIV/Appell.lean}, line 48.\par}\smallskip
\Needspace{7\baselineskip}
\subsection{Row 20: 9.9 Formal differential-operator representation}
\label{row:20}
\textbf{Exact row scope.} Exact finite action on monomials and exact operator inversion on every polynomial of bounded degree; the direct residue inverse corollary is the \code{r = 1} case.\par\smallskip
\textbf{Manuscript labels.} \code{prop:appell_operator_form}\par\smallskip
{\small\code{Komorebi.PaperIV.egfConstantCoeffAction}\par}
{\footnotesize\color{muted}definition; \code{KomorebiPaperIV/AlgebraicCompletion.lean}, line 325.\par}\smallskip
{\small\code{Komorebi.PaperIV.egfConstantCoeffAction_X_pow}\par}
{\footnotesize\color{muted}theorem; \code{KomorebiPaperIV/AlgebraicCompletion.lean}, line 331.\par}\smallskip
{\small\code{Komorebi.PaperIV.egfConstantCoeffAction_residue_X_pow}\par}
{\footnotesize\color{muted}theorem; \code{KomorebiPaperIV/AlgebraicCompletion.lean}, line 345.\par}\smallskip
{\small\code{Komorebi.PaperIV.polynomialDifferentialEval}\par}
{\footnotesize\color{muted}definition; \code{KomorebiPaperIV/AlgebraicCompletion.lean}, line 363.\par}\smallskip
{\small\code{Komorebi.PaperIV.polynomialDifferentialEval_X}\par}
{\footnotesize\color{muted}theorem; \code{KomorebiPaperIV/AlgebraicCompletion.lean}, line 370.\par}\smallskip
{\small\code{Komorebi.PaperIV.truncatedDifferentialOperator}\par}
{\footnotesize\color{muted}definition; \code{KomorebiPaperIV/AlgebraicCompletion.lean}, line 375.\par}\smallskip
{\small\code{Komorebi.PaperIV.coeff_trunc_mul_trunc_of_lt}\par}
{\footnotesize\color{muted}theorem; \code{KomorebiPaperIV/AlgebraicCompletion.lean}, line 379.\par}\smallskip
{\small\code{Komorebi.PaperIV.X_pow_dvd_trunc_mul_trunc_sub_one}\par}
{\footnotesize\color{muted}theorem; \code{KomorebiPaperIV/AlgebraicCompletion.lean}, line 392.\par}\smallskip
{\small\code{Komorebi.PaperIV.truncatedDifferentialOperator_mul_apply_of_natDegree_le}\par}
{\footnotesize\color{muted}theorem; \code{KomorebiPaperIV/AlgebraicCompletion.lean}, line 405.\par}\smallskip
{\small\code{Komorebi.PaperIV.truncatedDifferentialOperator_normalizedResidueSection_inverse}\par}
{\footnotesize\color{muted}theorem; \code{KomorebiPaperIV/AlgebraicCompletion.lean}, line 439.\par}\smallskip
\Needspace{7\baselineskip}
\subsection{Row 21: 9.10 Appell monicity and covariance}
\label{row:21}
\textbf{Exact row scope.} Exact monicity, exact degree, and \code{omega ^ m} covariance for block-supported Appell polynomials; residue specialization is \code{r = 1}.\par\smallskip
\textbf{Manuscript labels.} \code{cor:appell_symmetry}\par\smallskip
{\small\code{Komorebi.PaperIV.appellPolynomial_isMonicOfDegree}\par}
{\footnotesize\color{muted}theorem; \code{KomorebiPaperIV/AlgebraicCompletion.lean}, line 218.\par}\smallskip
{\small\code{Komorebi.PaperIV.blockAppellPolynomial_isMonicOfDegree}\par}
{\footnotesize\color{muted}theorem; \code{KomorebiPaperIV/AlgebraicCompletion.lean}, line 239.\par}\smallskip
{\small\code{Komorebi.PaperIV.eval_appellPolynomial_block_covariance}\par}
{\footnotesize\color{muted}theorem; \code{KomorebiPaperIV/AlgebraicCompletion.lean}, line 248.\par}\smallskip
{\small\code{Komorebi.PaperIV.residueAppellPolynomial_isMonicOfDegree}\par}
{\footnotesize\color{muted}theorem; \code{KomorebiPaperIV/AlgebraicCompletion.lean}, line 293.\par}\smallskip
{\small\code{Komorebi.PaperIV.eval_residueAppellPolynomial_covariance}\par}
{\footnotesize\color{muted}theorem; \code{KomorebiPaperIV/AlgebraicCompletion.lean}, line 303.\par}\smallskip
\Needspace{7\baselineskip}
\subsection{Row 22: 10.1 and 10.2 Theta definition and reconstruction}
\label{row:22}
\textbf{Exact row scope.} Exact formal theta series in the reduced variable and exact signed reconstruction from the named reciprocal coefficients.\par\smallskip
\textbf{Manuscript labels.} \code{def:unsigned_denominator}, \code{thm:theta_reconstruction}\par\smallskip
{\small\code{Komorebi.PaperIV.blockCoeff_rescale}\par}
{\footnotesize\color{muted}theorem; \code{KomorebiPaperIV/ThetaCore.lean}, line 26.\par}\smallskip
{\small\code{Komorebi.PaperIV.thetaSeries}\par}
{\footnotesize\color{muted}definition; \code{KomorebiPaperIV/ThetaCore.lean}, line 37.\par}\smallskip
{\small\code{Komorebi.PaperIV.thetaCoeff}\par}
{\footnotesize\color{muted}definition; \code{KomorebiPaperIV/ThetaCore.lean}, line 42.\par}\smallskip
{\small\code{Komorebi.PaperIV.thetaCoeff_reconstruction_unfactored}\par}
{\footnotesize\color{muted}theorem; \code{KomorebiPaperIV/ThetaCore.lean}, line 63.\par}\smallskip
{\small\code{Komorebi.PaperIV.thetaCoeff_reconstruction}\par}
{\footnotesize\color{muted}theorem; \code{KomorebiPaperIV/ThetaCore.lean}, line 77.\par}\smallskip
\Needspace{7\baselineskip}
\subsection{Row 23: 10.3 Theta trivial parameter factors}
\label{row:23}
\textbf{Exact row scope.} Exact vanishing at \code{u = 0}, and at \code{u = -1} for even \code{N}; these are the root statements underlying the polynomial divisibility claim.\par\smallskip
\textbf{Manuscript labels.} \code{prop:theta_factors}\par\smallskip
{\small\code{Komorebi.PaperIV.thetaSeries_zero}\par}
{\footnotesize\color{muted}theorem; \code{KomorebiPaperIV/ThetaCore.lean}, line 46.\par}\smallskip
{\small\code{Komorebi.PaperIV.thetaCoeff_zero_parameter}\par}
{\footnotesize\color{muted}theorem; \code{KomorebiPaperIV/ThetaCore.lean}, line 51.\par}\smallskip
{\small\code{Komorebi.PaperIV.thetaCoeff_zero_parameter_of_pos}\par}
{\footnotesize\color{muted}theorem; \code{KomorebiPaperIV/ThetaCore.lean}, line 56.\par}\smallskip
{\small\code{Komorebi.PaperIV.thetaSeries_neg_one_of_even}\par}
{\footnotesize\color{muted}theorem; \code{KomorebiPaperIV/ThetaCore.lean}, line 106.\par}\smallskip
{\small\code{Komorebi.PaperIV.thetaCoeff_neg_one_of_even_of_pos}\par}
{\footnotesize\color{muted}theorem; \code{KomorebiPaperIV/ThetaCore.lean}, line 112.\par}\smallskip
\Needspace{7\baselineskip}
\subsection{Row 24: 3.2 Covariance characterization}
\label{row:24}
\textbf{Exact row scope.} Exact formal fixed-space and range/covariance characterizations for a primitive root. Holomorphic-domain and pointwise versions remain outside this row.\par\smallskip
\textbf{Manuscript labels.} \code{prop:covariance}\par\smallskip
{\small\code{Komorebi.PaperIV.cyclicProjector_fixed_iff_rescale_covariant}\par}
{\footnotesize\color{muted}theorem; \code{KomorebiPaperIV/EarlyCompletion.lean}, line 35.\par}\smallskip
{\small\code{Komorebi.PaperIV.cyclicProjector_mem_range_iff_rescale_covariant}\par}
{\footnotesize\color{muted}theorem; \code{KomorebiPaperIV/EarlyCompletion.lean}, line 74.\par}\smallskip
\Needspace{7\baselineskip}
\subsection{Row 25: 4.4 Bivariate formal-series setup}
\label{row:25}
\textbf{Exact row scope.} Exact coefficient encodings for substitution into a sum of two formal variables and for separated products. A bridge to a canonical multivariate substitution API is not asserted.\par\smallskip
\textbf{Manuscript labels.} \code{cor:exponential_addition}\par\smallskip
{\small\code{Komorebi.PaperIV.bivariateExponentialGeneratingSeries}\par}
{\footnotesize\color{muted}definition; \code{KomorebiPaperIV/EarlyCompletion.lean}, line 90.\par}\smallskip
{\small\code{Komorebi.PaperIV.coeff_bivariateExponentialGeneratingSeries}\par}
{\footnotesize\color{muted}theorem; \code{KomorebiPaperIV/EarlyCompletion.lean}, line 97.\par}\smallskip
{\small\code{Komorebi.PaperIV.separatedExponentialGeneratingSeries}\par}
{\footnotesize\color{muted}definition; \code{KomorebiPaperIV/EarlyCompletion.lean}, line 105.\par}\smallskip
{\small\code{Komorebi.PaperIV.coeff_separatedExponentialGeneratingSeries}\par}
{\footnotesize\color{muted}theorem; \code{KomorebiPaperIV/EarlyCompletion.lean}, line 112.\par}\smallskip
\Needspace{7\baselineskip}
\subsection{Row 26: 4.4 Addition law}
\label{row:26}
\textbf{Exact row scope.} Exact coefficient-encoded bivariate formal addition law. Pointwise analytic evaluation is outside this row.\par\smallskip
\textbf{Manuscript labels.} \code{cor:exponential_addition}\par\smallskip
{\small\code{Komorebi.PaperIV.transferSection_addition}\par}
{\footnotesize\color{muted}theorem; \code{KomorebiPaperIV/EarlyCompletion.lean}, line 144.\par}\smallskip
\Needspace{7\baselineskip}
\subsection{Row 27: 4.5 Hypergeometric form}
\label{row:27}
\textbf{Exact row scope.} Exact rational factorial/Pochhammer normalization and formal-power-series identity. The parameter is rational; the explicit erased parameter multiset, complex analytic summation, and Mathlib regularized-hypergeometric bridge remain outside this row.\par\smallskip
\textbf{Manuscript labels.} \code{thm:hypergeometric_form}\par\smallskip
{\small\code{Komorebi.PaperIV.residuePochhammerProduct}\par}
{\footnotesize\color{muted}definition; \code{KomorebiPaperIV/HypergeometricCompletion.lean}, line 36.\par}\smallskip
{\small\code{Komorebi.PaperIV.factorial_eq_residuePochhammerProduct}\par}
{\footnotesize\color{muted}theorem; \code{KomorebiPaperIV/HypergeometricCompletion.lean}, line 61.\par}\smallskip
{\small\code{Komorebi.PaperIV.residueZeroFCoeff}\par}
{\footnotesize\color{muted}definition; \code{KomorebiPaperIV/HypergeometricCompletion.lean}, line 97.\par}\smallskip
{\small\code{Komorebi.PaperIV.residueZeroFSeries}\par}
{\footnotesize\color{muted}definition; \code{KomorebiPaperIV/HypergeometricCompletion.lean}, line 101.\par}\smallskip
{\small\code{Komorebi.PaperIV.coeff_residueZeroFSeries}\par}
{\footnotesize\color{muted}theorem; \code{KomorebiPaperIV/HypergeometricCompletion.lean}, line 105.\par}\smallskip
{\small\code{Komorebi.PaperIV.residueZeroFCoeff_normalization}\par}
{\footnotesize\color{muted}theorem; \code{KomorebiPaperIV/HypergeometricCompletion.lean}, line 113.\par}\smallskip
{\small\code{Komorebi.PaperIV.residueParameter_eq_one}\par}
{\footnotesize\color{muted}theorem; \code{KomorebiPaperIV/HypergeometricCompletion.lean}, line 123.\par}\smallskip
{\small\code{Komorebi.PaperIV.omitted_residuePochhammer_eq_factorial}\par}
{\footnotesize\color{muted}theorem; \code{KomorebiPaperIV/HypergeometricCompletion.lean}, line 130.\par}\smallskip
{\small\code{Komorebi.PaperIV.transferSection_eq_residueZeroFSeries}\par}
{\footnotesize\color{muted}theorem; \code{KomorebiPaperIV/HypergeometricCompletion.lean}, line 145.\par}\smallskip
\Needspace{7\baselineskip}
\subsection{Row 28: 5.5 Twisted addition law}
\label{row:28}
\textbf{Exact row scope.} Exact coefficient-encoded bivariate formal law, including the carry sign \code{(i+j)/N}. Canonical substitution and pointwise evaluation remain outside this row.\par\smallskip
\textbf{Manuscript labels.} \code{prop:signed_addition}\par\smallskip
{\small\code{Komorebi.PaperIV.signedSection_addition}\par}
{\footnotesize\color{muted}theorem; \code{KomorebiPaperIV/EarlyCompletion.lean}, line 217.\par}\smallskip
\Needspace{7\baselineskip}
\subsection{Row 29: 5.8 Denominator symmetries}
\label{row:29}
\textbf{Exact row scope.} Exact formal rotational invariance and conjugation-fixed coefficients. Pointwise conjugation, zero orbits, and multiplicity preservation remain outside this row.\par\smallskip
\textbf{Manuscript labels.} \code{cor:denominator_symmetries}\par\smallskip
{\small\code{Komorebi.PaperIV.signedDenominator_rescale_of_pow_eq_one}\par}
{\footnotesize\color{muted}theorem; \code{KomorebiPaperIV/EarlyCompletion.lean}, line 247.\par}\smallskip
{\small\code{Komorebi.PaperIV.signedDenominator_map_conj}\par}
{\footnotesize\color{muted}theorem; \code{KomorebiPaperIV/EarlyCompletion.lean}, line 263.\par}\smallskip
\Needspace{7\baselineskip}
\subsection{Row 30: 6.1 Finite cyclic algebras}
\label{row:30}
\textbf{Exact row scope.} Exact \code{AdjoinRoot} quotient algebras, generator relations, power reductions, and \code{N}-dimensional power bases for \code{N > 0}.\par\smallskip
\textbf{Manuscript labels.} \code{def:cyclic_algebras}\par\smallskip
{\small\code{Komorebi.PaperIV.cyclicAlgebraPolynomial}\par}
{\footnotesize\color{muted}definition; \code{KomorebiPaperIV/FiniteModels.lean}, line 33.\par}\smallskip
{\small\code{Komorebi.PaperIV.signedCyclicAlgebraPolynomial}\par}
{\footnotesize\color{muted}definition; \code{KomorebiPaperIV/FiniteModels.lean}, line 37.\par}\smallskip
{\small\code{Komorebi.PaperIV.CyclicAlgebra}\par}
{\footnotesize\color{muted}definition; \code{KomorebiPaperIV/FiniteModels.lean}, line 41.\par}\smallskip
{\small\code{Komorebi.PaperIV.SignedCyclicAlgebra}\par}
{\footnotesize\color{muted}definition; \code{KomorebiPaperIV/FiniteModels.lean}, line 45.\par}\smallskip
{\small\code{Komorebi.PaperIV.cyclicGenerator}\par}
{\footnotesize\color{muted}definition; \code{KomorebiPaperIV/FiniteModels.lean}, line 49.\par}\smallskip
{\small\code{Komorebi.PaperIV.signedCyclicGenerator}\par}
{\footnotesize\color{muted}definition; \code{KomorebiPaperIV/FiniteModels.lean}, line 53.\par}\smallskip
{\small\code{Komorebi.PaperIV.cyclicGenerator_pow}\par}
{\footnotesize\color{muted}theorem; \code{KomorebiPaperIV/FiniteModels.lean}, line 57.\par}\smallskip
{\small\code{Komorebi.PaperIV.signedCyclicGenerator_pow}\par}
{\footnotesize\color{muted}theorem; \code{KomorebiPaperIV/FiniteModels.lean}, line 65.\par}\smallskip
{\small\code{Komorebi.PaperIV.cyclicGenerator_pow_eq_mod}\par}
{\footnotesize\color{muted}theorem; \code{KomorebiPaperIV/FiniteModels.lean}, line 73.\par}\smallskip
{\small\code{Komorebi.PaperIV.signedCyclicGenerator_pow_eq_mod}\par}
{\footnotesize\color{muted}theorem; \code{KomorebiPaperIV/FiniteModels.lean}, line 86.\par}\smallskip
{\small\code{Komorebi.PaperIV.cyclicAlgebraPowerBasis}\par}
{\footnotesize\color{muted}definition; \code{KomorebiPaperIV/FiniteModels.lean}, line 104.\par}\smallskip
{\small\code{Komorebi.PaperIV.signedCyclicAlgebraPowerBasis}\par}
{\footnotesize\color{muted}definition; \code{KomorebiPaperIV/FiniteModels.lean}, line 112.\par}\smallskip
{\small\code{Komorebi.PaperIV.cyclicAlgebraPowerBasis_dim}\par}
{\footnotesize\color{muted}theorem; \code{KomorebiPaperIV/FiniteModels.lean}, line 119.\par}\smallskip
{\small\code{Komorebi.PaperIV.signedCyclicAlgebraPowerBasis_dim}\par}
{\footnotesize\color{muted}theorem; \code{KomorebiPaperIV/FiniteModels.lean}, line 127.\par}\smallskip
\Needspace{7\baselineskip}
\subsection{Row 31: 6.2 Finite algebra exponential identities}
\label{row:31}
\textbf{Exact row scope.} Exact identities of formal power series in the quotient algebras. Pointwise quotient-algebra convergence/evaluation and the packaged derivation of prior differential laws are outside this row.\par\smallskip
\textbf{Manuscript labels.} \code{thm:finite_algebra_exponentials}\par\smallskip
{\small\code{Komorebi.PaperIV.cyclicAlgebraExponentialSeries}\par}
{\footnotesize\color{muted}definition; \code{KomorebiPaperIV/FiniteModels.lean}, line 138.\par}\smallskip
{\small\code{Komorebi.PaperIV.signedCyclicAlgebraExponentialSeries}\par}
{\footnotesize\color{muted}definition; \code{KomorebiPaperIV/FiniteModels.lean}, line 145.\par}\smallskip
{\small\code{Komorebi.PaperIV.cyclicAlgebra_exponential_eq_sections}\par}
{\footnotesize\color{muted}theorem; \code{KomorebiPaperIV/FiniteModels.lean}, line 152.\par}\smallskip
{\small\code{Komorebi.PaperIV.signedCyclicAlgebra_exponential_eq_sections}\par}
{\footnotesize\color{muted}theorem; \code{KomorebiPaperIV/FiniteModels.lean}, line 178.\par}\smallskip
\Needspace{7\baselineskip}
\subsection{Row 32: 6.3 Signed shift matrix}
\label{row:32}
\textbf{Exact row scope.} Exact closed signed-shift construction and additive power semigroup for \code{N > 0}.\par\smallskip
\textbf{Manuscript labels.} \code{def:signed_shift_matrix}\par\smallskip
{\small\code{Komorebi.PaperIV.signedShiftPowerMatrix}\par}
{\footnotesize\color{muted}definition; \code{KomorebiPaperIV/FiniteModels.lean}, line 213.\par}\smallskip
{\small\code{Komorebi.PaperIV.signedShiftMatrix}\par}
{\footnotesize\color{muted}definition; \code{KomorebiPaperIV/FiniteModels.lean}, line 221.\par}\smallskip
{\small\code{Komorebi.PaperIV.signedShiftPowerMatrix_apply}\par}
{\footnotesize\color{muted}theorem; \code{KomorebiPaperIV/FiniteModels.lean}, line 225.\par}\smallskip
{\small\code{Komorebi.PaperIV.signedShiftPowerMatrix_zero}\par}
{\footnotesize\color{muted}theorem; \code{KomorebiPaperIV/FiniteModels.lean}, line 232.\par}\smallskip
{\small\code{Komorebi.PaperIV.signedShiftPowerMatrix_mul}\par}
{\footnotesize\color{muted}theorem; \code{KomorebiPaperIV/FiniteModels.lean}, line 258.\par}\smallskip
{\small\code{Komorebi.PaperIV.signedShiftMatrix_pow}\par}
{\footnotesize\color{muted}theorem; \code{KomorebiPaperIV/FiniteModels.lean}, line 297.\par}\smallskip
\Needspace{7\baselineskip}
\subsection{Row 33: 6.4 and 6.5 Companion realization and path entries}
\label{row:33}
\textbf{Exact row scope.} Exact formula for every entry of every matrix power and exact \code{J_N^N = -I}, plus the matrix-exponential construction. The exponential-section entry bridge, first-column and trace formulas, characteristic polynomial, and simple spectrum remain outside this row.\par\smallskip
\textbf{Manuscript labels.} \code{cor:fundamental_matrix_entries}, \code{thm:companion_matrix_realization}\par\smallskip
{\small\code{Komorebi.PaperIV.signedShiftMatrix_pow_apply}\par}
{\footnotesize\color{muted}theorem; \code{KomorebiPaperIV/FiniteModels.lean}, line 305.\par}\smallskip
{\small\code{Komorebi.PaperIV.signedShiftMatrix_pow_order}\par}
{\footnotesize\color{muted}theorem; \code{KomorebiPaperIV/FiniteModels.lean}, line 328.\par}\smallskip
{\small\code{Komorebi.PaperIV.signedFundamentalMatrix}\par}
{\footnotesize\color{muted}definition; \code{KomorebiPaperIV/FiniteModels.lean}, line 337.\par}\smallskip
\Needspace{7\baselineskip}
\subsection{Row 34: 6.6 Companion addition and inversion}
\label{row:34}
\textbf{Exact row scope.} Exact matrix-exponential addition, inverse, cancellation, and unit laws. The determinant-one and signed-section entry clauses remain outside this row.\par\smallskip
\textbf{Manuscript labels.} \code{cor:companion_addition_inversion}\par\smallskip
{\small\code{Komorebi.PaperIV.signedFundamentalMatrix_add}\par}
{\footnotesize\color{muted}theorem; \code{KomorebiPaperIV/FiniteModels.lean}, line 341.\par}\smallskip
{\small\code{Komorebi.PaperIV.signedFundamentalMatrix_neg}\par}
{\footnotesize\color{muted}theorem; \code{KomorebiPaperIV/FiniteModels.lean}, line 349.\par}\smallskip
{\small\code{Komorebi.PaperIV.signedFundamentalMatrix_mul_neg}\par}
{\footnotesize\color{muted}theorem; \code{KomorebiPaperIV/FiniteModels.lean}, line 355.\par}\smallskip
{\small\code{Komorebi.PaperIV.signedFundamentalMatrix_isUnit}\par}
{\footnotesize\color{muted}theorem; \code{KomorebiPaperIV/FiniteModels.lean}, line 361.\par}\smallskip
\Needspace{7\baselineskip}
\subsection{Row 35: 7.3 Parameter recurrence}
\label{row:35}
\textbf{Exact row scope.} Exact formal differential equation, factorial/binomial coefficient equation, solved recurrence, and initial value; the second sum uses the explicit reindexing \code{k=q+1}.\par\smallskip
\textbf{Manuscript labels.} \code{thm:parameter_recurrence}\par\smallskip
{\small\code{Komorebi.PaperIV.reducedDenominator_reciprocalSeries_differentialEquation}\par}
{\footnotesize\color{muted}theorem; \code{KomorebiPaperIV/ReciprocalCompletion.lean}, line 32.\par}\smallskip
{\small\code{Komorebi.PaperIV.reciprocalOrdCoeff_parameter_equation}\par}
{\footnotesize\color{muted}theorem; \code{KomorebiPaperIV/ReciprocalCompletion.lean}, line 58.\par}\smallskip
{\small\code{Komorebi.PaperIV.reciprocalCoeff_parameter_equation}\par}
{\footnotesize\color{muted}theorem; \code{KomorebiPaperIV/ReciprocalCompletion.lean}, line 198.\par}\smallskip
{\small\code{Komorebi.PaperIV.reciprocalCoeff_parameter_recurrence}\par}
{\footnotesize\color{muted}theorem; \code{KomorebiPaperIV/ReciprocalCompletion.lean}, line 317.\par}\smallskip
{\small\code{Komorebi.PaperIV.reciprocalCoeff_block_zero}\par}
{\footnotesize\color{muted}theorem; \code{KomorebiPaperIV/ReciprocalCompletion.lean}, line 346.\par}\smallskip
\Needspace{7\baselineskip}
\subsection{Row 36: 7.7 First three reciprocal polynomials}
\label{row:36}
\textbf{Exact row scope.} Exact \code{P_(N,1)=r} and exact recurrence-normal forms for \code{P_(N,2)} and \code{P_(N,3)}. The manuscript's compact factorial rewrites are not separately proved.\par\smallskip
\textbf{Manuscript labels.} \code{cor:first_reciprocal_polynomials}\par\smallskip
{\small\code{Komorebi.PaperIV.reciprocalCoeff_one}\par}
{\footnotesize\color{muted}theorem; \code{KomorebiPaperIV/ReciprocalCompletion.lean}, line 351.\par}\smallskip
{\small\code{Komorebi.PaperIV.reciprocalCoeff_two_recurrence_form}\par}
{\footnotesize\color{muted}theorem; \code{KomorebiPaperIV/ReciprocalCompletion.lean}, line 360.\par}\smallskip
{\small\code{Komorebi.PaperIV.reciprocalCoeff_three_recurrence_form}\par}
{\footnotesize\color{muted}theorem; \code{KomorebiPaperIV/ReciprocalCompletion.lean}, line 376.\par}\smallskip
\Needspace{7\baselineskip}
\subsection{Row 37: 7.8 Divisibility by the parameter}
\label{row:37}
\textbf{Exact row scope.} Exact divisibility by \code{X} for the universal polynomial over \code{ℚ}, plus a separately defined generic-field polynomial with exact scalar evaluation. A named base-change theorem connecting those two polynomial lifts, or generic-field divisibility itself, remains outside this row.\par\smallskip
\textbf{Manuscript labels.} \code{cor:factor_r}\par\smallskip
{\small\code{Komorebi.PaperIV.reciprocalCoeffPolynomial}\par}
{\footnotesize\color{muted}definition; \code{KomorebiPaperIV/ReciprocalCompletion.lean}, line 397.\par}\smallskip
{\small\code{Komorebi.PaperIV.X_dvd_reciprocalCoeffPolynomial}\par}
{\footnotesize\color{muted}theorem; \code{KomorebiPaperIV/ReciprocalCompletion.lean}, line 402.\par}\smallskip
{\small\code{Komorebi.PaperIV.coeff_zero_reciprocalCoeffPolynomial}\par}
{\footnotesize\color{muted}theorem; \code{KomorebiPaperIV/ReciprocalCompletion.lean}, line 430.\par}\smallskip
{\small\code{Komorebi.PaperIV.reciprocalCoeffPolynomialOver}\par}
{\footnotesize\color{muted}definition; \code{KomorebiPaperIV/ParameterThetaCompletion.lean}, line 718.\par}\smallskip
{\small\code{Komorebi.PaperIV.eval_reciprocalCoeffPolynomialOver}\par}
{\footnotesize\color{muted}theorem; \code{KomorebiPaperIV/ParameterThetaCompletion.lean}, line 721.\par}\smallskip
\Needspace{7\baselineskip}
\subsection{Row 38: 8.1, 8.2, and 8.3 Transfer tangent and recurrence}
\label{row:38}
\textbf{Exact row scope.} Exact formal quotient definition and exact factorial/binomial recurrence for the kappa-scaled sequence. The bridge identifying these named coefficients with the full-variable tangent series remains outside this row.\par\smallskip
\textbf{Manuscript labels.} \code{def:transfer_tangent}, \code{prop:tangent_kappa}, \code{thm:tangent_recurrence}\par\smallskip
{\small\code{Komorebi.PaperIV.formalTransferTangent}\par}
{\footnotesize\color{muted}definition; \code{KomorebiPaperIV/ReciprocalCompletion.lean}, line 579.\par}\smallskip
{\small\code{Komorebi.PaperIV.transferTangentCoeff}\par}
{\footnotesize\color{muted}definition; \code{KomorebiPaperIV/ReciprocalCompletion.lean}, line 444.\par}\smallskip
{\small\code{Komorebi.PaperIV.transferTangentCoeff_eq_kappa}\par}
{\footnotesize\color{muted}theorem; \code{KomorebiPaperIV/ReciprocalCompletion.lean}, line 448.\par}\smallskip
{\small\code{Komorebi.PaperIV.transferTangentCoeff_convolution}\par}
{\footnotesize\color{muted}theorem; \code{KomorebiPaperIV/ReciprocalCompletion.lean}, line 510.\par}\smallskip
{\small\code{Komorebi.PaperIV.transferTangentCoeff_recurrence}\par}
{\footnotesize\color{muted}theorem; \code{KomorebiPaperIV/ReciprocalCompletion.lean}, line 538.\par}\smallskip
{\small\code{Komorebi.PaperIV.transferTangentCoeff_zero}\par}
{\footnotesize\color{muted}theorem; \code{KomorebiPaperIV/ReciprocalCompletion.lean}, line 566.\par}\smallskip
\Needspace{7\baselineskip}
\subsection{Row 39: 8.5 Scalar higher Riccati identity}
\label{row:39}
\textbf{Exact row scope.} Exact nontrivial formal-germ differential-Bell/Riccati identity. Pointwise analytic validity on the entire zero-free set remains outside this row.\par\smallskip
\textbf{Manuscript labels.} \code{thm:higher_riccati}\par\smallskip
{\small\code{Komorebi.PaperIV.differentialBellSeries}\par}
{\footnotesize\color{muted}definition; \code{KomorebiPaperIV/ReciprocalCompletion.lean}, line 586.\par}\smallskip
{\small\code{Komorebi.PaperIV.differentialBellSeries_zero}\par}
{\footnotesize\color{muted}theorem; \code{KomorebiPaperIV/ReciprocalCompletion.lean}, line 594.\par}\smallskip
{\small\code{Komorebi.PaperIV.differentialBellSeries_succ}\par}
{\footnotesize\color{muted}theorem; \code{KomorebiPaperIV/ReciprocalCompletion.lean}, line 599.\par}\smallskip
{\small\code{Komorebi.PaperIV.iterate_derivative_eq_differentialBellSeries_mul}\par}
{\footnotesize\color{muted}theorem; \code{KomorebiPaperIV/ReciprocalCompletion.lean}, line 607.\par}\smallskip
{\small\code{Komorebi.PaperIV.derivative_signedSection_zero_eq_neg_formalTransferTangent_mul}\par}
{\footnotesize\color{muted}theorem; \code{KomorebiPaperIV/ReciprocalCompletion.lean}, line 621.\par}\smallskip
{\small\code{Komorebi.PaperIV.differentialBellSeries_neg_formalTransferTangent_eq_neg_one}\par}
{\footnotesize\color{muted}theorem; \code{KomorebiPaperIV/ReciprocalCompletion.lean}, line 640.\par}\smallskip
\Needspace{7\baselineskip}
\subsection{Row 40: 9.3 Arbitrary-order residue-section powers}
\label{row:40}
\textbf{Exact row scope.} Exact branch-free formal family \code{Phi_(N,j)^(-r)}, its factorially normalized coefficients, and the \code{r=1} bridge to the triangular inverse. No analytic logarithm branch is asserted.\par\smallskip
\textbf{Manuscript labels.} \code{def:residue_section_power}\par\smallskip
{\small\code{Komorebi.PaperIV.normalizedResidueSectionAlg}\par}
{\footnotesize\color{muted}definition; \code{KomorebiPaperIV/ParameterThetaCompletion.lean}, line 33.\par}\smallskip
{\small\code{Komorebi.PaperIV.coeff_normalizedResidueSectionAlg}\par}
{\footnotesize\color{muted}theorem; \code{KomorebiPaperIV/ParameterThetaCompletion.lean}, line 38.\par}\smallskip
{\small\code{Komorebi.PaperIV.constantCoeff_normalizedResidueSectionAlg}\par}
{\footnotesize\color{muted}theorem; \code{KomorebiPaperIV/ParameterThetaCompletion.lean}, line 45.\par}\smallskip
{\small\code{Komorebi.PaperIV.normalizedResidueSectionAlg_eq_normalizedResidueSection}\par}
{\footnotesize\color{muted}theorem; \code{KomorebiPaperIV/ParameterThetaCompletion.lean}, line 53.\par}\smallskip
{\small\code{Komorebi.PaperIV.normalizedResidueNegativeLog}\par}
{\footnotesize\color{muted}definition; \code{KomorebiPaperIV/ParameterThetaCompletion.lean}, line 68.\par}\smallskip
{\small\code{Komorebi.PaperIV.constantCoeff_normalizedResidueNegativeLog}\par}
{\footnotesize\color{muted}theorem; \code{KomorebiPaperIV/ParameterThetaCompletion.lean}, line 72.\par}\smallskip
{\small\code{Komorebi.PaperIV.residuePowerSeries}\par}
{\footnotesize\color{muted}definition; \code{KomorebiPaperIV/ParameterThetaCompletion.lean}, line 141.\par}\smallskip
{\small\code{Komorebi.PaperIV.constantCoeff_residuePowerSeries}\par}
{\footnotesize\color{muted}theorem; \code{KomorebiPaperIV/ParameterThetaCompletion.lean}, line 145.\par}\smallskip
{\small\code{Komorebi.PaperIV.residuePowerSeries_zero}\par}
{\footnotesize\color{muted}theorem; \code{KomorebiPaperIV/ParameterThetaCompletion.lean}, line 157.\par}\smallskip
{\small\code{Komorebi.PaperIV.residuePowerSeries_neg_one}\par}
{\footnotesize\color{muted}theorem; \code{KomorebiPaperIV/ParameterThetaCompletion.lean}, line 161.\par}\smallskip
{\small\code{Komorebi.PaperIV.residuePowerSeries_one_eq_inverse}\par}
{\footnotesize\color{muted}theorem; \code{KomorebiPaperIV/ParameterThetaCompletion.lean}, line 191.\par}\smallskip
{\small\code{Komorebi.PaperIV.residuePowerCoeff}\par}
{\footnotesize\color{muted}definition; \code{KomorebiPaperIV/ParameterThetaCompletion.lean}, line 202.\par}\smallskip
{\small\code{Komorebi.PaperIV.residuePowerCoeff_zero}\par}
{\footnotesize\color{muted}theorem; \code{KomorebiPaperIV/ParameterThetaCompletion.lean}, line 206.\par}\smallskip
{\small\code{Komorebi.PaperIV.residuePowerCoeff_one_eq_residueInverseCoeff}\par}
{\footnotesize\color{muted}theorem; \code{KomorebiPaperIV/ParameterThetaCompletion.lean}, line 218.\par}\smallskip
\Needspace{7\baselineskip}
\subsection{Row 41: 9.4 Residue parameter law, polynomiality, and leading term}
\label{row:41}
\textbf{Exact row scope.} Exact parameter addition, scalar evaluation, leading coefficient, and exact degree of the honest polynomial lift. Analytic complex powers remain outside this row.\par\smallskip
\textbf{Manuscript labels.} \code{thm:residue_section_parameter_addition}\par\smallskip
{\small\code{Komorebi.PaperIV.residuePowerSeries_add}\par}
{\footnotesize\color{muted}theorem; \code{KomorebiPaperIV/ParameterThetaCompletion.lean}, line 172.\par}\smallskip
{\small\code{Komorebi.PaperIV.residuePowerCoeff_addition}\par}
{\footnotesize\color{muted}theorem; \code{KomorebiPaperIV/ParameterThetaCompletion.lean}, line 210.\par}\smallskip
{\small\code{Komorebi.PaperIV.residuePowerCoeffPolynomial}\par}
{\footnotesize\color{muted}definition; \code{KomorebiPaperIV/ParameterThetaCompletion.lean}, line 232.\par}\smallskip
{\small\code{Komorebi.PaperIV.eval_residuePowerCoeffPolynomial}\par}
{\footnotesize\color{muted}theorem; \code{KomorebiPaperIV/ParameterThetaCompletion.lean}, line 237.\par}\smallskip
{\small\code{Komorebi.PaperIV.coeff_residuePowerCoeffPolynomial_parameterDegree}\par}
{\footnotesize\color{muted}theorem; \code{KomorebiPaperIV/ParameterThetaCompletion.lean}, line 463.\par}\smallskip
{\small\code{Komorebi.PaperIV.natDegree_residuePowerCoeffPolynomial}\par}
{\footnotesize\color{muted}theorem; \code{KomorebiPaperIV/ParameterThetaCompletion.lean}, line 503.\par}\smallskip
\Needspace{7\baselineskip}
\subsection{Row 42: 9.6, 9.7, 9.8, 9.9, and 9.10 Arbitrary-order residue Appell system}
\label{row:42}
\textbf{Exact row scope.} Exact arbitrary-order formal Appell specialization: sparse formula, lowering, translation, joint convolution, monicity, covariance, monomial symbol action, and opposite-order finite operator inversion. The two operator encodings are not yet bridged into the manuscript's single displayed named inverse relation.\par\smallskip
\textbf{Manuscript labels.} \code{cor:appell_symmetry}, \code{prop:appell_operator_form}, \code{prop:explicit_appell}, \code{thm:appell_laws}, \code{thm:appell_two_parameter_convolution}\par\smallskip
{\small\code{Komorebi.PaperIV.expandedResidueNegativeLog}\par}
{\footnotesize\color{muted}definition; \code{KomorebiPaperIV/ParameterThetaCompletion.lean}, line 528.\par}\smallskip
{\small\code{Komorebi.PaperIV.constantCoeff_expandedResidueNegativeLog}\par}
{\footnotesize\color{muted}theorem; \code{KomorebiPaperIV/ParameterThetaCompletion.lean}, line 533.\par}\smallskip
{\small\code{Komorebi.PaperIV.expand_residuePowerSeries}\par}
{\footnotesize\color{muted}theorem; \code{KomorebiPaperIV/ParameterThetaCompletion.lean}, line 539.\par}\smallskip
{\small\code{Komorebi.PaperIV.egfCoeff_expand_residuePowerSeries}\par}
{\footnotesize\color{muted}theorem; \code{KomorebiPaperIV/ParameterThetaCompletion.lean}, line 552.\par}\smallskip
{\small\code{Komorebi.PaperIV.residueAppellPolynomialParam}\par}
{\footnotesize\color{muted}definition; \code{KomorebiPaperIV/ParameterThetaCompletion.lean}, line 564.\par}\smallskip
{\small\code{Komorebi.PaperIV.residueAppellPolynomialParam_eq_block}\par}
{\footnotesize\color{muted}theorem; \code{KomorebiPaperIV/ParameterThetaCompletion.lean}, line 568.\par}\smallskip
{\small\code{Komorebi.PaperIV.residueAppellPolynomialParam_explicit}\par}
{\footnotesize\color{muted}theorem; \code{KomorebiPaperIV/ParameterThetaCompletion.lean}, line 581.\par}\smallskip
{\small\code{Komorebi.PaperIV.derivative_residueAppellPolynomialParam}\par}
{\footnotesize\color{muted}theorem; \code{KomorebiPaperIV/ParameterThetaCompletion.lean}, line 592.\par}\smallskip
{\small\code{Komorebi.PaperIV.eval_residueAppellPolynomialParam_add}\par}
{\footnotesize\color{muted}theorem; \code{KomorebiPaperIV/ParameterThetaCompletion.lean}, line 601.\par}\smallskip
{\small\code{Komorebi.PaperIV.eval_residueAppellPolynomialParam_two_parameter}\par}
{\footnotesize\color{muted}theorem; \code{KomorebiPaperIV/ParameterThetaCompletion.lean}, line 613.\par}\smallskip
{\small\code{Komorebi.PaperIV.residueAppellPolynomialParam_isMonicOfDegree}\par}
{\footnotesize\color{muted}theorem; \code{KomorebiPaperIV/ParameterThetaCompletion.lean}, line 627.\par}\smallskip
{\small\code{Komorebi.PaperIV.eval_residueAppellPolynomialParam_covariance}\par}
{\footnotesize\color{muted}theorem; \code{KomorebiPaperIV/ParameterThetaCompletion.lean}, line 636.\par}\smallskip
{\small\code{Komorebi.PaperIV.egfConstantCoeffAction_residuePowerSeries_X_pow}\par}
{\footnotesize\color{muted}theorem; \code{KomorebiPaperIV/ParameterThetaCompletion.lean}, line 645.\par}\smallskip
{\small\code{Komorebi.PaperIV.truncatedDifferentialOperator_residuePowerSeries_inverse}\par}
{\footnotesize\color{muted}theorem; \code{KomorebiPaperIV/ParameterThetaCompletion.lean}, line 660.\par}\smallskip
\Needspace{7\baselineskip}
\subsection{Row 43: 10.2 Theta polynomial and degree}
\label{row:43}
\textbf{Exact row scope.} Exact polynomial lift and scalar evaluation with degree at most \code{N*n}, completing the formal polynomial clause of theta reconstruction.\par\smallskip
\textbf{Manuscript labels.} \code{thm:theta_reconstruction}\par\smallskip
{\small\code{Komorebi.PaperIV.thetaPolynomial}\par}
{\footnotesize\color{muted}definition; \code{KomorebiPaperIV/ParameterThetaCompletion.lean}, line 714.\par}\smallskip
{\small\code{Komorebi.PaperIV.eval_thetaPolynomial}\par}
{\footnotesize\color{muted}theorem; \code{KomorebiPaperIV/ParameterThetaCompletion.lean}, line 763.\par}\smallskip
{\small\code{Komorebi.PaperIV.natDegree_thetaPolynomial_le}\par}
{\footnotesize\color{muted}theorem; \code{KomorebiPaperIV/ParameterThetaCompletion.lean}, line 779.\par}\smallskip
\Needspace{7\baselineskip}
\subsection{Row 44: 10.3 Theta parameter factors}
\label{row:44}
\textbf{Exact row scope.} Exact literal polynomial divisibility by \code{X}, and by \code{X+1} when \code{N} is even, for positive coefficient index.\par\smallskip
\textbf{Manuscript labels.} \code{prop:theta_factors}\par\smallskip
{\small\code{Komorebi.PaperIV.X_dvd_thetaPolynomial}\par}
{\footnotesize\color{muted}theorem; \code{KomorebiPaperIV/ParameterThetaCompletion.lean}, line 831.\par}\smallskip
{\small\code{Komorebi.PaperIV.X_add_one_dvd_thetaPolynomial_of_even}\par}
{\footnotesize\color{muted}theorem; \code{KomorebiPaperIV/ParameterThetaCompletion.lean}, line 839.\par}\smallskip
\Needspace{7\baselineskip}
\subsection{Row 45: 10.4 Theta Newton recurrence}
\label{row:45}
\textbf{Exact row scope.} Exact formal logarithmic derivative, Newton recurrence in equivalent antidiagonal indexing, initial value, and \code{Theta_(N,1)=-(u^N+u)}.\par\smallskip
\textbf{Manuscript labels.} \code{prop:theta_newton_recurrence}\par\smallskip
{\small\code{Komorebi.PaperIV.derivative_rescale_series}\par}
{\footnotesize\color{muted}theorem; \code{KomorebiPaperIV/ParameterThetaCompletion.lean}, line 859.\par}\smallskip
{\small\code{Komorebi.PaperIV.rescale_C_series}\par}
{\footnotesize\color{muted}theorem; \code{KomorebiPaperIV/ParameterThetaCompletion.lean}, line 870.\par}\smallskip
{\small\code{Komorebi.PaperIV.scaledReciprocalLogDerivative}\par}
{\footnotesize\color{muted}definition; \code{KomorebiPaperIV/ParameterThetaCompletion.lean}, line 877.\par}\smallskip
{\small\code{Komorebi.PaperIV.derivative_rescale_reciprocalSeries}\par}
{\footnotesize\color{muted}theorem; \code{KomorebiPaperIV/ParameterThetaCompletion.lean}, line 882.\par}\smallskip
{\small\code{Komorebi.PaperIV.thetaLogDerivativeSeries}\par}
{\footnotesize\color{muted}definition; \code{KomorebiPaperIV/ParameterThetaCompletion.lean}, line 895.\par}\smallskip
{\small\code{Komorebi.PaperIV.derivative_thetaSeries}\par}
{\footnotesize\color{muted}theorem; \code{KomorebiPaperIV/ParameterThetaCompletion.lean}, line 899.\par}\smallskip
{\small\code{Komorebi.PaperIV.coeff_thetaLogDerivativeSeries}\par}
{\footnotesize\color{muted}theorem; \code{KomorebiPaperIV/ParameterThetaCompletion.lean}, line 911.\par}\smallskip
{\small\code{Komorebi.PaperIV.thetaOrdCoeff}\par}
{\footnotesize\color{muted}definition; \code{KomorebiPaperIV/ParameterThetaCompletion.lean}, line 925.\par}\smallskip
{\small\code{Komorebi.PaperIV.thetaOrdCoeff_zero}\par}
{\footnotesize\color{muted}theorem; \code{KomorebiPaperIV/ParameterThetaCompletion.lean}, line 929.\par}\smallskip
{\small\code{Komorebi.PaperIV.thetaOrdCoeff_newton}\par}
{\footnotesize\color{muted}theorem; \code{KomorebiPaperIV/ParameterThetaCompletion.lean}, line 938.\par}\smallskip
{\small\code{Komorebi.PaperIV.factorial_mul_kappaCoeff_one}\par}
{\footnotesize\color{muted}theorem; \code{KomorebiPaperIV/ParameterThetaCompletion.lean}, line 952.\par}\smallskip
{\small\code{Komorebi.PaperIV.thetaCoeff_one}\par}
{\footnotesize\color{muted}theorem; \code{KomorebiPaperIV/ParameterThetaCompletion.lean}, line 979.\par}\smallskip
\Needspace{7\baselineskip}
\subsection{Row 46: 7.2 Bell formula and 7.8 parameter divisibility}
\label{row:46}
\textbf{Exact row scope.} Exact standard Bell generating-function specialization, every reciprocal coefficient, exact polynomial degree and leading coefficient over any characteristic-zero field, and generic-field parameter divisibility. This closes the full formal statement of 7.2; it asserts no global logarithm branch.\par\smallskip
\textbf{Manuscript labels.} \code{cor:factor_r}, \code{thm:reciprocal_bell}\par\smallskip
{\small\code{Komorebi.PaperIV.bellInput}\par}
{\footnotesize\color{muted}definition; \code{KomorebiPaperIV/BellCompletion.lean}, line 21.\par}\smallskip
{\small\code{Komorebi.PaperIV.completeExponentialBell}\par}
{\footnotesize\color{muted}definition; \code{KomorebiPaperIV/BellCompletion.lean}, line 25.\par}\smallskip
{\small\code{Komorebi.PaperIV.constantCoeff_bellInput}\par}
{\footnotesize\color{muted}theorem; \code{KomorebiPaperIV/BellCompletion.lean}, line 30.\par}\smallskip
{\small\code{Komorebi.PaperIV.completeExponentialBell_generating_coefficient}\par}
{\footnotesize\color{muted}theorem; \code{KomorebiPaperIV/BellCompletion.lean}, line 35.\par}\smallskip
{\small\code{Komorebi.PaperIV.completeExponentialBell_zero}\par}
{\footnotesize\color{muted}theorem; \code{KomorebiPaperIV/BellCompletion.lean}, line 41.\par}\smallskip
{\small\code{Komorebi.PaperIV.bellInput_reciprocal}\par}
{\footnotesize\color{muted}theorem; \code{KomorebiPaperIV/BellCompletion.lean}, line 50.\par}\smallskip
{\small\code{Komorebi.PaperIV.reciprocalCoeff_bell_formula}\par}
{\footnotesize\color{muted}theorem; \code{KomorebiPaperIV/BellCompletion.lean}, line 66.\par}\smallskip
{\small\code{Komorebi.PaperIV.coeff_one_negativeLogDenominator}\par}
{\footnotesize\color{muted}theorem; \code{KomorebiPaperIV/BellCompletion.lean}, line 76.\par}\smallskip
{\small\code{Komorebi.PaperIV.coeff_reciprocalCoeffPolynomialOver_parameterDegree}\par}
{\footnotesize\color{muted}theorem; \code{KomorebiPaperIV/BellCompletion.lean}, line 86.\par}\smallskip
{\small\code{Komorebi.PaperIV.natDegree_reciprocalCoeffPolynomialOver}\par}
{\footnotesize\color{muted}theorem; \code{KomorebiPaperIV/BellCompletion.lean}, line 111.\par}\smallskip
{\small\code{Komorebi.PaperIV.X_dvd_reciprocalCoeffPolynomialOver}\par}
{\footnotesize\color{muted}theorem; \code{KomorebiPaperIV/BellCompletion.lean}, line 124.\par}\smallskip
\Needspace{7\baselineskip}
\subsection{Row 47: 4.1, 4.2, 5.1, 5.2, 5.3, and 5.6 Entire realization}
\label{row:47}
\textbf{Exact row scope.} Actual entire functions, convergence of the existing formal transfer/signed series at every complex argument, independence of root choices, pointwise finite exponential formulas, iterated derivative expressions, full-period ODEs, and exact initial data. Existence of a basis of the full ODE solution space and the one-step section-to-section derivative identities are not asserted by this row.\par\smallskip
\textbf{Manuscript labels.} \code{def:signed_sections}, \code{def:transfer_exponential}, \code{prop:signed_exponential_relation}, \code{thm:exponential_cycle}, \code{thm:finite_exponential_representation}, \code{thm:signed_cycle}\par\smallskip
{\small\code{Komorebi.PaperIV.exponentialSumJet}\par}
{\footnotesize\color{muted}definition; \code{KomorebiPaperIV/EntireBridge.lean}, line 21.\par}\smallskip
{\small\code{Komorebi.PaperIV.analyticAt_exponentialSumJet}\par}
{\footnotesize\color{muted}theorem; \code{KomorebiPaperIV/EntireBridge.lean}, line 24.\par}\smallskip
{\small\code{Komorebi.PaperIV.hasDerivAt_exponentialSumJet}\par}
{\footnotesize\color{muted}theorem; \code{KomorebiPaperIV/EntireBridge.lean}, line 29.\par}\smallskip
{\small\code{Komorebi.PaperIV.iteratedDeriv_exponentialSum}\par}
{\footnotesize\color{muted}theorem; \code{KomorebiPaperIV/EntireBridge.lean}, line 43.\par}\smallskip
{\small\code{Komorebi.PaperIV.exponentialSum_ode}\par}
{\footnotesize\color{muted}theorem; \code{KomorebiPaperIV/EntireBridge.lean}, line 53.\par}\smallskip
{\small\code{Komorebi.PaperIV.entireTransfer}\par}
{\footnotesize\color{muted}definition; \code{KomorebiPaperIV/EntireBridge.lean}, line 66.\par}\smallskip
{\small\code{Komorebi.PaperIV.analyticAt_entireTransfer}\par}
{\footnotesize\color{muted}theorem; \code{KomorebiPaperIV/EntireBridge.lean}, line 71.\par}\smallskip
{\small\code{Komorebi.PaperIV.entireTransfer_eq_exponentialSum}\par}
{\footnotesize\color{muted}theorem; \code{KomorebiPaperIV/EntireBridge.lean}, line 76.\par}\smallskip
{\small\code{Komorebi.PaperIV.hasSum_finiteRootAverage}\par}
{\footnotesize\color{muted}theorem; \code{KomorebiPaperIV/EntireBridge.lean}, line 85.\par}\smallskip
{\small\code{Komorebi.PaperIV.hasSum_entireTransfer}\par}
{\footnotesize\color{muted}theorem; \code{KomorebiPaperIV/EntireBridge.lean}, line 106.\par}\smallskip
{\small\code{Komorebi.PaperIV.entireTransfer_eq_tsum}\par}
{\footnotesize\color{muted}theorem; \code{KomorebiPaperIV/EntireBridge.lean}, line 113.\par}\smallskip
{\small\code{Komorebi.PaperIV.entireTransfer_independent_root}\par}
{\footnotesize\color{muted}theorem; \code{KomorebiPaperIV/EntireBridge.lean}, line 120.\par}\smallskip
{\small\code{Komorebi.PaperIV.entireTransfer_ode}\par}
{\footnotesize\color{muted}theorem; \code{KomorebiPaperIV/EntireBridge.lean}, line 127.\par}\smallskip
{\small\code{Komorebi.PaperIV.entireSigned}\par}
{\footnotesize\color{muted}definition; \code{KomorebiPaperIV/EntireBridge.lean}, line 138.\par}\smallskip
{\small\code{Komorebi.PaperIV.analyticAt_entireSigned}\par}
{\footnotesize\color{muted}theorem; \code{KomorebiPaperIV/EntireBridge.lean}, line 141.\par}\smallskip
{\small\code{Komorebi.PaperIV.hasSum_entireSigned}\par}
{\footnotesize\color{muted}theorem; \code{KomorebiPaperIV/EntireBridge.lean}, line 146.\par}\smallskip
{\small\code{Komorebi.PaperIV.entireSigned_eq_tsum}\par}
{\footnotesize\color{muted}theorem; \code{KomorebiPaperIV/EntireBridge.lean}, line 155.\par}\smallskip
{\small\code{Komorebi.PaperIV.entireSigned_independent_roots}\par}
{\footnotesize\color{muted}theorem; \code{KomorebiPaperIV/EntireBridge.lean}, line 162.\par}\smallskip
{\small\code{Komorebi.PaperIV.entireSigned_zero_eq_average}\par}
{\footnotesize\color{muted}theorem; \code{KomorebiPaperIV/EntireBridge.lean}, line 171.\par}\smallskip
{\small\code{Komorebi.PaperIV.entireTransfer_derivative_at_zero}\par}
{\footnotesize\color{muted}theorem; \code{KomorebiPaperIV/EntireBridge.lean}, line 177.\par}\smallskip
{\small\code{Komorebi.PaperIV.entireTransfer_initial_data}\par}
{\footnotesize\color{muted}theorem; \code{KomorebiPaperIV/EntireBridge.lean}, line 194.\par}\smallskip
{\small\code{Komorebi.PaperIV.entireSigned_eq_exponentialSum}\par}
{\footnotesize\color{muted}theorem; \code{KomorebiPaperIV/EntireBridge.lean}, line 207.\par}\smallskip
{\small\code{Komorebi.PaperIV.entireSigned_ode}\par}
{\footnotesize\color{muted}theorem; \code{KomorebiPaperIV/EntireBridge.lean}, line 215.\par}\smallskip
{\small\code{Komorebi.PaperIV.entireSigned_derivative_eq_scaled_transfer}\par}
{\footnotesize\color{muted}theorem; \code{KomorebiPaperIV/EntireBridge.lean}, line 227.\par}\smallskip
{\small\code{Komorebi.PaperIV.entireSigned_initial_data}\par}
{\footnotesize\color{muted}theorem; \code{KomorebiPaperIV/EntireBridge.lean}, line 234.\par}\smallskip
\Needspace{7\baselineskip}
\subsection{Row 48: 4.3 and 5.4 Analytic Wronskians}
\label{row:48}
\textbf{Exact row scope.} Exact analytic Wronskians using actual iterated derivatives: lambda\textasciicircum{}(N*(N-1)/2) for the ordered transfer family and 1 for the ordered signed family, for N>1. The transfer equality also holds at lambda=0. The proofs derive determinant factorization and constancy; no Wronskian or determinant value is assumed.\par\smallskip
\textbf{Manuscript labels.} \code{cor:signed_wronskian}, \code{cor:wronskians}\par\smallskip
{\small\code{Komorebi.PaperIV.exponentialJetMatrix}\par}
{\footnotesize\color{muted}definition; \code{KomorebiPaperIV/Wronskians.lean}, line 21.\par}\smallskip
{\small\code{Komorebi.PaperIV.exponentialJetMatrix_factor}\par}
{\footnotesize\color{muted}theorem; \code{KomorebiPaperIV/Wronskians.lean}, line 25.\par}\smallskip
{\small\code{Komorebi.PaperIV.det_exponentialJetMatrix}\par}
{\footnotesize\color{muted}theorem; \code{KomorebiPaperIV/Wronskians.lean}, line 40.\par}\smallskip
{\small\code{Komorebi.PaperIV.det_exponentialJetMatrix_constant}\par}
{\footnotesize\color{muted}theorem; \code{KomorebiPaperIV/Wronskians.lean}, line 51.\par}\smallskip
{\small\code{Komorebi.PaperIV.analyticWronskian}\par}
{\footnotesize\color{muted}definition; \code{KomorebiPaperIV/Wronskians.lean}, line 58.\par}\smallskip
{\small\code{Komorebi.PaperIV.transferWronskian_constant}\par}
{\footnotesize\color{muted}theorem; \code{KomorebiPaperIV/Wronskians.lean}, line 61.\par}\smallskip
{\small\code{Komorebi.PaperIV.transferWronskian_eq}\par}
{\footnotesize\color{muted}theorem; \code{KomorebiPaperIV/Wronskians.lean}, line 73.\par}\smallskip
{\small\code{Komorebi.PaperIV.signedWronskian_constant}\par}
{\footnotesize\color{muted}theorem; \code{KomorebiPaperIV/Wronskians.lean}, line 90.\par}\smallskip
{\small\code{Komorebi.PaperIV.signedWronskian_eq_one}\par}
{\footnotesize\color{muted}theorem; \code{KomorebiPaperIV/Wronskians.lean}, line 102.\par}\smallskip
\Needspace{7\baselineskip}
\subsection{Row 49: 6.4 Companion first column and trace}
\label{row:49}
\textbf{Exact row scope.} Convergent scalar series for each entry of the actual matrix exponential, exact first-column identification with the signed entire sections, every diagonal entry equal to the denominator, and trace=N*denominator. The full off-diagonal entry formula, characteristic polynomial, spectrum, and determinant-one claim remain outside this row.\par\smallskip
\textbf{Manuscript labels.} \code{thm:companion_matrix_realization}\par\smallskip
{\small\code{Komorebi.PaperIV.hasSum_signedFundamentalMatrix_entry}\par}
{\footnotesize\color{muted}theorem; \code{KomorebiPaperIV/CompanionBridge.lean}, line 21.\par}\smallskip
{\small\code{Komorebi.PaperIV.signedFundamentalMatrix_first_column}\par}
{\footnotesize\color{muted}theorem; \code{KomorebiPaperIV/CompanionBridge.lean}, line 35.\par}\smallskip
{\small\code{Komorebi.PaperIV.signedFundamentalMatrix_diagonal}\par}
{\footnotesize\color{muted}theorem; \code{KomorebiPaperIV/CompanionBridge.lean}, line 72.\par}\smallskip
{\small\code{Komorebi.PaperIV.signedFundamentalMatrix_trace}\par}
{\footnotesize\color{muted}theorem; \code{KomorebiPaperIV/CompanionBridge.lean}, line 85.\par}\smallskip
\Needspace{7\baselineskip}
\subsection{Row 50: 6.5 and 6.6 CompanionCompletion}
\label{row:50}
\textbf{Exact row scope.} Exact convergent derivative-series bridge, every signed off-diagonal entry, and the determinant-one identity derived from the actual Wronskian.\par\smallskip
\textbf{Manuscript labels.} \code{cor:companion_addition_inversion}, \code{cor:fundamental_matrix_entries}\par\smallskip
{\small\code{Komorebi.PaperIV.hasSum_exponentialSumJet}\par}
{\footnotesize\color{muted}theorem; \code{KomorebiPaperIV/CompanionCompletion.lean}, line 20.\par}\smallskip
{\small\code{Komorebi.PaperIV.entireSigned_derivative_at_zero}\par}
{\footnotesize\color{muted}theorem; \code{KomorebiPaperIV/CompanionCompletion.lean}, line 36.\par}\smallskip
{\small\code{Komorebi.PaperIV.hasSum_entireSigned_derivative}\par}
{\footnotesize\color{muted}theorem; \code{KomorebiPaperIV/CompanionCompletion.lean}, line 48.\par}\smallskip
{\small\code{Komorebi.PaperIV.signedFundamentalMatrix_entry_eq_derivative}\par}
{\footnotesize\color{muted}theorem; \code{KomorebiPaperIV/CompanionCompletion.lean}, line 66.\par}\smallskip
{\small\code{Komorebi.PaperIV.signedFundamentalMatrix_det}\par}
{\footnotesize\color{muted}theorem; \code{KomorebiPaperIV/CompanionCompletion.lean}, line 84.\par}\smallskip
{\small\code{Komorebi.PaperIV.companionDisplacement}\par}
{\footnotesize\color{muted}definition; \code{KomorebiPaperIV/CompanionCompletion.lean}, line 97.\par}\smallskip
{\small\code{Komorebi.PaperIV.companionDisplacement_lt}\par}
{\footnotesize\color{muted}theorem; \code{KomorebiPaperIV/CompanionCompletion.lean}, line 100.\par}\smallskip
{\small\code{Komorebi.PaperIV.companionEntrySign}\par}
{\footnotesize\color{muted}definition; \code{KomorebiPaperIV/CompanionCompletion.lean}, line 108.\par}\smallskip
{\small\code{Komorebi.PaperIV.signedShiftMatrix_pow_entry_packet}\par}
{\footnotesize\color{muted}theorem; \code{KomorebiPaperIV/CompanionCompletion.lean}, line 112.\par}\smallskip
{\small\code{Komorebi.PaperIV.signedFundamentalMatrix_entry}\par}
{\footnotesize\color{muted}theorem; \code{KomorebiPaperIV/CompanionCompletion.lean}, line 159.\par}\smallskip
\Needspace{7\baselineskip}
\subsection{Row 51: 6.4 CompanionSpectrum}
\label{row:51}
\textbf{Exact row scope.} Exact characteristic polynomial X\textasciicircum{}N+1, spectrum equal to the Nth roots of -1, N roots counted with multiplicity, and no repeated roots. No spectral certificate is assumed.\par\smallskip
\textbf{Manuscript labels.} \code{thm:companion_matrix_realization}\par\smallskip
{\small\code{Komorebi.PaperIV.aeval_signedShiftMatrix_first_column}\par}
{\footnotesize\color{muted}theorem; \code{KomorebiPaperIV/CompanionSpectrum.lean}, line 23.\par}\smallskip
{\small\code{Komorebi.PaperIV.signedShiftMatrix_no_small_annihilator}\par}
{\footnotesize\color{muted}theorem; \code{KomorebiPaperIV/CompanionSpectrum.lean}, line 38.\par}\smallskip
{\small\code{Komorebi.PaperIV.aeval_signedShiftMatrix_X_pow_add_one}\par}
{\footnotesize\color{muted}theorem; \code{KomorebiPaperIV/CompanionSpectrum.lean}, line 49.\par}\smallskip
{\small\code{Komorebi.PaperIV.signedShiftMatrix_charpoly}\par}
{\footnotesize\color{muted}theorem; \code{KomorebiPaperIV/CompanionSpectrum.lean}, line 55.\par}\smallskip
{\small\code{Komorebi.PaperIV.signedShiftMatrix_mem_spectrum_iff}\par}
{\footnotesize\color{muted}theorem; \code{KomorebiPaperIV/CompanionSpectrum.lean}, line 69.\par}\smallskip
{\small\code{Komorebi.PaperIV.signedShiftMatrix_charpoly_separable}\par}
{\footnotesize\color{muted}theorem; \code{KomorebiPaperIV/CompanionSpectrum.lean}, line 75.\par}\smallskip
{\small\code{Komorebi.PaperIV.signedShiftMatrix_charpoly_roots_card}\par}
{\footnotesize\color{muted}theorem; \code{KomorebiPaperIV/CompanionSpectrum.lean}, line 82.\par}\smallskip
{\small\code{Komorebi.PaperIV.signedShiftMatrix_charpoly_roots_nodup}\par}
{\footnotesize\color{muted}theorem; \code{KomorebiPaperIV/CompanionSpectrum.lean}, line 89.\par}\smallskip
\Needspace{7\baselineskip}
\subsection{Row 52: 4.2 ODESolutions}
\label{row:52}
\textbf{Exact row scope.} Exact uniqueness and explicit expansion for every entire analytic solution of y\textasciicircum{}(N)=lambda\textasciicircum{}N y, with unique coefficients when lambda is nonzero. This is the finite-basis assertion for the entire solution space.\par\smallskip
\textbf{Manuscript labels.} \code{thm:exponential_cycle}\par\smallskip
{\small\code{Komorebi.PaperIV.ode_all_derivatives_zero}\par}
{\footnotesize\color{muted}theorem; \code{KomorebiPaperIV/ODESolutions.lean}, line 14.\par}\smallskip
{\small\code{Komorebi.PaperIV.entire_ode_zero_of_initial_zero}\par}
{\footnotesize\color{muted}theorem; \code{KomorebiPaperIV/ODESolutions.lean}, line 30.\par}\smallskip
{\small\code{Komorebi.PaperIV.entire_ode_unique}\par}
{\footnotesize\color{muted}theorem; \code{KomorebiPaperIV/ODESolutions.lean}, line 43.\par}\smallskip
{\small\code{Komorebi.PaperIV.transferCombination}\par}
{\footnotesize\color{muted}definition; \code{KomorebiPaperIV/ODESolutions.lean}, line 60.\par}\smallskip
{\small\code{Komorebi.PaperIV.analyticAt_transferCombination}\par}
{\footnotesize\color{muted}theorem; \code{KomorebiPaperIV/ODESolutions.lean}, line 63.\par}\smallskip
{\small\code{Komorebi.PaperIV.transferCombination_derivative}\par}
{\footnotesize\color{muted}theorem; \code{KomorebiPaperIV/ODESolutions.lean}, line 71.\par}\smallskip
{\small\code{Komorebi.PaperIV.transferCombination_ode}\par}
{\footnotesize\color{muted}theorem; \code{KomorebiPaperIV/ODESolutions.lean}, line 82.\par}\smallskip
{\small\code{Komorebi.PaperIV.transferCombination_initial}\par}
{\footnotesize\color{muted}theorem; \code{KomorebiPaperIV/ODESolutions.lean}, line 93.\par}\smallskip
{\small\code{Komorebi.PaperIV.entireTransfer_solution_expansion}\par}
{\footnotesize\color{muted}theorem; \code{KomorebiPaperIV/ODESolutions.lean}, line 108.\par}\smallskip
{\small\code{Komorebi.PaperIV.transferCombination_injective}\par}
{\footnotesize\color{muted}theorem; \code{KomorebiPaperIV/ODESolutions.lean}, line 124.\par}\smallskip
\Needspace{7\baselineskip}
\subsection{Row 53: 4.2 and 5.3 AnalyticCycles}
\label{row:53}
\textbf{Exact row scope.} Exact pointwise analytic derivative cycles, including both wraparound cases. The generic bridge derives analytic differential identities from the previously defined formal Taylor coefficients.\par\smallskip
\textbf{Manuscript labels.} \code{thm:exponential_cycle}, \code{thm:signed_cycle}\par\smallskip
{\small\code{Komorebi.PaperIV.entire_eq_of_all_derivatives}\par}
{\footnotesize\color{muted}theorem; \code{KomorebiPaperIV/AnalyticCycles.lean}, line 15.\par}\smallskip
{\small\code{Komorebi.PaperIV.entire_derivative_of_formal}\par}
{\footnotesize\color{muted}theorem; \code{KomorebiPaperIV/AnalyticCycles.lean}, line 31.\par}\smallskip
{\small\code{Komorebi.PaperIV.entireTransfer_derivative_of_pos}\par}
{\footnotesize\color{muted}theorem; \code{KomorebiPaperIV/AnalyticCycles.lean}, line 49.\par}\smallskip
{\small\code{Komorebi.PaperIV.entireTransfer_derivative_zero}\par}
{\footnotesize\color{muted}theorem; \code{KomorebiPaperIV/AnalyticCycles.lean}, line 61.\par}\smallskip
{\small\code{Komorebi.PaperIV.entireSigned_derivative_of_pos}\par}
{\footnotesize\color{muted}theorem; \code{KomorebiPaperIV/AnalyticCycles.lean}, line 72.\par}\smallskip
{\small\code{Komorebi.PaperIV.entireSigned_derivative_zero}\par}
{\footnotesize\color{muted}theorem; \code{KomorebiPaperIV/AnalyticCycles.lean}, line 84.\par}\smallskip
\Needspace{7\baselineskip}
\subsection{Row 54: 6.6 and 11.2 ExponentialBounds}
\label{row:54}
\textbf{Exact row scope.} Explicit root witnesses and unconditional companion determinant for N>1; exact global bound norm(D\_N(z)) <= exp(norm(z)). The order and matching type lower bound in Proposition 11.2 remain outside this row.\par\smallskip
\textbf{Manuscript labels.} \code{cor:companion_addition_inversion}, \code{prop:orders}\par\smallskip
{\small\code{Komorebi.PaperIV.canonicalPrimitiveRoot}\par}
{\footnotesize\color{muted}definition; \code{KomorebiPaperIV/ExponentialBounds.lean}, line 12.\par}\smallskip
{\small\code{Komorebi.PaperIV.canonicalSignedRoot}\par}
{\footnotesize\color{muted}definition; \code{KomorebiPaperIV/ExponentialBounds.lean}, line 15.\par}\smallskip
{\small\code{Komorebi.PaperIV.canonicalPrimitiveRoot_isPrimitive}\par}
{\footnotesize\color{muted}theorem; \code{KomorebiPaperIV/ExponentialBounds.lean}, line 18.\par}\smallskip
{\small\code{Komorebi.PaperIV.canonicalSignedRoot_pow}\par}
{\footnotesize\color{muted}theorem; \code{KomorebiPaperIV/ExponentialBounds.lean}, line 22.\par}\smallskip
{\small\code{Komorebi.PaperIV.signedFundamentalMatrix_det_one}\par}
{\footnotesize\color{muted}theorem; \code{KomorebiPaperIV/ExponentialBounds.lean}, line 32.\par}\smallskip
{\small\code{Komorebi.PaperIV.norm_signedRoot}\par}
{\footnotesize\color{muted}theorem; \code{KomorebiPaperIV/ExponentialBounds.lean}, line 38.\par}\smallskip
{\small\code{Komorebi.PaperIV.norm_entireSigned_zero_le}\par}
{\footnotesize\color{muted}theorem; \code{KomorebiPaperIV/ExponentialBounds.lean}, line 46.\par}\smallskip
\Needspace{7\baselineskip}
\subsection{Row 55: 11.6 and 11.10 ZeroGeometry}
\label{row:55}
\textbf{Exact row scope.} Full near-cancellation lemma with universal epsilon=1/2 and C=3/2, including uniqueness, with no strip restriction. Full order-two cosine zero set in integer and positive/negative natural forms, analytic multiplicity one, and exact lattice spacing.\par\smallskip
\textbf{Manuscript labels.} \code{cor:case_two_zeros}, \code{lem:near_cancellation}\par\smallskip
{\small\code{Komorebi.PaperIV.entireSigned_two_eq_cos}\par}
{\footnotesize\color{muted}theorem; \code{KomorebiPaperIV/ZeroGeometry.lean}, line 20.\par}\smallskip
{\small\code{Komorebi.PaperIV.entireSigned_two_zero_iff}\par}
{\footnotesize\color{muted}theorem; \code{KomorebiPaperIV/ZeroGeometry.lean}, line 30.\par}\smallskip
{\small\code{Komorebi.PaperIV.entireSigned_two_zero_derivative_ne_zero}\par}
{\footnotesize\color{muted}theorem; \code{KomorebiPaperIV/ZeroGeometry.lean}, line 37.\par}\smallskip
{\small\code{Komorebi.PaperIV.entireSigned_two_zero_iff_nat}\par}
{\footnotesize\color{muted}theorem; \code{KomorebiPaperIV/ZeroGeometry.lean}, line 52.\par}\smallskip
{\small\code{Komorebi.PaperIV.entireSigned_two_zero_order_one}\par}
{\footnotesize\color{muted}theorem; \code{KomorebiPaperIV/ZeroGeometry.lean}, line 78.\par}\smallskip
{\small\code{Komorebi.PaperIV.orderTwo_zero_spacing}\par}
{\footnotesize\color{muted}theorem; \code{KomorebiPaperIV/ZeroGeometry.lean}, line 86.\par}\smallskip
{\small\code{Komorebi.PaperIV.oddExponentialCenter}\par}
{\footnotesize\color{muted}definition; \code{KomorebiPaperIV/ZeroGeometry.lean}, line 90.\par}\smallskip
{\small\code{Komorebi.PaperIV.near_cancellation_exists}\par}
{\footnotesize\color{muted}theorem; \code{KomorebiPaperIV/ZeroGeometry.lean}, line 93.\par}\smallskip
{\small\code{Komorebi.PaperIV.oddExponentialCenter_separation}\par}
{\footnotesize\color{muted}theorem; \code{KomorebiPaperIV/ZeroGeometry.lean}, line 114.\par}\smallskip
{\small\code{Komorebi.PaperIV.near_cancellation_unique}\par}
{\footnotesize\color{muted}theorem; \code{KomorebiPaperIV/ZeroGeometry.lean}, line 129.\par}\smallskip
\Needspace{7\baselineskip}
\subsection{Row 56: 5.3 Signed ODE basis}
\label{row:56}
\textbf{Exact row scope.} Exact spanning and unique coefficients for every entire analytic solution of y\textasciicircum{}(N)=-y, in the signed-section basis with its stated initial normalization.\par\smallskip
\textbf{Manuscript labels.} \code{thm:signed_cycle}\par\smallskip
{\small\code{Komorebi.PaperIV.signedCombination}\par}
{\footnotesize\color{muted}definition; \code{KomorebiPaperIV/ODESolutions.lean}, line 134.\par}\smallskip
{\small\code{Komorebi.PaperIV.signedCombination_eq_transfer}\par}
{\footnotesize\color{muted}theorem; \code{KomorebiPaperIV/ODESolutions.lean}, line 137.\par}\smallskip
{\small\code{Komorebi.PaperIV.entireSigned_solution_expansion}\par}
{\footnotesize\color{muted}theorem; \code{KomorebiPaperIV/ODESolutions.lean}, line 148.\par}\smallskip
{\small\code{Komorebi.PaperIV.signedCombination_injective}\par}
{\footnotesize\color{muted}theorem; \code{KomorebiPaperIV/ODESolutions.lean}, line 161.\par}\smallskip
\Needspace{7\baselineskip}
\subsection{Row 57: 11.7 Local zero existence and 11.8 leading correction}
\label{row:57}
\textbf{Exact row scope.} General two-exponential perturbation theorem, proved by the actual logarithmic contraction. Explicit remainder and derivative bounds imply unique local zero, multiplicity one, displacement and quadratic correction bounds. These helper theorems do not assume or prove global exhaustion.\par\smallskip
\textbf{Manuscript labels.} \code{thm:first_zero_correction}, \code{thm:large_zero_strings}\par\smallskip
{\small\code{Komorebi.PaperIV.norm_one_add_lower_half}\par}
{\footnotesize\color{muted}theorem; \code{KomorebiPaperIV/ZeroPerturbation.lean}, line 19.\par}\smallskip
{\small\code{Komorebi.PaperIV.log_one_add_lipschitz_half}\par}
{\footnotesize\color{muted}theorem; \code{KomorebiPaperIV/ZeroPerturbation.lean}, line 25.\par}\smallskip
{\small\code{Komorebi.PaperIV.logarithmicZeroMap}\par}
{\footnotesize\color{muted}definition; \code{KomorebiPaperIV/ZeroPerturbation.lean}, line 40.\par}\smallskip
{\small\code{Komorebi.PaperIV.logarithmicZeroMap_mapsTo}\par}
{\footnotesize\color{muted}theorem; \code{KomorebiPaperIV/ZeroPerturbation.lean}, line 43.\par}\smallskip
{\small\code{Komorebi.PaperIV.logarithmicZeroMap_lipschitz}\par}
{\footnotesize\color{muted}theorem; \code{KomorebiPaperIV/ZeroPerturbation.lean}, line 53.\par}\smallskip
{\small\code{Komorebi.PaperIV.logarithmicZeroMap_fixed_iff}\par}
{\footnotesize\color{muted}theorem; \code{KomorebiPaperIV/ZeroPerturbation.lean}, line 71.\par}\smallskip
{\small\code{Komorebi.PaperIV.perturbed_exponential_unique_zero}\par}
{\footnotesize\color{muted}theorem; \code{KomorebiPaperIV/ZeroPerturbation.lean}, line 113.\par}\smallskip
{\small\code{Komorebi.PaperIV.perturbed_exponential_zero_distance}\par}
{\footnotesize\color{muted}theorem; \code{KomorebiPaperIV/ZeroPerturbation.lean}, line 139.\par}\smallskip
{\small\code{Komorebi.PaperIV.perturbed_exponential_pointwise_distance}\par}
{\footnotesize\color{muted}theorem; \code{KomorebiPaperIV/ZeroPerturbation.lean}, line 155.\par}\smallskip
{\small\code{Komorebi.PaperIV.hasDerivAt_perturbed_exponential}\par}
{\footnotesize\color{muted}theorem; \code{KomorebiPaperIV/ZeroPerturbation.lean}, line 168.\par}\smallskip
{\small\code{Komorebi.PaperIV.perturbed_exponential_zero_derivative_lower}\par}
{\footnotesize\color{muted}theorem; \code{KomorebiPaperIV/ZeroPerturbation.lean}, line 180.\par}\smallskip
{\small\code{Komorebi.PaperIV.perturbed_exponential_zero_order_one}\par}
{\footnotesize\color{muted}theorem; \code{KomorebiPaperIV/ZeroPerturbation.lean}, line 194.\par}\smallskip
{\small\code{Komorebi.PaperIV.perturbed_exponential_derivative_error}\par}
{\footnotesize\color{muted}theorem; \code{KomorebiPaperIV/ZeroPerturbation.lean}, line 207.\par}\smallskip
{\small\code{Komorebi.PaperIV.norm_log_one_add_sub_self_half}\par}
{\footnotesize\color{muted}theorem; \code{KomorebiPaperIV/ZeroPerturbation.lean}, line 215.\par}\smallskip
{\small\code{Komorebi.PaperIV.perturbed_exponential_first_correction}\par}
{\footnotesize\color{muted}theorem; \code{KomorebiPaperIV/ZeroPerturbation.lean}, line 224.\par}\smallskip
\Needspace{7\baselineskip}
\subsection{Row 58: 11.5 Two-term normalization and remainder jets}
\label{row:58}
\textbf{Exact row scope.} Exact normalized finite sum and actual signed remainder, derivative jets of every order, zero equivalence and strip estimate from stated real-part gaps. The sector-to-strip confinement clause is outside this row.\par\smallskip
\textbf{Manuscript labels.} \code{lem:two_term_dominance}\par\smallskip
{\small\code{Komorebi.PaperIV.finiteExponentialRemainder}\par}
{\footnotesize\color{muted}definition; \code{KomorebiPaperIV/ExponentialRemainder.lean}, line 17.\par}\smallskip
{\small\code{Komorebi.PaperIV.analyticAt_finiteExponentialRemainder}\par}
{\footnotesize\color{muted}theorem; \code{KomorebiPaperIV/ExponentialRemainder.lean}, line 21.\par}\smallskip
{\small\code{Komorebi.PaperIV.hasDerivAt_finiteExponentialRemainder}\par}
{\footnotesize\color{muted}theorem; \code{KomorebiPaperIV/ExponentialRemainder.lean}, line 27.\par}\smallskip
{\small\code{Komorebi.PaperIV.iteratedDeriv_finiteExponentialRemainder}\par}
{\footnotesize\color{muted}theorem; \code{KomorebiPaperIV/ExponentialRemainder.lean}, line 42.\par}\smallskip
{\small\code{Komorebi.PaperIV.finiteExponentialRemainder_strip_bound}\par}
{\footnotesize\color{muted}theorem; \code{KomorebiPaperIV/ExponentialRemainder.lean}, line 53.\par}\smallskip
{\small\code{Komorebi.PaperIV.finite_two_term_normalization}\par}
{\footnotesize\color{muted}theorem; \code{KomorebiPaperIV/ExponentialRemainder.lean}, line 82.\par}\smallskip
{\small\code{Komorebi.PaperIV.signedExponentialRemainder}\par}
{\footnotesize\color{muted}definition; \code{KomorebiPaperIV/ExponentialRemainder.lean}, line 105.\par}\smallskip
{\small\code{Komorebi.PaperIV.signed_two_term_normalization}\par}
{\footnotesize\color{muted}theorem; \code{KomorebiPaperIV/ExponentialRemainder.lean}, line 109.\par}\smallskip
{\small\code{Komorebi.PaperIV.signed_two_term_zero_iff}\par}
{\footnotesize\color{muted}theorem; \code{KomorebiPaperIV/ExponentialRemainder.lean}, line 128.\par}\smallskip
\Needspace{7\baselineskip}
\subsection{Row 59: 11.5 Canonical polygon gap and strip bounds}
\label{row:59}
\textbf{Exact row scope.} The canonical roots have the precise positive gap gamma\_N. The actual signed remainder satisfies the explicit strip estimate for every derivative order. No polygon-gap certificate is assumed. Global sector confinement is outside this row and supplied by rows 62--67.\par\smallskip
\textbf{Manuscript labels.} \code{lem:two_term_dominance}\par\smallskip
{\small\code{Komorebi.PaperIV.canonicalSignedVertex}\par}
{\footnotesize\color{muted}definition; \code{KomorebiPaperIV/RegularPolygonGap.lean}, line 12.\par}\smallskip
{\small\code{Komorebi.PaperIV.signedPolygonGap}\par}
{\footnotesize\color{muted}definition; \code{KomorebiPaperIV/RegularPolygonGap.lean}, line 15.\par}\smallskip
{\small\code{Komorebi.PaperIV.canonicalSignedVertex_exp}\par}
{\footnotesize\color{muted}theorem; \code{KomorebiPaperIV/RegularPolygonGap.lean}, line 18.\par}\smallskip
{\small\code{Komorebi.PaperIV.canonicalSignedVertex_re}\par}
{\footnotesize\color{muted}theorem; \code{KomorebiPaperIV/RegularPolygonGap.lean}, line 27.\par}\smallskip
{\small\code{Komorebi.PaperIV.canonicalSignedVertex_im}\par}
{\footnotesize\color{muted}theorem; \code{KomorebiPaperIV/RegularPolygonGap.lean}, line 31.\par}\smallskip
{\small\code{Komorebi.PaperIV.norm_canonicalSignedVertex}\par}
{\footnotesize\color{muted}theorem; \code{KomorebiPaperIV/RegularPolygonGap.lean}, line 35.\par}\smallskip
{\small\code{Komorebi.PaperIV.canonicalSignedVertex_last_angle}\par}
{\footnotesize\color{muted}theorem; \code{KomorebiPaperIV/RegularPolygonGap.lean}, line 39.\par}\smallskip
{\small\code{Komorebi.PaperIV.canonicalSignedVertex_last_re}\par}
{\footnotesize\color{muted}theorem; \code{KomorebiPaperIV/RegularPolygonGap.lean}, line 48.\par}\smallskip
{\small\code{Komorebi.PaperIV.canonicalSignedVertex_last_im}\par}
{\footnotesize\color{muted}theorem; \code{KomorebiPaperIV/RegularPolygonGap.lean}, line 52.\par}\smallskip
{\small\code{Komorebi.PaperIV.signedPolygonGap_pos}\par}
{\footnotesize\color{muted}theorem; \code{KomorebiPaperIV/RegularPolygonGap.lean}, line 56.\par}\smallskip
{\small\code{Komorebi.PaperIV.canonicalSignedVertex_re_gap}\par}
{\footnotesize\color{muted}theorem; \code{KomorebiPaperIV/RegularPolygonGap.lean}, line 64.\par}\smallskip
{\small\code{Komorebi.PaperIV.canonicalSignedRemainder_strip_bound}\par}
{\footnotesize\color{muted}theorem; \code{KomorebiPaperIV/RegularPolygonGap.lean}, line 95.\par}\smallskip
{\small\code{Komorebi.PaperIV.canonicalSignedVertex_adjacent_difference}\par}
{\footnotesize\color{muted}theorem; \code{KomorebiPaperIV/RegularPolygonGap.lean}, line 117.\par}\smallskip
\Needspace{7\baselineskip}
\subsection{Row 60: 5.8 Conjugation, 11.5 disk bounds, and 11.7 local zero strings}
\label{row:60}
\textbf{Exact row scope.} For every N >= 3, constructs a unique positive real simple zero near every sufficiently large lattice center with exponential displacement. All smallness hypotheses follow from the actual canonical remainder. Also proves conjugation and local lattice assignment. Exhaustion of all large complex zeros, orbit uniqueness and the original-denominator derivative asymptotic remain outside this row.\par\smallskip
\textbf{Manuscript labels.} \code{cor:denominator_symmetries}, \code{lem:two_term_dominance}, \code{thm:large_zero_strings}\par\smallskip
{\small\code{Komorebi.PaperIV.canonicalSignedDenominator}\par}
{\footnotesize\color{muted}definition; \code{KomorebiPaperIV/SignedZeroLocalization.lean}, line 20.\par}\smallskip
{\small\code{Komorebi.PaperIV.signedZeroFrequency}\par}
{\footnotesize\color{muted}definition; \code{KomorebiPaperIV/SignedZeroLocalization.lean}, line 23.\par}\smallskip
{\small\code{Komorebi.PaperIV.signedZeroCenter}\par}
{\footnotesize\color{muted}definition; \code{KomorebiPaperIV/SignedZeroLocalization.lean}, line 26.\par}\smallskip
{\small\code{Komorebi.PaperIV.signed_sine_pos}\par}
{\footnotesize\color{muted}theorem; \code{KomorebiPaperIV/SignedZeroLocalization.lean}, line 29.\par}\smallskip
{\small\code{Komorebi.PaperIV.signedZeroFrequency_ne_zero}\par}
{\footnotesize\color{muted}theorem; \code{KomorebiPaperIV/SignedZeroLocalization.lean}, line 35.\par}\smallskip
{\small\code{Komorebi.PaperIV.signedZeroCenter_cancellation}\par}
{\footnotesize\color{muted}theorem; \code{KomorebiPaperIV/SignedZeroLocalization.lean}, line 39.\par}\smallskip
{\small\code{Komorebi.PaperIV.canonicalSignedRemainder_disk_bound}\par}
{\footnotesize\color{muted}theorem; \code{KomorebiPaperIV/SignedZeroLocalization.lean}, line 55.\par}\smallskip
{\small\code{Komorebi.PaperIV.signedRemainder_eventually_small_on_disks}\par}
{\footnotesize\color{muted}theorem; \code{KomorebiPaperIV/SignedZeroLocalization.lean}, line 80.\par}\smallskip
{\small\code{Komorebi.PaperIV.canonicalSignedDenominator_conj}\par}
{\footnotesize\color{muted}theorem; \code{KomorebiPaperIV/SignedZeroLocalization.lean}, line 119.\par}\smallskip
{\small\code{Komorebi.PaperIV.canonicalSignedDenominator_eventual_local_zero}\par}
{\footnotesize\color{muted}theorem; \code{KomorebiPaperIV/SignedZeroLocalization.lean}, line 134.\par}\smallskip
{\small\code{Komorebi.PaperIV.canonicalSignedDenominator_unique_zero_real}\par}
{\footnotesize\color{muted}theorem; \code{KomorebiPaperIV/SignedZeroLocalization.lean}, line 161.\par}\smallskip
{\small\code{Komorebi.PaperIV.canonicalSignedDenominator_local_zero_simple}\par}
{\footnotesize\color{muted}theorem; \code{KomorebiPaperIV/SignedZeroLocalization.lean}, line 175.\par}\smallskip
{\small\code{Komorebi.PaperIV.canonicalSignedDenominator_eventual_positive_simple_zeros}\par}
{\footnotesize\color{muted}theorem; \code{KomorebiPaperIV/SignedZeroLocalization.lean}, line 206.\par}\smallskip
{\small\code{Komorebi.PaperIV.canonicalSignedDenominator_local_zero_exponential_distance}\par}
{\footnotesize\color{muted}theorem; \code{KomorebiPaperIV/SignedZeroLocalization.lean}, line 248.\par}\smallskip
{\small\code{Komorebi.PaperIV.canonicalSignedDenominator_local_strings}\par}
{\footnotesize\color{muted}theorem; \code{KomorebiPaperIV/SignedZeroLocalization.lean}, line 273.\par}\smallskip
{\small\code{Komorebi.PaperIV.canonicalSignedDenominator_zero_lattice_distance}\par}
{\footnotesize\color{muted}theorem; \code{KomorebiPaperIV/SignedZeroLocalization.lean}, line 309.\par}\smallskip
\Needspace{7\baselineskip}
\subsection{Row 61: 11.8 Leading correction and 11.7 normalized derivative}
\label{row:61}
\textbf{Exact row scope.} Actual signed-remainder leading correction has an explicit quadratic envelope error, exactly a constant times exp(-2 gamma\_N t). Proves normalized derivative error. Assumes the stated envelope smallness and local zero membership. Does not cover global exhaustion, derivative sign, reality of the displayed correction, or the full original-denominator asymptotic.\par\smallskip
\textbf{Manuscript labels.} \code{thm:first_zero_correction}, \code{thm:large_zero_strings}\par\smallskip
{\small\code{Komorebi.PaperIV.signedRemainderScale}\par}
{\footnotesize\color{muted}definition; \code{KomorebiPaperIV/SignedZeroCorrection.lean}, line 14.\par}\smallskip
{\small\code{Komorebi.PaperIV.signedRemainderEnvelope}\par}
{\footnotesize\color{muted}definition; \code{KomorebiPaperIV/SignedZeroCorrection.lean}, line 18.\par}\smallskip
{\small\code{Komorebi.PaperIV.signedRemainderEnvelope_pos}\par}
{\footnotesize\color{muted}theorem; \code{KomorebiPaperIV/SignedZeroCorrection.lean}, line 21.\par}\smallskip
{\small\code{Komorebi.PaperIV.signedRemainderEnvelope_bounds}\par}
{\footnotesize\color{muted}theorem; \code{KomorebiPaperIV/SignedZeroCorrection.lean}, line 27.\par}\smallskip
{\small\code{Komorebi.PaperIV.canonicalSignedDenominator_first_correction}\par}
{\footnotesize\color{muted}theorem; \code{KomorebiPaperIV/SignedZeroCorrection.lean}, line 39.\par}\smallskip
{\small\code{Komorebi.PaperIV.signedRemainderEnvelope_sq}\par}
{\footnotesize\color{muted}theorem; \code{KomorebiPaperIV/SignedZeroCorrection.lean}, line 62.\par}\smallskip
{\small\code{Komorebi.PaperIV.canonicalSignedDenominator_normalized_derivative_error}\par}
{\footnotesize\color{muted}theorem; \code{KomorebiPaperIV/SignedZeroCorrection.lean}, line 71.\par}\smallskip
\Needspace{7\baselineskip}
\subsection{Row 62: 11.5 Exact sector confinement}
\label{row:62}
\textbf{Exact row scope.} Exact Cartesian sector and explicit strip width log(N-1)/(2 sin(pi/N)) for every sector zero; derived finite polygon projection inequalities and one-term normalization. No large-radius or projection-gap certificate is assumed.\par\smallskip
\textbf{Manuscript labels.} \code{lem:two_term_dominance}\par\smallskip
{\small\code{Komorebi.PaperIV.signedFundamentalSector}\par}
{\footnotesize\color{muted}definition; \code{KomorebiPaperIV/SignedSectorConfinement.lean}, line 17.\par}\smallskip
{\small\code{Komorebi.PaperIV.signed_cosine_pos}\par}
{\footnotesize\color{muted}theorem; \code{KomorebiPaperIV/SignedSectorConfinement.lean}, line 20.\par}\smallskip
{\small\code{Komorebi.PaperIV.cosine_arc_upper}\par}
{\footnotesize\color{muted}theorem; \code{KomorebiPaperIV/SignedSectorConfinement.lean}, line 28.\par}\smallskip
{\small\code{Komorebi.PaperIV.canonicalSignedVertex_upper_boundary}\par}
{\footnotesize\color{muted}theorem; \code{KomorebiPaperIV/SignedSectorConfinement.lean}, line 38.\par}\smallskip
{\small\code{Komorebi.PaperIV.canonicalSignedVertex_upper_projection}\par}
{\footnotesize\color{muted}theorem; \code{KomorebiPaperIV/SignedSectorConfinement.lean}, line 65.\par}\smallskip
{\small\code{Komorebi.PaperIV.canonicalSignedVertex_sector_gap}\par}
{\footnotesize\color{muted}theorem; \code{KomorebiPaperIV/SignedSectorConfinement.lean}, line 94.\par}\smallskip
{\small\code{Komorebi.PaperIV.finite_one_term_normalization}\par}
{\footnotesize\color{muted}theorem; \code{KomorebiPaperIV/SignedSectorConfinement.lean}, line 120.\par}\smallskip
{\small\code{Komorebi.PaperIV.finite_exponential_zero_balance}\par}
{\footnotesize\color{muted}theorem; \code{KomorebiPaperIV/SignedSectorConfinement.lean}, line 131.\par}\smallskip
{\small\code{Komorebi.PaperIV.signedSectorWidth}\par}
{\footnotesize\color{muted}definition; \code{KomorebiPaperIV/SignedSectorConfinement.lean}, line 147.\par}\smallskip
{\small\code{Komorebi.PaperIV.canonicalSignedDenominator_upper_sector_zero_bound}\par}
{\footnotesize\color{muted}theorem; \code{KomorebiPaperIV/SignedSectorConfinement.lean}, line 150.\par}\smallskip
{\small\code{Komorebi.PaperIV.canonicalSignedDenominator_sector_zero_bound}\par}
{\footnotesize\color{muted}theorem; \code{KomorebiPaperIV/SignedSectorConfinement.lean}, line 173.\par}\smallskip
\Needspace{7\baselineskip}
\subsection{Row 63: 11.7 Sector exhaustion}
\label{row:63}
\textbf{Exact row scope.} Uniform strip smallness is derived from the actual remainder. Every sufficiently large sector zero belongs to a natural-index local disk, is positive real and simple, and is the only denominator zero in that disk.\par\smallskip
\textbf{Manuscript labels.} \code{thm:large_zero_strings}\par\smallskip
{\small\code{Komorebi.PaperIV.signedSectorWidth_pos}\par}
{\footnotesize\color{muted}theorem; \code{KomorebiPaperIV/SignedSectorExhaustion.lean}, line 13.\par}\smallskip
{\small\code{Komorebi.PaperIV.signedRemainder_eventually_small_on_strip}\par}
{\footnotesize\color{muted}theorem; \code{KomorebiPaperIV/SignedSectorExhaustion.lean}, line 18.\par}\smallskip
{\small\code{Komorebi.PaperIV.signed_integer_center}\par}
{\footnotesize\color{muted}theorem; \code{KomorebiPaperIV/SignedSectorExhaustion.lean}, line 32.\par}\smallskip
{\small\code{Komorebi.PaperIV.canonicalSignedDenominator_sector_exhaustion}\par}
{\footnotesize\color{muted}theorem; \code{KomorebiPaperIV/SignedSectorExhaustion.lean}, line 43.\par}\smallskip
\Needspace{7\baselineskip}
\subsection{Row 64: 5.8 Rotation and 11.7 global exhaustion}
\label{row:64}
\textbf{Exact row scope.} Actual entire rotation invariance under every Nth root of unity; existence of a rotation into the fundamental sector with a real-part/norm bound; every sufficiently large complex zero rotates into a constructed positive real simple local zero.\par\smallskip
\textbf{Manuscript labels.} \code{cor:denominator_symmetries}, \code{thm:large_zero_strings}\par\smallskip
{\small\code{Komorebi.PaperIV.canonicalSignedDenominator_rotation}\par}
{\footnotesize\color{muted}theorem; \code{KomorebiPaperIV/SignedGlobalExhaustion.lean}, line 14.\par}\smallskip
{\small\code{Komorebi.PaperIV.signedFundamentalSector_polar}\par}
{\footnotesize\color{muted}theorem; \code{KomorebiPaperIV/SignedGlobalExhaustion.lean}, line 32.\par}\smallskip
{\small\code{Komorebi.PaperIV.exists_signed_sector_rotation}\par}
{\footnotesize\color{muted}theorem; \code{KomorebiPaperIV/SignedGlobalExhaustion.lean}, line 73.\par}\smallskip
{\small\code{Komorebi.PaperIV.canonicalSignedDenominator_global_exhaustion}\par}
{\footnotesize\color{muted}theorem; \code{KomorebiPaperIV/SignedGlobalExhaustion.lean}, line 110.\par}\smallskip
\Needspace{7\baselineskip}
\subsection{Row 65: 11.7 Rotation, multiplicity, and orbit parameters}
\label{row:65}
\textbf{Exact row scope.} Exact frequency norm, rotation of analytic multiplicity and derivative, uniqueness of a unit-modulus positive real rotation, disjoint lattice disks, and global exhaustion using the canonical finite powers of omega. The full sequence/Big-O presentation is not encoded as a single theorem.\par\smallskip
\textbf{Manuscript labels.} \code{thm:large_zero_strings}\par\smallskip
{\small\code{Komorebi.PaperIV.signedZeroFrequency_norm}\par}
{\footnotesize\color{muted}theorem; \code{KomorebiPaperIV/SignedZeroOrbits.lean}, line 11.\par}\smallskip
{\small\code{Komorebi.PaperIV.canonicalSignedDenominator_order_rotation}\par}
{\footnotesize\color{muted}theorem; \code{KomorebiPaperIV/SignedZeroOrbits.lean}, line 19.\par}\smallskip
{\small\code{Komorebi.PaperIV.canonicalSignedDenominator_derivative_rotation}\par}
{\footnotesize\color{muted}theorem; \code{KomorebiPaperIV/SignedZeroOrbits.lean}, line 34.\par}\smallskip
{\small\code{Komorebi.PaperIV.positive_real_rotation_unique}\par}
{\footnotesize\color{muted}theorem; \code{KomorebiPaperIV/SignedZeroOrbits.lean}, line 50.\par}\smallskip
{\small\code{Komorebi.PaperIV.signedZeroCenter_disk_index_unique}\par}
{\footnotesize\color{muted}theorem; \code{KomorebiPaperIV/SignedZeroOrbits.lean}, line 66.\par}\smallskip
{\small\code{Komorebi.PaperIV.canonicalSignedDenominator_large_zero_orbits}\par}
{\footnotesize\color{muted}theorem; \code{KomorebiPaperIV/SignedZeroOrbits.lean}, line 102.\par}\smallskip
\Needspace{7\baselineskip}
\subsection{Row 66: 11.7 Original denominator derivative}
\label{row:66}
\textbf{Exact row scope.} Exact product-rule bridge at zeros, quantified relative error for the actual derivative, exact leading modulus (2 sin(pi/N)/N) exp(cos(pi/N) x), and a modulus-ratio error bound. These explicit estimates assume the stated local disk membership and center threshold; signs and the final sequence-level Big-O packaging remain outside this row.\par\smallskip
\textbf{Manuscript labels.} \code{thm:large_zero_strings}\par\smallskip
{\small\code{Komorebi.PaperIV.signedNormalizationFactor}\par}
{\footnotesize\color{muted}definition; \code{KomorebiPaperIV/SignedDerivativeAsymptotics.lean}, line 13.\par}\smallskip
{\small\code{Komorebi.PaperIV.signedNormalizationFactor_ne_zero}\par}
{\footnotesize\color{muted}theorem; \code{KomorebiPaperIV/SignedDerivativeAsymptotics.lean}, line 16.\par}\smallskip
{\small\code{Komorebi.PaperIV.canonicalSignedDenominator_derivative_normalization}\par}
{\footnotesize\color{muted}theorem; \code{KomorebiPaperIV/SignedDerivativeAsymptotics.lean}, line 22.\par}\smallskip
{\small\code{Komorebi.PaperIV.canonicalSignedDenominator_relative_derivative_error}\par}
{\footnotesize\color{muted}theorem; \code{KomorebiPaperIV/SignedDerivativeAsymptotics.lean}, line 47.\par}\smallskip
{\small\code{Komorebi.PaperIV.signed_derivative_leading_norm}\par}
{\footnotesize\color{muted}theorem; \code{KomorebiPaperIV/SignedDerivativeAsymptotics.lean}, line 74.\par}\smallskip
{\small\code{Komorebi.PaperIV.canonicalSignedDenominator_derivative_modulus_error}\par}
{\footnotesize\color{muted}theorem; \code{KomorebiPaperIV/SignedDerivativeAsymptotics.lean}, line 86.\par}\smallskip
\Needspace{7\baselineskip}
\subsection{Row 67: 11.5 Full two-term dominance lemma}
\label{row:67}
\textbf{Exact row scope.} Full argument-sector confinement, exact normalization, and a common positive constant for derivative jets q=0,1,2. Every constant is derived; Y=log(N-1)/(2 sin(pi/N)), R=1, C=4(N-2) exp(2Y).\par\smallskip
\textbf{Manuscript labels.} \code{lem:two_term_dominance}\par\smallskip
{\small\code{Komorebi.PaperIV.signedFundamentalSector_of_arg}\par}
{\footnotesize\color{muted}theorem; \code{KomorebiPaperIV/SignedDominance.lean}, line 11.\par}\smallskip
{\small\code{Komorebi.PaperIV.canonicalSignedRemainder_three_jet_bound}\par}
{\footnotesize\color{muted}theorem; \code{KomorebiPaperIV/SignedDominance.lean}, line 17.\par}\smallskip
{\small\code{Komorebi.PaperIV.canonicalSignedDenominator_two_term_dominance}\par}
{\footnotesize\color{muted}theorem; \code{KomorebiPaperIV/SignedDominance.lean}, line 33.\par}\smallskip
\Needspace{7\baselineskip}
\subsection{Row 68: 11.11 Eventual simplicity and exact rays}
\label{row:68}
\textbf{Exact row scope.} Full statement for every N >= 2: compact zero sets are finite, large zeros are simple and lie on canonical positive rays, and the exceptional set is finite.\par\smallskip
\textbf{Manuscript labels.} \code{cor:eventual_simplicity}\par\smallskip
{\small\code{Komorebi.PaperIV.canonicalSignedDenominator_zero}\par}
{\footnotesize\color{muted}theorem; \code{KomorebiPaperIV/SignedFiniteExceptions.lean}, line 12.\par}\smallskip
{\small\code{Komorebi.PaperIV.canonicalSignedDenominator_finite_zeros_on_compact}\par}
{\footnotesize\color{muted}theorem; \code{KomorebiPaperIV/SignedFiniteExceptions.lean}, line 18.\par}\smallskip
{\small\code{Komorebi.PaperIV.onSignedZeroRay}\par}
{\footnotesize\color{muted}definition; \code{KomorebiPaperIV/SignedFiniteExceptions.lean}, line 32.\par}\smallskip
{\small\code{Komorebi.PaperIV.canonicalSignedDenominator_eventually_simple_on_rays}\par}
{\footnotesize\color{muted}theorem; \code{KomorebiPaperIV/SignedFiniteExceptions.lean}, line 35.\par}\smallskip
{\small\code{Komorebi.PaperIV.canonicalSignedDenominator_finite_exceptional_zeros}\par}
{\footnotesize\color{muted}theorem; \code{KomorebiPaperIV/SignedFiniteExceptions.lean}, line 68.\par}\smallskip
\Needspace{7\baselineskip}
\subsection{Row 69: 11.13 Entire reduced denominator and multiplicity}
\label{row:69}
\textbf{Exact row scope.} The exact complex series is entire for N > 0. Convergent reindexing proves D\_N(z)=Delta\_N(z\textasciicircum{}N); the power substitution preserves analytic multiplicity away from zero.\par\smallskip
\textbf{Manuscript labels.} \code{cor:reduced_zeros}\par\smallskip
{\small\code{Komorebi.PaperIV.entireReducedDenominator}\par}
{\footnotesize\color{muted}definition; \code{KomorebiPaperIV/ReducedEntire.lean}, line 14.\par}\smallskip
{\small\code{Komorebi.PaperIV.reducedDenominator_coeff_complex}\par}
{\footnotesize\color{muted}theorem; \code{KomorebiPaperIV/ReducedEntire.lean}, line 17.\par}\smallskip
{\small\code{Komorebi.PaperIV.reducedDenominator_coeff_norm}\par}
{\footnotesize\color{muted}theorem; \code{KomorebiPaperIV/ReducedEntire.lean}, line 22.\par}\smallskip
{\small\code{Komorebi.PaperIV.reducedDenominator_entire_radius}\par}
{\footnotesize\color{muted}theorem; \code{KomorebiPaperIV/ReducedEntire.lean}, line 28.\par}\smallskip
{\small\code{Komorebi.PaperIV.analyticAt_entireReducedDenominator}\par}
{\footnotesize\color{muted}theorem; \code{KomorebiPaperIV/ReducedEntire.lean}, line 38.\par}\smallskip
{\small\code{Komorebi.PaperIV.hasSum_entireReducedDenominator}\par}
{\footnotesize\color{muted}theorem; \code{KomorebiPaperIV/ReducedEntire.lean}, line 43.\par}\smallskip
{\small\code{Komorebi.PaperIV.entireReducedDenominator_zero}\par}
{\footnotesize\color{muted}theorem; \code{KomorebiPaperIV/ReducedEntire.lean}, line 50.\par}\smallskip
{\small\code{Komorebi.PaperIV.canonicalSignedDenominator_eq_reduced_power}\par}
{\footnotesize\color{muted}theorem; \code{KomorebiPaperIV/ReducedEntire.lean}, line 54.\par}\smallskip
{\small\code{Komorebi.PaperIV.canonicalSignedDenominator_order_eq_reduced}\par}
{\footnotesize\color{muted}theorem; \code{KomorebiPaperIV/ReducedEntire.lean}, line 73.\par}\smallskip
\Needspace{7\baselineskip}
\subsection{Row 70: 11.10 Order two and 11.13 reduced zeros}
\label{row:70}
\textbf{Exact row scope.} Finite exceptional reduced zeros, eventual positivity and simplicity for every N >= 2, and the exact positive half-integer squares with multiplicity one for N=2.\par\smallskip
\textbf{Manuscript labels.} \code{cor:case_two_zeros}, \code{cor:reduced_zeros}\par\smallskip
{\small\code{Komorebi.PaperIV.exists_complex_nth_root}\par}
{\footnotesize\color{muted}theorem; \code{KomorebiPaperIV/ReducedZeros.lean}, line 11.\par}\smallskip
{\small\code{Komorebi.PaperIV.entireReducedDenominator_zero_ne_zero}\par}
{\footnotesize\color{muted}theorem; \code{KomorebiPaperIV/ReducedZeros.lean}, line 18.\par}\smallskip
{\small\code{Komorebi.PaperIV.entireReducedDenominator_finite_zeros_on_compact}\par}
{\footnotesize\color{muted}theorem; \code{KomorebiPaperIV/ReducedZeros.lean}, line 22.\par}\smallskip
{\small\code{Komorebi.PaperIV.onSignedZeroRay_power_positive}\par}
{\footnotesize\color{muted}theorem; \code{KomorebiPaperIV/ReducedZeros.lean}, line 35.\par}\smallskip
{\small\code{Komorebi.PaperIV.entireReducedDenominator_finite_exceptional_zeros}\par}
{\footnotesize\color{muted}theorem; \code{KomorebiPaperIV/ReducedZeros.lean}, line 44.\par}\smallskip
{\small\code{Komorebi.PaperIV.entireReducedDenominator_eventually_positive_simple}\par}
{\footnotesize\color{muted}theorem; \code{KomorebiPaperIV/ReducedZeros.lean}, line 68.\par}\smallskip
{\small\code{Komorebi.PaperIV.entireReducedDenominator_two_zero_iff}\par}
{\footnotesize\color{muted}theorem; \code{KomorebiPaperIV/ReducedZeros.lean}, line 92.\par}\smallskip
{\small\code{Komorebi.PaperIV.entireReducedDenominator_two_zero_order_one}\par}
{\footnotesize\color{muted}theorem; \code{KomorebiPaperIV/ReducedZeros.lean}, line 109.\par}\smallskip
\Needspace{7\baselineskip}
\subsection{Row 71: 11.7 Named strings and 11.13 reduced displacement}
\label{row:71}
\textbf{Exact row scope.} Actual choice-defined positive sequences have local and global exhaustion and tail injectivity for N >= 3. The reduced sequence is the Nth power and satisfies the explicit C t\_m\textasciicircum{}(N-1) exp(-gamma t\_m) error. Definitions use a documented fallback only before the proved tail. Together with rows 69-70 this covers all clauses of 11.13; the remaining 11.7 assembly is supplied by row 81. Alternating derivative signs belong to 11.8 and are now supplied by rows 82--83.\par\smallskip
\textbf{Manuscript labels.} \code{cor:reduced_zeros}, \code{thm:large_zero_strings}\par\smallskip
{\small\code{Komorebi.PaperIV.isLocalizedPositiveZero}\par}
{\footnotesize\color{muted}definition; \code{KomorebiPaperIV/SignedZeroSequence.lean}, line 12.\par}\smallskip
{\small\code{Komorebi.PaperIV.signedPositiveZero}\par}
{\footnotesize\color{muted}definition; \code{KomorebiPaperIV/SignedZeroSequence.lean}, line 17.\par}\smallskip
{\small\code{Komorebi.PaperIV.signedPositiveZero_spec}\par}
{\footnotesize\color{muted}theorem; \code{KomorebiPaperIV/SignedZeroSequence.lean}, line 21.\par}\smallskip
{\small\code{Komorebi.PaperIV.signedPositiveZero_eq_of_unique}\par}
{\footnotesize\color{muted}theorem; \code{KomorebiPaperIV/SignedZeroSequence.lean}, line 27.\par}\smallskip
{\small\code{Komorebi.PaperIV.signedPositiveZero_local_strings}\par}
{\footnotesize\color{muted}theorem; \code{KomorebiPaperIV/SignedZeroSequence.lean}, line 38.\par}\smallskip
{\small\code{Komorebi.PaperIV.signedPositiveZero_global_strings}\par}
{\footnotesize\color{muted}theorem; \code{KomorebiPaperIV/SignedZeroSequence.lean}, line 59.\par}\smallskip
{\small\code{Komorebi.PaperIV.signedZeroCenter_ge_radius}\par}
{\footnotesize\color{muted}theorem; \code{KomorebiPaperIV/SignedZeroSequence.lean}, line 72.\par}\smallskip
{\small\code{Komorebi.PaperIV.signedPositiveZero_pos}\par}
{\footnotesize\color{muted}theorem; \code{KomorebiPaperIV/SignedZeroSequence.lean}, line 83.\par}\smallskip
{\small\code{Komorebi.PaperIV.signedPositiveZero_injective_on_tail}\par}
{\footnotesize\color{muted}theorem; \code{KomorebiPaperIV/SignedZeroSequence.lean}, line 93.\par}\smallskip
{\small\code{Komorebi.PaperIV.reducedPositiveZero}\par}
{\footnotesize\color{muted}definition; \code{KomorebiPaperIV/SignedZeroSequence.lean}, line 103.\par}\smallskip
{\small\code{Komorebi.PaperIV.reducedPositiveZero_global_strings}\par}
{\footnotesize\color{muted}theorem; \code{KomorebiPaperIV/SignedZeroSequence.lean}, line 106.\par}\smallskip
{\small\code{Komorebi.PaperIV.reducedPositiveZero_injective_on_tail}\par}
{\footnotesize\color{muted}theorem; \code{KomorebiPaperIV/SignedZeroSequence.lean}, line 127.\par}\smallskip
{\small\code{Komorebi.PaperIV.reducedPositiveZero_local_strings}\par}
{\footnotesize\color{muted}theorem; \code{KomorebiPaperIV/SignedZeroSequence.lean}, line 138.\par}\smallskip
\Needspace{7\baselineskip}
\subsection{Row 72: 11.15 Sharp summability thresholds}
\label{row:72}
\textbf{Exact row scope.} Actual positive sequences satisfy tau\_m \textasciitilde{} t\_m and lambda\_m \textasciitilde{} t\_m\textasciicircum{}N. Reciprocal real powers are summable exactly for s>1 and N*s>1 respectively, for N >= 3. This row does not assemble the full complex zero multiset, exceptional multiplicities, or the N=2 exponent statement.\par\smallskip
\textbf{Manuscript labels.} \code{cor:exponent_of_convergence}\par\smallskip
{\small\code{Komorebi.PaperIV.signedZeroCenter_tendsto_atTop}\par}
{\footnotesize\color{muted}theorem; \code{KomorebiPaperIV/ZeroSequenceAsymptotics.lean}, line 14.\par}\smallskip
{\small\code{Komorebi.PaperIV.signedPositiveZero_sub_center_isBigO_one}\par}
{\footnotesize\color{muted}theorem; \code{KomorebiPaperIV/ZeroSequenceAsymptotics.lean}, line 19.\par}\smallskip
{\small\code{Komorebi.PaperIV.signedPositiveZero_equivalent_center}\par}
{\footnotesize\color{muted}theorem; \code{KomorebiPaperIV/ZeroSequenceAsymptotics.lean}, line 29.\par}\smallskip
{\small\code{Komorebi.PaperIV.signedPositiveZero_tendsto_atTop}\par}
{\footnotesize\color{muted}theorem; \code{KomorebiPaperIV/ZeroSequenceAsymptotics.lean}, line 40.\par}\smallskip
{\small\code{Komorebi.PaperIV.reducedPositiveZero_equivalent_center_power}\par}
{\footnotesize\color{muted}theorem; \code{KomorebiPaperIV/ZeroSequenceAsymptotics.lean}, line 45.\par}\smallskip
{\small\code{Komorebi.PaperIV.signedPositiveZero_reciprocal_rpow_summable_iff}\par}
{\footnotesize\color{muted}theorem; \code{KomorebiPaperIV/ZeroSequenceAsymptotics.lean}, line 50.\par}\smallskip
{\small\code{Komorebi.PaperIV.reducedPositiveZero_reciprocal_rpow_summable_iff}\par}
{\footnotesize\color{muted}theorem; \code{KomorebiPaperIV/ZeroSequenceAsymptotics.lean}, line 67.\par}\smallskip
{\small\code{Komorebi.PaperIV.reducedPositiveZero_reciprocal_summable}\par}
{\footnotesize\color{muted}theorem; \code{KomorebiPaperIV/ZeroSequenceAsymptotics.lean}, line 76.\par}\smallskip
\Needspace{7\baselineskip}
\subsection{Row 73: 11.14 Positive-tail counting}
\label{row:73}
\textbf{Exact row scope.} A finite lattice-window counting lemma is applied to the actual positive zero tail, giving density sin(pi/N)/pi with uniform bounded error for N >= 3. This row is limited to tails; rows 77--78 supply both full multiplicity-weighted counts.\par\smallskip
\textbf{Manuscript labels.} \code{thm:zero_counting}\par\smallskip
{\small\code{Komorebi.PaperIV.sequenceCounting}\par}
{\footnotesize\color{muted}definition; \code{KomorebiPaperIV/ZeroSequenceCounting.lean}, line 14.\par}\smallskip
{\small\code{Komorebi.PaperIV.lattice_window_counting}\par}
{\footnotesize\color{muted}theorem; \code{KomorebiPaperIV/ZeroSequenceCounting.lean}, line 17.\par}\smallskip
{\small\code{Komorebi.PaperIV.signedPositiveZero_lattice_windows}\par}
{\footnotesize\color{muted}theorem; \code{KomorebiPaperIV/ZeroSequenceCounting.lean}, line 49.\par}\smallskip
{\small\code{Komorebi.PaperIV.signedPositiveZero_tail_counting}\par}
{\footnotesize\color{muted}theorem; \code{KomorebiPaperIV/ZeroSequenceCounting.lean}, line 70.\par}\smallskip
\Needspace{7\baselineskip}
\subsection{Row 74: 12.1 Convergent entire tail product}
\label{row:74}
\textbf{Exact row scope.} Reciprocal summability proves compact-uniform convergence and entire analyticity of the genus-zero product over any shifted chosen positive sequence for N >= 3. Its zero set is exactly that sequence. A sufficiently late shift consists of actual simple reduced zeros by row 71. The finite exceptional factor and equality with Delta\_N are not asserted.\par\smallskip
\textbf{Manuscript labels.} \code{thm:genus_zero_product}\par\smallskip
{\small\code{Komorebi.PaperIV.reducedTailProduct}\par}
{\footnotesize\color{muted}definition; \code{KomorebiPaperIV/ReducedCanonicalProduct.lean}, line 16.\par}\smallskip
{\small\code{Komorebi.PaperIV.reducedPositiveZero_pos}\par}
{\footnotesize\color{muted}theorem; \code{KomorebiPaperIV/ReducedCanonicalProduct.lean}, line 19.\par}\smallskip
{\small\code{Komorebi.PaperIV.reducedTail_reciprocal_summable}\par}
{\footnotesize\color{muted}theorem; \code{KomorebiPaperIV/ReducedCanonicalProduct.lean}, line 22.\par}\smallskip
{\small\code{Komorebi.PaperIV.reducedTail_factors_norm_summable}\par}
{\footnotesize\color{muted}theorem; \code{KomorebiPaperIV/ReducedCanonicalProduct.lean}, line 26.\par}\smallskip
{\small\code{Komorebi.PaperIV.reducedTailProduct_multipliable}\par}
{\footnotesize\color{muted}theorem; \code{KomorebiPaperIV/ReducedCanonicalProduct.lean}, line 32.\par}\smallskip
{\small\code{Komorebi.PaperIV.reducedTailProduct_uniform_on_compact}\par}
{\footnotesize\color{muted}theorem; \code{KomorebiPaperIV/ReducedCanonicalProduct.lean}, line 36.\par}\smallskip
{\small\code{Komorebi.PaperIV.reducedTailProduct_locally_uniform}\par}
{\footnotesize\color{muted}theorem; \code{KomorebiPaperIV/ReducedCanonicalProduct.lean}, line 50.\par}\smallskip
{\small\code{Komorebi.PaperIV.analyticAt_reducedTailProduct}\par}
{\footnotesize\color{muted}theorem; \code{KomorebiPaperIV/ReducedCanonicalProduct.lean}, line 57.\par}\smallskip
{\small\code{Komorebi.PaperIV.reducedTailProduct_zero}\par}
{\footnotesize\color{muted}theorem; \code{KomorebiPaperIV/ReducedCanonicalProduct.lean}, line 67.\par}\smallskip
{\small\code{Komorebi.PaperIV.reducedTailProduct_ne_zero}\par}
{\footnotesize\color{muted}theorem; \code{KomorebiPaperIV/ReducedCanonicalProduct.lean}, line 70.\par}\smallskip
{\small\code{Komorebi.PaperIV.reducedTailProduct_zero_iff}\par}
{\footnotesize\color{muted}theorem; \code{KomorebiPaperIV/ReducedCanonicalProduct.lean}, line 84.\par}\smallskip
\Needspace{7\baselineskip}
\subsection{Row 75: 11.14 Analytic multiplicity-weighted counting}
\label{row:75}
\textbf{Exact row scope.} Actual disk counts sum analyticOrderNatAt over the finite zero set. The original and reduced entire denominators have finite order at every point, so the Nat-valued order loses no infinite multiplicity. Derived finite-weight comparison lemmas retain all exceptional weights.\par\smallskip
\textbf{Manuscript labels.} \code{thm:zero_counting}\par\smallskip
{\small\code{Komorebi.PaperIV.finiteWeight}\par}
{\footnotesize\color{muted}definition; \code{KomorebiPaperIV/WeightedZeroCounts.lean}, line 13.\par}\smallskip
{\small\code{Komorebi.PaperIV.finiteWeight_of_finite}\par}
{\footnotesize\color{muted}theorem; \code{KomorebiPaperIV/WeightedZeroCounts.lean}, line 17.\par}\smallskip
{\small\code{Komorebi.PaperIV.finiteWeight_mono}\par}
{\footnotesize\color{muted}theorem; \code{KomorebiPaperIV/WeightedZeroCounts.lean}, line 23.\par}\smallskip
{\small\code{Komorebi.PaperIV.finiteWeight_union_le}\par}
{\footnotesize\color{muted}theorem; \code{KomorebiPaperIV/WeightedZeroCounts.lean}, line 29.\par}\smallskip
{\small\code{Komorebi.PaperIV.finiteWeight_eq_ncard}\par}
{\footnotesize\color{muted}theorem; \code{KomorebiPaperIV/WeightedZeroCounts.lean}, line 38.\par}\smallskip
{\small\code{Komorebi.PaperIV.finiteWeight_between_tail_and_exceptions}\par}
{\footnotesize\color{muted}theorem; \code{KomorebiPaperIV/WeightedZeroCounts.lean}, line 47.\par}\smallskip
{\small\code{Komorebi.PaperIV.zeroDisk}\par}
{\footnotesize\color{muted}definition; \code{KomorebiPaperIV/WeightedZeroCounts.lean}, line 60.\par}\smallskip
{\small\code{Komorebi.PaperIV.analyticZeroCount}\par}
{\footnotesize\color{muted}definition; \code{KomorebiPaperIV/WeightedZeroCounts.lean}, line 63.\par}\smallskip
{\small\code{Komorebi.PaperIV.entire_analyticOrderAt_ne_top}\par}
{\footnotesize\color{muted}theorem; \code{KomorebiPaperIV/WeightedZeroCounts.lean}, line 67.\par}\smallskip
{\small\code{Komorebi.PaperIV.canonicalSignedDenominator_order_ne_top}\par}
{\footnotesize\color{muted}theorem; \code{KomorebiPaperIV/WeightedZeroCounts.lean}, line 75.\par}\smallskip
{\small\code{Komorebi.PaperIV.entireReducedDenominator_order_ne_top}\par}
{\footnotesize\color{muted}theorem; \code{KomorebiPaperIV/WeightedZeroCounts.lean}, line 80.\par}\smallskip
{\small\code{Komorebi.PaperIV.canonicalSignedDenominator_zeroDisk_finite}\par}
{\footnotesize\color{muted}theorem; \code{KomorebiPaperIV/WeightedZeroCounts.lean}, line 85.\par}\smallskip
{\small\code{Komorebi.PaperIV.entireReducedDenominator_zeroDisk_finite}\par}
{\footnotesize\color{muted}theorem; \code{KomorebiPaperIV/WeightedZeroCounts.lean}, line 90.\par}\smallskip
\Needspace{7\baselineskip}
\subsection{Row 76: 11.14 Exhaustive zero enumeration}
\label{row:76}
\textbf{Exact row scope.} Constructs a positive simple tail for every N >= 2, including the exact cosine case. Proves injectivity, lattice windows, exhaustive finite root-of-unity orbits, orbit norms and multiplicities. The structure is constructed from the proved zero theorems, never assumed in the final counting results.\par\smallskip
\textbf{Manuscript labels.} \code{thm:zero_counting}\par\smallskip
{\small\code{Komorebi.PaperIV.signedZeroEnumeration_exists_ge_three}\par}
{\footnotesize\color{muted}theorem; \code{KomorebiPaperIV/ZeroEnumeration.lean}, line 30.\par}\smallskip
{\small\code{Komorebi.PaperIV.signedZeroEnumeration_exists_two}\par}
{\footnotesize\color{muted}theorem; \code{KomorebiPaperIV/ZeroEnumeration.lean}, line 81.\par}\smallskip
{\small\code{Komorebi.PaperIV.signedZeroEnumeration_exists}\par}
{\footnotesize\color{muted}theorem; \code{KomorebiPaperIV/ZeroEnumeration.lean}, line 121.\par}\smallskip
{\small\code{Komorebi.PaperIV.canonicalZeroEnumeration}\par}
{\footnotesize\color{muted}definition; \code{KomorebiPaperIV/ZeroEnumeration.lean}, line 127.\par}\smallskip
{\small\code{Komorebi.PaperIV.SignedZeroEnumeration.orbit}\par}
{\footnotesize\color{muted}definition; \code{KomorebiPaperIV/ZeroEnumeration.lean}, line 130.\par}\smallskip
{\small\code{Komorebi.PaperIV.SignedZeroEnumeration.orbit_norm}\par}
{\footnotesize\color{muted}theorem; \code{KomorebiPaperIV/ZeroEnumeration.lean}, line 133.\par}\smallskip
{\small\code{Komorebi.PaperIV.SignedZeroEnumeration.orbit_injective}\par}
{\footnotesize\color{muted}theorem; \code{KomorebiPaperIV/ZeroEnumeration.lean}, line 139.\par}\smallskip
{\small\code{Komorebi.PaperIV.SignedZeroEnumeration.orbit_zero_simple}\par}
{\footnotesize\color{muted}theorem; \code{KomorebiPaperIV/ZeroEnumeration.lean}, line 153.\par}\smallskip
\Needspace{7\baselineskip}
\subsection{Row 77: 11.14 Full original-denominator counting}
\label{row:77}
\textbf{Exact row scope.} For every N >= 2, the actual disk count, including all exceptional analytic multiplicities, differs from (N sin(pi/N)/pi) R by a uniform bounded amount for all R >= 0.\par\smallskip
\textbf{Manuscript labels.} \code{thm:zero_counting}\par\smallskip
{\small\code{Komorebi.PaperIV.SignedZeroEnumeration.orbitDisk}\par}
{\footnotesize\color{muted}definition; \code{KomorebiPaperIV/FullZeroCounting.lean}, line 11.\par}\smallskip
{\small\code{Komorebi.PaperIV.SignedZeroEnumeration.exceptions}\par}
{\footnotesize\color{muted}definition; \code{KomorebiPaperIV/FullZeroCounting.lean}, line 14.\par}\smallskip
{\small\code{Komorebi.PaperIV.SignedZeroEnumeration.exceptions_finite}\par}
{\footnotesize\color{muted}theorem; \code{KomorebiPaperIV/FullZeroCounting.lean}, line 17.\par}\smallskip
{\small\code{Komorebi.PaperIV.SignedZeroEnumeration.orbitDisk_finite_card}\par}
{\footnotesize\color{muted}theorem; \code{KomorebiPaperIV/FullZeroCounting.lean}, line 20.\par}\smallskip
{\small\code{Komorebi.PaperIV.SignedZeroEnumeration.orbitDisk_subset}\par}
{\footnotesize\color{muted}theorem; \code{KomorebiPaperIV/FullZeroCounting.lean}, line 28.\par}\smallskip
{\small\code{Komorebi.PaperIV.SignedZeroEnumeration.zeroDisk_subset_union}\par}
{\footnotesize\color{muted}theorem; \code{KomorebiPaperIV/FullZeroCounting.lean}, line 34.\par}\smallskip
{\small\code{Komorebi.PaperIV.SignedZeroEnumeration.weighted_count_bounds}\par}
{\footnotesize\color{muted}theorem; \code{KomorebiPaperIV/FullZeroCounting.lean}, line 49.\par}\smallskip
{\small\code{Komorebi.PaperIV.canonicalSignedDenominator_zero_counting}\par}
{\footnotesize\color{muted}theorem; \code{KomorebiPaperIV/FullZeroCounting.lean}, line 64.\par}\smallskip
\Needspace{7\baselineskip}
\subsection{Row 78: 11.14 Full reduced-denominator counting}
\label{row:78}
\textbf{Exact row scope.} For every N >= 2, the actual reduced disk count differs from (sin(pi/N)/pi) R\textasciicircum{}(1/N) by a uniform bounded amount for all R >= 0, including exceptional multiplicities. Together with row 77, this fully encodes Theorem 11.14.\par\smallskip
\textbf{Manuscript labels.} \code{thm:zero_counting}\par\smallskip
{\small\code{Komorebi.PaperIV.SignedZeroEnumeration.reducedRoot}\par}
{\footnotesize\color{muted}definition; \code{KomorebiPaperIV/ReducedZeroCounting.lean}, line 11.\par}\smallskip
{\small\code{Komorebi.PaperIV.SignedZeroEnumeration.reducedDisk}\par}
{\footnotesize\color{muted}definition; \code{KomorebiPaperIV/ReducedZeroCounting.lean}, line 14.\par}\smallskip
{\small\code{Komorebi.PaperIV.SignedZeroEnumeration.reducedExceptions}\par}
{\footnotesize\color{muted}definition; \code{KomorebiPaperIV/ReducedZeroCounting.lean}, line 17.\par}\smallskip
{\small\code{Komorebi.PaperIV.SignedZeroEnumeration.reducedRoot_norm}\par}
{\footnotesize\color{muted}theorem; \code{KomorebiPaperIV/ReducedZeroCounting.lean}, line 20.\par}\smallskip
{\small\code{Komorebi.PaperIV.SignedZeroEnumeration.reducedRoot_injective}\par}
{\footnotesize\color{muted}theorem; \code{KomorebiPaperIV/ReducedZeroCounting.lean}, line 24.\par}\smallskip
{\small\code{Komorebi.PaperIV.SignedZeroEnumeration.reducedRoot_zero_simple}\par}
{\footnotesize\color{muted}theorem; \code{KomorebiPaperIV/ReducedZeroCounting.lean}, line 32.\par}\smallskip
{\small\code{Komorebi.PaperIV.SignedZeroEnumeration.reducedExceptions_finite}\par}
{\footnotesize\color{muted}theorem; \code{KomorebiPaperIV/ReducedZeroCounting.lean}, line 43.\par}\smallskip
{\small\code{Komorebi.PaperIV.SignedZeroEnumeration.reducedDisk_finite_card}\par}
{\footnotesize\color{muted}theorem; \code{KomorebiPaperIV/ReducedZeroCounting.lean}, line 47.\par}\smallskip
{\small\code{Komorebi.PaperIV.SignedZeroEnumeration.reducedDisk_subset}\par}
{\footnotesize\color{muted}theorem; \code{KomorebiPaperIV/ReducedZeroCounting.lean}, line 54.\par}\smallskip
{\small\code{Komorebi.PaperIV.SignedZeroEnumeration.reducedZeroDisk_subset_union}\par}
{\footnotesize\color{muted}theorem; \code{KomorebiPaperIV/ReducedZeroCounting.lean}, line 62.\par}\smallskip
{\small\code{Komorebi.PaperIV.SignedZeroEnumeration.reduced_weighted_count_bounds}\par}
{\footnotesize\color{muted}theorem; \code{KomorebiPaperIV/ReducedZeroCounting.lean}, line 94.\par}\smallskip
{\small\code{Komorebi.PaperIV.entireReducedDenominator_zero_counting_power}\par}
{\footnotesize\color{muted}theorem; \code{KomorebiPaperIV/ReducedZeroCounting.lean}, line 108.\par}\smallskip
{\small\code{Komorebi.PaperIV.entireReducedDenominator_zero_counting}\par}
{\footnotesize\color{muted}theorem; \code{KomorebiPaperIV/ReducedZeroCounting.lean}, line 128.\par}\smallskip
\Needspace{7\baselineskip}
\subsection{Row 79: 11.15 Full zero convergence exponents}
\label{row:79}
\textbf{Exact row scope.} The reciprocal-power series over all distinct zeros, weighted by actual analytic multiplicity, is summable exactly for s>1 for D\_N and s>1/N for Delta\_N, for every N >= 2. Both boundary divergences are included. Finite exceptional sets are removed only through a proved summability equivalence. This fully encodes Corollary 11.15.\par\smallskip
\textbf{Manuscript labels.} \code{cor:exponent_of_convergence}\par\smallskip
{\small\code{Komorebi.PaperIV.ZeroPowerSummable}\par}
{\footnotesize\color{muted}definition; \code{KomorebiPaperIV/FullZeroSummability.lean}, line 14.\par}\smallskip
{\small\code{Komorebi.PaperIV.summable_iff_comp_of_finite_complement}\par}
{\footnotesize\color{muted}theorem; \code{KomorebiPaperIV/FullZeroSummability.lean}, line 18.\par}\smallskip
{\small\code{Komorebi.PaperIV.SignedZeroEnumeration.root_equivalent_lattice}\par}
{\footnotesize\color{muted}theorem; \code{KomorebiPaperIV/FullZeroSummability.lean}, line 28.\par}\smallskip
{\small\code{Komorebi.PaperIV.SignedZeroEnumeration.root_reciprocal_rpow_summable_iff}\par}
{\footnotesize\color{muted}theorem; \code{KomorebiPaperIV/FullZeroSummability.lean}, line 46.\par}\smallskip
{\small\code{Komorebi.PaperIV.SignedZeroEnumeration.reducedRoot_reciprocal_rpow_summable_iff}\par}
{\footnotesize\color{muted}theorem; \code{KomorebiPaperIV/FullZeroSummability.lean}, line 64.\par}\smallskip
{\small\code{Komorebi.PaperIV.SignedZeroEnumeration.orbitZero}\par}
{\footnotesize\color{muted}definition; \code{KomorebiPaperIV/FullZeroSummability.lean}, line 74.\par}\smallskip
{\small\code{Komorebi.PaperIV.SignedZeroEnumeration.orbitZero_injective}\par}
{\footnotesize\color{muted}theorem; \code{KomorebiPaperIV/FullZeroSummability.lean}, line 78.\par}\smallskip
{\small\code{Komorebi.PaperIV.SignedZeroEnumeration.orbitZero_finite_complement}\par}
{\footnotesize\color{muted}theorem; \code{KomorebiPaperIV/FullZeroSummability.lean}, line 83.\par}\smallskip
{\small\code{Komorebi.PaperIV.SignedZeroEnumeration.reducedZero}\par}
{\footnotesize\color{muted}definition; \code{KomorebiPaperIV/FullZeroSummability.lean}, line 95.\par}\smallskip
{\small\code{Komorebi.PaperIV.SignedZeroEnumeration.reducedZero_injective}\par}
{\footnotesize\color{muted}theorem; \code{KomorebiPaperIV/FullZeroSummability.lean}, line 99.\par}\smallskip
{\small\code{Komorebi.PaperIV.SignedZeroEnumeration.reducedZero_finite_complement}\par}
{\footnotesize\color{muted}theorem; \code{KomorebiPaperIV/FullZeroSummability.lean}, line 104.\par}\smallskip
{\small\code{Komorebi.PaperIV.canonicalSignedDenominator_zero_power_summable_iff}\par}
{\footnotesize\color{muted}theorem; \code{KomorebiPaperIV/FullZeroSummability.lean}, line 122.\par}\smallskip
{\small\code{Komorebi.PaperIV.entireReducedDenominator_zero_power_summable_iff}\par}
{\footnotesize\color{muted}theorem; \code{KomorebiPaperIV/FullZeroSummability.lean}, line 148.\par}\smallskip
{\small\code{Komorebi.PaperIV.entireReducedDenominator_reciprocal_zeros_summable}\par}
{\footnotesize\color{muted}theorem; \code{KomorebiPaperIV/FullZeroSummability.lean}, line 164.\par}\smallskip
\Needspace{7\baselineskip}
\subsection{Row 80: 14.2 Absolute convergence and holomorphy}
\label{row:80}
\textbf{Exact row scope.} Defines the exact complex summands m(z) exp(-s log(norm z)) over every zero. Proves both sharp absolute-convergence half-planes and failure of absolute convergence on and to their left, local uniform convergence by an explicit strip majorant, and holomorphy of both sums for every N >= 2. This fully encodes Theorem 14.2; meromorphic continuation in 14.3 is not asserted.\par\smallskip
\textbf{Manuscript labels.} \code{thm:zeta_convergence}\par\smallskip
{\small\code{Komorebi.PaperIV.weightedDirichletTerm}\par}
{\footnotesize\color{muted}definition; \code{KomorebiPaperIV/ZeroZetaConvergence.lean}, line 13.\par}\smallskip
{\small\code{Komorebi.PaperIV.weightedDirichletTerm_norm}\par}
{\footnotesize\color{muted}theorem; \code{KomorebiPaperIV/ZeroZetaConvergence.lean}, line 16.\par}\smallskip
{\small\code{Komorebi.PaperIV.weightedDirichletTerm_differentiable}\par}
{\footnotesize\color{muted}theorem; \code{KomorebiPaperIV/ZeroZetaConvergence.lean}, line 26.\par}\smallskip
{\small\code{Komorebi.PaperIV.weightedDirichletTerm_strip_bound}\par}
{\footnotesize\color{muted}theorem; \code{KomorebiPaperIV/ZeroZetaConvergence.lean}, line 31.\par}\smallskip
{\small\code{Komorebi.PaperIV.weightedDirichlet_locally_uniform}\par}
{\footnotesize\color{muted}theorem; \code{KomorebiPaperIV/ZeroZetaConvergence.lean}, line 49.\par}\smallskip
{\small\code{Komorebi.PaperIV.weightedDirichlet_differentiableOn}\par}
{\footnotesize\color{muted}theorem; \code{KomorebiPaperIV/ZeroZetaConvergence.lean}, line 71.\par}\smallskip
{\small\code{Komorebi.PaperIV.absoluteZeroZetaTerm}\par}
{\footnotesize\color{muted}definition; \code{KomorebiPaperIV/ZeroZetaConvergence.lean}, line 81.\par}\smallskip
{\small\code{Komorebi.PaperIV.absoluteZeroZeta}\par}
{\footnotesize\color{muted}definition; \code{KomorebiPaperIV/ZeroZetaConvergence.lean}, line 85.\par}\smallskip
{\small\code{Komorebi.PaperIV.zero_norm_pos_of_value_at_zero}\par}
{\footnotesize\color{muted}theorem; \code{KomorebiPaperIV/ZeroZetaConvergence.lean}, line 87.\par}\smallskip
{\small\code{Komorebi.PaperIV.absoluteZeroZetaTerm_norm}\par}
{\footnotesize\color{muted}theorem; \code{KomorebiPaperIV/ZeroZetaConvergence.lean}, line 95.\par}\smallskip
{\small\code{Komorebi.PaperIV.absoluteZeroZeta_summable_norm_iff}\par}
{\footnotesize\color{muted}theorem; \code{KomorebiPaperIV/ZeroZetaConvergence.lean}, line 100.\par}\smallskip
{\small\code{Komorebi.PaperIV.absoluteZeroZeta_locally_uniform}\par}
{\footnotesize\color{muted}theorem; \code{KomorebiPaperIV/ZeroZetaConvergence.lean}, line 104.\par}\smallskip
{\small\code{Komorebi.PaperIV.absoluteZeroZeta_differentiableOn}\par}
{\footnotesize\color{muted}theorem; \code{KomorebiPaperIV/ZeroZetaConvergence.lean}, line 110.\par}\smallskip
{\small\code{Komorebi.PaperIV.canonicalSignedDenominator_zeta_absolute_convergence}\par}
{\footnotesize\color{muted}theorem; \code{KomorebiPaperIV/ZeroZetaConvergence.lean}, line 116.\par}\smallskip
{\small\code{Komorebi.PaperIV.canonicalSignedDenominator_zeta_locally_uniform}\par}
{\footnotesize\color{muted}theorem; \code{KomorebiPaperIV/ZeroZetaConvergence.lean}, line 121.\par}\smallskip
{\small\code{Komorebi.PaperIV.canonicalSignedDenominator_zeta_holomorphic}\par}
{\footnotesize\color{muted}theorem; \code{KomorebiPaperIV/ZeroZetaConvergence.lean}, line 128.\par}\smallskip
{\small\code{Komorebi.PaperIV.entireReducedDenominator_zeta_absolute_convergence}\par}
{\footnotesize\color{muted}theorem; \code{KomorebiPaperIV/ZeroZetaConvergence.lean}, line 134.\par}\smallskip
{\small\code{Komorebi.PaperIV.entireReducedDenominator_zeta_locally_uniform}\par}
{\footnotesize\color{muted}theorem; \code{KomorebiPaperIV/ZeroZetaConvergence.lean}, line 139.\par}\smallskip
{\small\code{Komorebi.PaperIV.entireReducedDenominator_zeta_holomorphic}\par}
{\footnotesize\color{muted}theorem; \code{KomorebiPaperIV/ZeroZetaConvergence.lean}, line 146.\par}\smallskip
\Needspace{7\baselineskip}
\subsection{Row 81: 11.7 Complete large zero strings}
\label{row:81}
\textbf{Exact row scope.} One theorem supplies a common tail cutoff, positive simple zeros, exponential displacement, unique global orbit/index representation, exact derivative covariance and the original derivative modulus asymptotic with exponential error measured at the actual zero. This fully encodes Theorem 11.7 for N >= 3. The alternating derivative sign in 11.8 is supplied by rows 82--83.\par\smallskip
\textbf{Manuscript labels.} \code{thm:large_zero_strings}\par\smallskip
{\small\code{Komorebi.PaperIV.signedPositiveZero_derivative_asymptotic}\par}
{\footnotesize\color{muted}theorem; \code{KomorebiPaperIV/LargeZeroCompletion.lean}, line 14.\par}\smallskip
{\small\code{Komorebi.PaperIV.signedPositiveZero_unique_global_strings}\par}
{\footnotesize\color{muted}theorem; \code{KomorebiPaperIV/LargeZeroCompletion.lean}, line 53.\par}\smallskip
{\small\code{Komorebi.PaperIV.signedPositiveZero_derivative_rotation}\par}
{\footnotesize\color{muted}theorem; \code{KomorebiPaperIV/LargeZeroCompletion.lean}, line 97.\par}\smallskip
{\small\code{Komorebi.PaperIV.canonicalSignedDenominator_large_zero_strings}\par}
{\footnotesize\color{muted}theorem; \code{KomorebiPaperIV/LargeZeroCompletion.lean}, line 114.\par}\smallskip
\Needspace{7\baselineskip}
\subsection{Row 82: 11.8 Reality and exact phase}
\label{row:82}
\textbf{Exact row scope.} Canonical signed sections and their denominator derivative are real on the real axis. The lattice phase and the exact signed leading derivative factor are proved, including the alternating power of minus one.\par\smallskip
\textbf{Manuscript labels.} \code{thm:first_zero_correction}\par\smallskip
{\small\code{Komorebi.PaperIV.entireSigned_canonical_conj}\par}
{\footnotesize\color{muted}theorem; \code{KomorebiPaperIV/SignedRealPhase.lean}, line 12.\par}\smallskip
{\small\code{Komorebi.PaperIV.entireSigned_canonical_real}\par}
{\footnotesize\color{muted}theorem; \code{KomorebiPaperIV/SignedRealPhase.lean}, line 24.\par}\smallskip
{\small\code{Komorebi.PaperIV.canonicalSignedDenominator_derivative_real}\par}
{\footnotesize\color{muted}theorem; \code{KomorebiPaperIV/SignedRealPhase.lean}, line 30.\par}\smallskip
{\small\code{Komorebi.PaperIV.signed_sine_mul_center}\par}
{\footnotesize\color{muted}theorem; \code{KomorebiPaperIV/SignedRealPhase.lean}, line 36.\par}\smallskip
{\small\code{Komorebi.PaperIV.signedZeroCenter_positive_phase}\par}
{\footnotesize\color{muted}theorem; \code{KomorebiPaperIV/SignedRealPhase.lean}, line 41.\par}\smallskip
{\small\code{Komorebi.PaperIV.signedZeroCenter_negative_phase}\par}
{\footnotesize\color{muted}theorem; \code{KomorebiPaperIV/SignedRealPhase.lean}, line 52.\par}\smallskip
{\small\code{Komorebi.PaperIV.signedDerivativeScale}\par}
{\footnotesize\color{muted}definition; \code{KomorebiPaperIV/SignedRealPhase.lean}, line 60.\par}\smallskip
{\small\code{Komorebi.PaperIV.signedDerivativeMainTerm}\par}
{\footnotesize\color{muted}definition; \code{KomorebiPaperIV/SignedRealPhase.lean}, line 63.\par}\smallskip
{\small\code{Komorebi.PaperIV.signedDerivativeScale_pos}\par}
{\footnotesize\color{muted}theorem; \code{KomorebiPaperIV/SignedRealPhase.lean}, line 66.\par}\smallskip
{\small\code{Komorebi.PaperIV.signedDerivativeMainTerm_ne_zero}\par}
{\footnotesize\color{muted}theorem; \code{KomorebiPaperIV/SignedRealPhase.lean}, line 70.\par}\smallskip
{\small\code{Komorebi.PaperIV.signed_derivative_leading_phase}\par}
{\footnotesize\color{muted}theorem; \code{KomorebiPaperIV/SignedRealPhase.lean}, line 75.\par}\smallskip
\Needspace{7\baselineskip}
\subsection{Row 83: 11.8 Full leading correction and derivative sign}
\label{row:83}
\textbf{Exact row scope.} At the actual named positive zeros, the leading correction is real and the displacement remainder is bounded by C exp(-2 gamma t). The signed derivative relative error is bounded by B exp(-gamma tau); the eventual alternating sign is also proved explicitly. Together with row 82 this fully encodes Theorem 11.8.\par\smallskip
\textbf{Manuscript labels.} \code{thm:first_zero_correction}\par\smallskip
{\small\code{Komorebi.PaperIV.signedLeadingCorrection}\par}
{\footnotesize\color{muted}definition; \code{KomorebiPaperIV/SignedCorrectionCompletion.lean}, line 12.\par}\smallskip
{\small\code{Komorebi.PaperIV.signedLeadingCorrection_eq_real_quotient}\par}
{\footnotesize\color{muted}theorem; \code{KomorebiPaperIV/SignedCorrectionCompletion.lean}, line 16.\par}\smallskip
{\small\code{Komorebi.PaperIV.signedLeadingCorrection_real}\par}
{\footnotesize\color{muted}theorem; \code{KomorebiPaperIV/SignedCorrectionCompletion.lean}, line 32.\par}\smallskip
{\small\code{Komorebi.PaperIV.signedRemainderEnvelope_center_tendsto_zero}\par}
{\footnotesize\color{muted}theorem; \code{KomorebiPaperIV/SignedCorrectionCompletion.lean}, line 38.\par}\smallskip
{\small\code{Komorebi.PaperIV.signedPositiveZero_first_correction}\par}
{\footnotesize\color{muted}theorem; \code{KomorebiPaperIV/SignedCorrectionCompletion.lean}, line 47.\par}\smallskip
{\small\code{Komorebi.PaperIV.norm_real_phase_sub_one_le}\par}
{\footnotesize\color{muted}theorem; \code{KomorebiPaperIV/SignedCorrectionCompletion.lean}, line 70.\par}\smallskip
{\small\code{Komorebi.PaperIV.relative_error_through_unit_phase}\par}
{\footnotesize\color{muted}theorem; \code{KomorebiPaperIV/SignedCorrectionCompletion.lean}, line 78.\par}\smallskip
{\small\code{Komorebi.PaperIV.signedPositiveZero_signed_derivative_asymptotic}\par}
{\footnotesize\color{muted}theorem; \code{KomorebiPaperIV/SignedCorrectionCompletion.lean}, line 87.\par}\smallskip
{\small\code{Komorebi.PaperIV.signedPositiveZero_eventual_derivative_sign}\par}
{\footnotesize\color{muted}theorem; \code{KomorebiPaperIV/SignedCorrectionCompletion.lean}, line 133.\par}\smallskip
\Needspace{7\baselineskip}
\subsection{Row 84: 11.9 Explicit cubic correction}
\label{row:84}
\textbf{Exact row scope.} The actual cubic remainder, spectral gap 3/2, lattice centers, and correction (-1)\textasciicircum{}m exp(-3t/2)/sqrt(3) are identified exactly. The actual positive zeros have error bounded by C exp(-3t). This fully encodes Corollary 11.9.\par\smallskip
\textbf{Manuscript labels.} \code{cor:N3_zero_correction}\par\smallskip
{\small\code{Komorebi.PaperIV.signedPolygonGap_three}\par}
{\footnotesize\color{muted}theorem; \code{KomorebiPaperIV/CubicZeroCorrection.lean}, line 10.\par}\smallskip
{\small\code{Komorebi.PaperIV.signedZeroCenter_three}\par}
{\footnotesize\color{muted}theorem; \code{KomorebiPaperIV/CubicZeroCorrection.lean}, line 15.\par}\smallskip
{\small\code{Komorebi.PaperIV.canonicalSignedVertex_three_one}\par}
{\footnotesize\color{muted}theorem; \code{KomorebiPaperIV/CubicZeroCorrection.lean}, line 21.\par}\smallskip
{\small\code{Komorebi.PaperIV.canonicalSignedVertex_three_two}\par}
{\footnotesize\color{muted}theorem; \code{KomorebiPaperIV/CubicZeroCorrection.lean}, line 29.\par}\smallskip
{\small\code{Komorebi.PaperIV.signedExponentialRemainder_three}\par}
{\footnotesize\color{muted}theorem; \code{KomorebiPaperIV/CubicZeroCorrection.lean}, line 41.\par}\smallskip
{\small\code{Komorebi.PaperIV.signedLeadingCorrection_three}\par}
{\footnotesize\color{muted}theorem; \code{KomorebiPaperIV/CubicZeroCorrection.lean}, line 52.\par}\smallskip
{\small\code{Komorebi.PaperIV.signedPositiveZero_three_correction}\par}
{\footnotesize\color{muted}theorem; \code{KomorebiPaperIV/CubicZeroCorrection.lean}, line 74.\par}\smallskip
\Needspace{7\baselineskip}
\subsection{Row 85: 14.8 Global derivative growth}
\label{row:85}
\textbf{Exact row scope.} Derives the exponential derivative lower bound at every sufficiently large zero on every symmetry ray. Its inverse positive integer powers dominate every real polynomial weight. No growth certificate is assumed.\par\smallskip
\textbf{Manuscript labels.} \code{thm:derivative_weighted_convergence}\par\smallskip
{\small\code{Komorebi.PaperIV.canonicalSignedDenominator_derivative_lower_bound}\par}
{\footnotesize\color{muted}theorem; \code{KomorebiPaperIV/DerivativeGrowth.lean}, line 16.\par}\smallskip
{\small\code{Komorebi.PaperIV.canonicalSignedDenominator_derivative_power_bound}\par}
{\footnotesize\color{muted}theorem; \code{KomorebiPaperIV/DerivativeGrowth.lean}, line 52.\par}\smallskip
\Needspace{7\baselineskip}
\subsection{Row 86: 14.8 Bounded logarithmic arguments}
\label{row:86}
\textbf{Exact row scope.} A proved Weierstrass criterion for complex coefficients and bounded logarithmic arguments. Includes absolute norm summability, local uniform convergence, entire sums and half-plane holomorphy. The final denominator specialization in row 87 proves the required summability hypotheses.\par\smallskip
\textbf{Manuscript labels.} \code{thm:derivative_weighted_convergence}\par\smallskip
{\small\code{Komorebi.PaperIV.branchedDirichletTerm}\par}
{\footnotesize\color{muted}definition; \code{KomorebiPaperIV/BranchedDirichlet.lean}, line 12.\par}\smallskip
{\small\code{Komorebi.PaperIV.branchedDirichletTerm_differentiable}\par}
{\footnotesize\color{muted}theorem; \code{KomorebiPaperIV/BranchedDirichlet.lean}, line 15.\par}\smallskip
{\small\code{Komorebi.PaperIV.branchedDirichletTerm_norm}\par}
{\footnotesize\color{muted}theorem; \code{KomorebiPaperIV/BranchedDirichlet.lean}, line 20.\par}\smallskip
{\small\code{Komorebi.PaperIV.branchedDirichletTerm_strip_bound}\par}
{\footnotesize\color{muted}theorem; \code{KomorebiPaperIV/BranchedDirichlet.lean}, line 30.\par}\smallskip
{\small\code{Komorebi.PaperIV.branchedDirichlet_summable_norm}\par}
{\footnotesize\color{muted}theorem; \code{KomorebiPaperIV/BranchedDirichlet.lean}, line 55.\par}\smallskip
{\small\code{Komorebi.PaperIV.branchedDirichlet_locally_uniform}\par}
{\footnotesize\color{muted}theorem; \code{KomorebiPaperIV/BranchedDirichlet.lean}, line 67.\par}\smallskip
{\small\code{Komorebi.PaperIV.branchedDirichlet_entire}\par}
{\footnotesize\color{muted}theorem; \code{KomorebiPaperIV/BranchedDirichlet.lean}, line 88.\par}\smallskip
{\small\code{Komorebi.PaperIV.branchedDirichlet_halfplane}\par}
{\footnotesize\color{muted}theorem; \code{KomorebiPaperIV/BranchedDirichlet.lean}, line 101.\par}\smallskip
\Needspace{7\baselineskip}
\subsection{Row 87: 14.8 Full derivative-weighted convergence}
\label{row:87}
\textbf{Exact row scope.} Defines the exact exp(-s Log z)/(D\_N prime(z)\textasciicircum{}q) series over all actual simple zeros, retaining exceptional simple zeros and the complex derivative weight. For every N >= 3 and q >= 1, proves absolute and locally uniform convergence everywhere and an entire sum. For N=2, proves absolute and locally uniform convergence and holomorphy on Re(s)>1. Logarithms must exponentiate to the zero and have arguments in a fixed bounded interval. This fully encodes Theorem 14.8.\par\smallskip
\textbf{Manuscript labels.} \code{thm:derivative_weighted_convergence}\par\smallskip
{\small\code{Komorebi.PaperIV.SimpleZeros}\par}
{\footnotesize\color{muted}definition; \code{KomorebiPaperIV/DerivativeWeightedZeta.lean}, line 15.\par}\smallskip
{\small\code{Komorebi.PaperIV.SimpleZeros.toZero}\par}
{\footnotesize\color{muted}definition; \code{KomorebiPaperIV/DerivativeWeightedZeta.lean}, line 17.\par}\smallskip
{\small\code{Komorebi.PaperIV.SimpleZeros.toZero_injective}\par}
{\footnotesize\color{muted}theorem; \code{KomorebiPaperIV/DerivativeWeightedZeta.lean}, line 20.\par}\smallskip
{\small\code{Komorebi.PaperIV.SimpleZeros.order_one}\par}
{\footnotesize\color{muted}theorem; \code{KomorebiPaperIV/DerivativeWeightedZeta.lean}, line 25.\par}\smallskip
{\small\code{Komorebi.PaperIV.canonicalSignedDenominator_simple_zero_power_summable}\par}
{\footnotesize\color{muted}theorem; \code{KomorebiPaperIV/DerivativeWeightedZeta.lean}, line 31.\par}\smallskip
{\small\code{Komorebi.PaperIV.simple_logarithm_re}\par}
{\footnotesize\color{muted}theorem; \code{KomorebiPaperIV/DerivativeWeightedZeta.lean}, line 43.\par}\smallskip
{\small\code{Komorebi.PaperIV.derivativeZeroWeight}\par}
{\footnotesize\color{muted}definition; \code{KomorebiPaperIV/DerivativeWeightedZeta.lean}, line 50.\par}\smallskip
{\small\code{Komorebi.PaperIV.derivativeZeroWeight_norm}\par}
{\footnotesize\color{muted}theorem; \code{KomorebiPaperIV/DerivativeWeightedZeta.lean}, line 53.\par}\smallskip
{\small\code{Komorebi.PaperIV.derivativeZeroWeight_power_summable}\par}
{\footnotesize\color{muted}theorem; \code{KomorebiPaperIV/DerivativeWeightedZeta.lean}, line 59.\par}\smallskip
{\small\code{Komorebi.PaperIV.canonicalSignedDenominator_two_derivative_norm}\par}
{\footnotesize\color{muted}theorem; \code{KomorebiPaperIV/DerivativeWeightedZeta.lean}, line 79.\par}\smallskip
{\small\code{Komorebi.PaperIV.derivativeZeroWeight_two_power_summable}\par}
{\footnotesize\color{muted}theorem; \code{KomorebiPaperIV/DerivativeWeightedZeta.lean}, line 92.\par}\smallskip
{\small\code{Komorebi.PaperIV.derivativeWeightedZeroZetaTerm}\par}
{\footnotesize\color{muted}definition; \code{KomorebiPaperIV/DerivativeWeightedZeta.lean}, line 100.\par}\smallskip
{\small\code{Komorebi.PaperIV.derivativeWeightedZeroZeta}\par}
{\footnotesize\color{muted}definition; \code{KomorebiPaperIV/DerivativeWeightedZeta.lean}, line 105.\par}\smallskip
{\small\code{Komorebi.PaperIV.derivativeWeightedZeroZetaTerm_eq}\par}
{\footnotesize\color{muted}theorem; \code{KomorebiPaperIV/DerivativeWeightedZeta.lean}, line 109.\par}\smallskip
{\small\code{Komorebi.PaperIV.derivativeWeightedZeroZeta_summable_norm}\par}
{\footnotesize\color{muted}theorem; \code{KomorebiPaperIV/DerivativeWeightedZeta.lean}, line 117.\par}\smallskip
{\small\code{Komorebi.PaperIV.derivativeWeightedZeroZeta_two_summable_norm}\par}
{\footnotesize\color{muted}theorem; \code{KomorebiPaperIV/DerivativeWeightedZeta.lean}, line 128.\par}\smallskip
{\small\code{Komorebi.PaperIV.derivativeWeightedZeroZeta_entire}\par}
{\footnotesize\color{muted}theorem; \code{KomorebiPaperIV/DerivativeWeightedZeta.lean}, line 141.\par}\smallskip
{\small\code{Komorebi.PaperIV.derivativeWeightedZeroZeta_two_holomorphic}\par}
{\footnotesize\color{muted}theorem; \code{KomorebiPaperIV/DerivativeWeightedZeta.lean}, line 156.\par}\smallskip
\Needspace{7\baselineskip}
\subsection{Row 88: 14.3 Entire exponential comparison}
\label{row:88}
\textbf{Exact row scope.} A locally uniform summable majorant on every bounded parameter set proves that exponentially displaced positive sequences have an entire difference of zeta sums. The needed displacement and summability are derived in row 90.\par\smallskip
\textbf{Manuscript labels.} \code{thm:zeta_continuation}\par\smallskip
{\small\code{Komorebi.PaperIV.positiveZetaTerm}\par}
{\footnotesize\color{muted}definition; \code{KomorebiPaperIV/ExponentialZetaComparison.lean}, line 15.\par}\smallskip
{\small\code{Komorebi.PaperIV.exponentialZetaDifference}\par}
{\footnotesize\color{muted}definition; \code{KomorebiPaperIV/ExponentialZetaComparison.lean}, line 17.\par}\smallskip
{\small\code{Komorebi.PaperIV.positiveZetaTerm_entire}\par}
{\footnotesize\color{muted}theorem; \code{KomorebiPaperIV/ExponentialZetaComparison.lean}, line 20.\par}\smallskip
{\small\code{Komorebi.PaperIV.exponentialZetaDifference_entire}\par}
{\footnotesize\color{muted}theorem; \code{KomorebiPaperIV/ExponentialZetaComparison.lean}, line 24.\par}\smallskip
{\small\code{Komorebi.PaperIV.real_log_lipschitz_one}\par}
{\footnotesize\color{muted}theorem; \code{KomorebiPaperIV/ExponentialZetaComparison.lean}, line 28.\par}\smallskip
{\small\code{Komorebi.PaperIV.positiveZetaTerm_norm_le_rpow}\par}
{\footnotesize\color{muted}theorem; \code{KomorebiPaperIV/ExponentialZetaComparison.lean}, line 38.\par}\smallskip
{\small\code{Komorebi.PaperIV.exponentialZetaDifference_factor}\par}
{\footnotesize\color{muted}theorem; \code{KomorebiPaperIV/ExponentialZetaComparison.lean}, line 49.\par}\smallskip
{\small\code{Komorebi.PaperIV.summable_polynomial_exponential}\par}
{\footnotesize\color{muted}theorem; \code{KomorebiPaperIV/ExponentialZetaComparison.lean}, line 58.\par}\smallskip
{\small\code{Komorebi.PaperIV.exponentialZetaDifference_majorant}\par}
{\footnotesize\color{muted}theorem; \code{KomorebiPaperIV/ExponentialZetaComparison.lean}, line 73.\par}\smallskip
{\small\code{Komorebi.PaperIV.exponentialZetaDifference_locally_uniform}\par}
{\footnotesize\color{muted}theorem; \code{KomorebiPaperIV/ExponentialZetaComparison.lean}, line 108.\par}\smallskip
{\small\code{Komorebi.PaperIV.exponentialZetaDifference_sum_entire}\par}
{\footnotesize\color{muted}theorem; \code{KomorebiPaperIV/ExponentialZetaComparison.lean}, line 127.\par}\smallskip
\Needspace{7\baselineskip}
\subsection{Row 89: 14.3 Hurwitz model and pole}
\label{row:89}
\textbf{Exact row scope.} The scaled Hurwitz model is meromorphic on the whole plane, holomorphic off its sole possible pole, with the exact residue. A nonzero residue gives meromorphic order -1. Uses the pinned proved Mathlib Hurwitz theory.\par\smallskip
\textbf{Manuscript labels.} \code{thm:zeta_continuation}\par\smallskip
{\small\code{Komorebi.PaperIV.hurwitzRegularPart}\par}
{\footnotesize\color{muted}definition; \code{KomorebiPaperIV/HurwitzLattice.lean}, line 15.\par}\smallskip
{\small\code{Komorebi.PaperIV.hurwitzRegularPart_entire}\par}
{\footnotesize\color{muted}theorem; \code{KomorebiPaperIV/HurwitzLattice.lean}, line 18.\par}\smallskip
{\small\code{Komorebi.PaperIV.hurwitzZeta_meromorphicAt}\par}
{\footnotesize\color{muted}theorem; \code{KomorebiPaperIV/HurwitzLattice.lean}, line 30.\par}\smallskip
{\small\code{Komorebi.PaperIV.hurwitzLatticeModel}\par}
{\footnotesize\color{muted}definition; \code{KomorebiPaperIV/HurwitzLattice.lean}, line 40.\par}\smallskip
{\small\code{Komorebi.PaperIV.hurwitzLatticeModel_meromorphicAt}\par}
{\footnotesize\color{muted}theorem; \code{KomorebiPaperIV/HurwitzLattice.lean}, line 44.\par}\smallskip
{\small\code{Komorebi.PaperIV.hurwitzLatticeModel_differentiableAt}\par}
{\footnotesize\color{muted}theorem; \code{KomorebiPaperIV/HurwitzLattice.lean}, line 51.\par}\smallskip
{\small\code{Komorebi.PaperIV.hurwitzLatticeModel_residue}\par}
{\footnotesize\color{muted}theorem; \code{KomorebiPaperIV/HurwitzLattice.lean}, line 63.\par}\smallskip
{\small\code{Komorebi.PaperIV.meromorphicOrderAt_eq_neg_one_of_residue}\par}
{\footnotesize\color{muted}theorem; \code{KomorebiPaperIV/HurwitzLattice.lean}, line 95.\par}\smallskip
\Needspace{7\baselineskip}
\subsection{Row 90: 14.3 Actual positive-zero comparison}
\label{row:90}
\textbf{Exact row scope.} The proved enumeration carries its actual exponential displacement and offset. Its entire correction and finite missing lattice prefix are constructed, and the exact root-zeta identity with the unshifted Hurwitz model is proved.\par\smallskip
\textbf{Manuscript labels.} \code{thm:zeta_continuation}\par\smallskip
{\small\code{Komorebi.PaperIV.positiveZetaTerm_mul}\par}
{\footnotesize\color{muted}theorem; \code{KomorebiPaperIV/ZeroLatticeComparison.lean}, line 13.\par}\smallskip
{\small\code{Komorebi.PaperIV.positiveZetaTerm_eq_cpow}\par}
{\footnotesize\color{muted}theorem; \code{KomorebiPaperIV/ZeroLatticeComparison.lean}, line 18.\par}\smallskip
{\small\code{Komorebi.PaperIV.positiveZetaTerm_norm}\par}
{\footnotesize\color{muted}theorem; \code{KomorebiPaperIV/ZeroLatticeComparison.lean}, line 26.\par}\smallskip
{\small\code{Komorebi.PaperIV.positiveZetaTerm_lattice_hasSum}\par}
{\footnotesize\color{muted}theorem; \code{KomorebiPaperIV/ZeroLatticeComparison.lean}, line 31.\par}\smallskip
{\small\code{Komorebi.PaperIV.signedZeroCenter_positive}\par}
{\footnotesize\color{muted}theorem; \code{KomorebiPaperIV/ZeroLatticeComparison.lean}, line 43.\par}\smallskip
{\small\code{Komorebi.PaperIV.signedZeroCenter_reciprocal_square_summable}\par}
{\footnotesize\color{muted}theorem; \code{KomorebiPaperIV/ZeroLatticeComparison.lean}, line 48.\par}\smallskip
{\small\code{Komorebi.PaperIV.SignedZeroEnumeration.root_tendsto_atTop}\par}
{\footnotesize\color{muted}theorem; \code{KomorebiPaperIV/ZeroLatticeComparison.lean}, line 55.\par}\smallskip
{\small\code{Komorebi.PaperIV.SignedZeroEnumeration.latticeCorrection}\par}
{\footnotesize\color{muted}definition; \code{KomorebiPaperIV/ZeroLatticeComparison.lean}, line 61.\par}\smallskip
{\small\code{Komorebi.PaperIV.SignedZeroEnumeration.latticeCorrection_locally_uniform}\par}
{\footnotesize\color{muted}theorem; \code{KomorebiPaperIV/ZeroLatticeComparison.lean}, line 64.\par}\smallskip
{\small\code{Komorebi.PaperIV.SignedZeroEnumeration.latticeCorrection_entire}\par}
{\footnotesize\color{muted}theorem; \code{KomorebiPaperIV/ZeroLatticeComparison.lean}, line 77.\par}\smallskip
{\small\code{Komorebi.PaperIV.SignedZeroEnumeration.positiveZetaTerm_summable}\par}
{\footnotesize\color{muted}theorem; \code{KomorebiPaperIV/ZeroLatticeComparison.lean}, line 83.\par}\smallskip
{\small\code{Komorebi.PaperIV.SignedZeroEnumeration.latticePrefix}\par}
{\footnotesize\color{muted}definition; \code{KomorebiPaperIV/ZeroLatticeComparison.lean}, line 90.\par}\smallskip
{\small\code{Komorebi.PaperIV.SignedZeroEnumeration.latticePrefix_entire}\par}
{\footnotesize\color{muted}theorem; \code{KomorebiPaperIV/ZeroLatticeComparison.lean}, line 93.\par}\smallskip
{\small\code{Komorebi.PaperIV.SignedZeroEnumeration.root_zeta_eq_model}\par}
{\footnotesize\color{muted}theorem; \code{KomorebiPaperIV/ZeroLatticeComparison.lean}, line 99.\par}\smallskip
\Needspace{7\baselineskip}
\subsection{Row 91: 14.3 All-zero decomposition}
\label{row:91}
\textbf{Exact row scope.} Both original and reduced zeta sums split into the actual simple tail and the finite exceptional zeros with their analytic multiplicities. The factor N and the scaled parameter Ns are proved.\par\smallskip
\textbf{Manuscript labels.} \code{thm:zeta_continuation}\par\smallskip
{\small\code{Komorebi.PaperIV.tsum_injection_add_complement}\par}
{\footnotesize\color{muted}theorem; \code{KomorebiPaperIV/ZeroZetaDecomposition.lean}, line 12.\par}\smallskip
{\small\code{Komorebi.PaperIV.finite_subtype_tsum_entire}\par}
{\footnotesize\color{muted}theorem; \code{KomorebiPaperIV/ZeroZetaDecomposition.lean}, line 18.\par}\smallskip
{\small\code{Komorebi.PaperIV.SignedZeroEnumeration.reducedExceptionZeta}\par}
{\footnotesize\color{muted}definition; \code{KomorebiPaperIV/ZeroZetaDecomposition.lean}, line 26.\par}\smallskip
{\small\code{Komorebi.PaperIV.SignedZeroEnumeration.absoluteExceptionZeta}\par}
{\footnotesize\color{muted}definition; \code{KomorebiPaperIV/ZeroZetaDecomposition.lean}, line 30.\par}\smallskip
{\small\code{Komorebi.PaperIV.absoluteZeroZetaTerm_entire}\par}
{\footnotesize\color{muted}theorem; \code{KomorebiPaperIV/ZeroZetaDecomposition.lean}, line 34.\par}\smallskip
{\small\code{Komorebi.PaperIV.SignedZeroEnumeration.reducedExceptionZeta_entire}\par}
{\footnotesize\color{muted}theorem; \code{KomorebiPaperIV/ZeroZetaDecomposition.lean}, line 39.\par}\smallskip
{\small\code{Komorebi.PaperIV.SignedZeroEnumeration.absoluteExceptionZeta_entire}\par}
{\footnotesize\color{muted}theorem; \code{KomorebiPaperIV/ZeroZetaDecomposition.lean}, line 44.\par}\smallskip
{\small\code{Komorebi.PaperIV.SignedZeroEnumeration.reducedZetaTerm}\par}
{\footnotesize\color{muted}theorem; \code{KomorebiPaperIV/ZeroZetaDecomposition.lean}, line 49.\par}\smallskip
{\small\code{Komorebi.PaperIV.SignedZeroEnumeration.orbitZetaTerm}\par}
{\footnotesize\color{muted}theorem; \code{KomorebiPaperIV/ZeroZetaDecomposition.lean}, line 60.\par}\smallskip
{\small\code{Komorebi.PaperIV.SignedZeroEnumeration.reducedZeta_decomposition}\par}
{\footnotesize\color{muted}theorem; \code{KomorebiPaperIV/ZeroZetaDecomposition.lean}, line 67.\par}\smallskip
{\small\code{Komorebi.PaperIV.SignedZeroEnumeration.absoluteZeta_decomposition}\par}
{\footnotesize\color{muted}theorem; \code{KomorebiPaperIV/ZeroZetaDecomposition.lean}, line 77.\par}\smallskip
\Needspace{7\baselineskip}
\subsection{Row 92: 14.3 Full unweighted continuation}
\label{row:92}
\textbf{Exact row scope.} Constructs both entire corrections, proves equality on the original convergence domains, global meromorphic continuation, holomorphy off the sole pole, meromorphic order -1 at the pole, and both exact manuscript residues. Fully encodes Theorem 14.3.\par\smallskip
\textbf{Manuscript labels.} \code{thm:zeta_continuation}\par\smallskip
{\small\code{Komorebi.PaperIV.reducedZetaCorrection}\par}
{\footnotesize\color{muted}definition; \code{KomorebiPaperIV/ZeroZetaContinuation.lean}, line 12.\par}\smallskip
{\small\code{Komorebi.PaperIV.absoluteZetaCorrection}\par}
{\footnotesize\color{muted}definition; \code{KomorebiPaperIV/ZeroZetaContinuation.lean}, line 17.\par}\smallskip
{\small\code{Komorebi.PaperIV.reducedZetaCorrection_entire}\par}
{\footnotesize\color{muted}theorem; \code{KomorebiPaperIV/ZeroZetaContinuation.lean}, line 22.\par}\smallskip
{\small\code{Komorebi.PaperIV.absoluteZetaCorrection_entire}\par}
{\footnotesize\color{muted}theorem; \code{KomorebiPaperIV/ZeroZetaContinuation.lean}, line 29.\par}\smallskip
{\small\code{Komorebi.PaperIV.entireReducedDenominator_zeta_eq_hurwitz}\par}
{\footnotesize\color{muted}theorem; \code{KomorebiPaperIV/ZeroZetaContinuation.lean}, line 38.\par}\smallskip
{\small\code{Komorebi.PaperIV.canonicalSignedDenominator_zeta_eq_hurwitz}\par}
{\footnotesize\color{muted}theorem; \code{KomorebiPaperIV/ZeroZetaContinuation.lean}, line 51.\par}\smallskip
{\small\code{Komorebi.PaperIV.reducedZetaContinuation}\par}
{\footnotesize\color{muted}definition; \code{KomorebiPaperIV/ZeroZetaContinuation.lean}, line 61.\par}\smallskip
{\small\code{Komorebi.PaperIV.absoluteZetaContinuation}\par}
{\footnotesize\color{muted}definition; \code{KomorebiPaperIV/ZeroZetaContinuation.lean}, line 64.\par}\smallskip
{\small\code{Komorebi.PaperIV.reducedZetaContinuation_meromorphic}\par}
{\footnotesize\color{muted}theorem; \code{KomorebiPaperIV/ZeroZetaContinuation.lean}, line 67.\par}\smallskip
{\small\code{Komorebi.PaperIV.absoluteZetaContinuation_meromorphic}\par}
{\footnotesize\color{muted}theorem; \code{KomorebiPaperIV/ZeroZetaContinuation.lean}, line 72.\par}\smallskip
{\small\code{Komorebi.PaperIV.reducedZetaContinuation_holomorphic_off_pole}\par}
{\footnotesize\color{muted}theorem; \code{KomorebiPaperIV/ZeroZetaContinuation.lean}, line 77.\par}\smallskip
{\small\code{Komorebi.PaperIV.absoluteZetaContinuation_holomorphic_off_pole}\par}
{\footnotesize\color{muted}theorem; \code{KomorebiPaperIV/ZeroZetaContinuation.lean}, line 81.\par}\smallskip
{\small\code{Komorebi.PaperIV.entire_add_preserves_residue}\par}
{\footnotesize\color{muted}theorem; \code{KomorebiPaperIV/ZeroZetaContinuation.lean}, line 87.\par}\smallskip
{\small\code{Komorebi.PaperIV.reducedZetaContinuation_residue}\par}
{\footnotesize\color{muted}theorem; \code{KomorebiPaperIV/ZeroZetaContinuation.lean}, line 96.\par}\smallskip
{\small\code{Komorebi.PaperIV.absoluteZetaContinuation_residue}\par}
{\footnotesize\color{muted}theorem; \code{KomorebiPaperIV/ZeroZetaContinuation.lean}, line 106.\par}\smallskip
{\small\code{Komorebi.PaperIV.reducedZetaContinuation_simple_pole}\par}
{\footnotesize\color{muted}theorem; \code{KomorebiPaperIV/ZeroZetaContinuation.lean}, line 120.\par}\smallskip
{\small\code{Komorebi.PaperIV.absoluteZetaContinuation_simple_pole}\par}
{\footnotesize\color{muted}theorem; \code{KomorebiPaperIV/ZeroZetaContinuation.lean}, line 127.\par}\smallskip
\Needspace{7\baselineskip}
\subsection{Row 93: 14.6 Full branch-dependent continuation}
\label{row:93}
\textbf{Exact row scope.} Defines the actual multiplicity-weighted branched zero series. For logarithms real on the proved positive tail, it agrees with the absolute zeta plus an entire finite correction. Proves global continuation, the same single simple pole/residue, and an entire difference between choices. Intermediate finite-correction claims hold even for arbitrary exceptional complex values. Fully encodes Corollary 14.6.\par\smallskip
\textbf{Manuscript labels.} \code{cor:branch_grouped_continuation}\par\smallskip
{\small\code{Komorebi.PaperIV.branchedZeroZetaTerm}\par}
{\footnotesize\color{muted}definition; \code{KomorebiPaperIV/BranchZetaContinuation.lean}, line 14.\par}\smallskip
{\small\code{Komorebi.PaperIV.branchedZeroZeta}\par}
{\footnotesize\color{muted}definition; \code{KomorebiPaperIV/BranchZetaContinuation.lean}, line 18.\par}\smallskip
{\small\code{Komorebi.PaperIV.branchedZeroZetaTerm_entire}\par}
{\footnotesize\color{muted}theorem; \code{KomorebiPaperIV/BranchZetaContinuation.lean}, line 21.\par}\smallskip
{\small\code{Komorebi.PaperIV.RealLogOnReducedTail}\par}
{\footnotesize\color{muted}definition; \code{KomorebiPaperIV/BranchZetaContinuation.lean}, line 29.\par}\smallskip
{\small\code{Komorebi.PaperIV.reducedBranchCorrection}\par}
{\footnotesize\color{muted}definition; \code{KomorebiPaperIV/BranchZetaContinuation.lean}, line 34.\par}\smallskip
{\small\code{Komorebi.PaperIV.reducedBranchCorrection_entire}\par}
{\footnotesize\color{muted}theorem; \code{KomorebiPaperIV/BranchZetaContinuation.lean}, line 40.\par}\smallskip
{\small\code{Komorebi.PaperIV.branchedZeroZetaTerm_eq_absolute_on_tail}\par}
{\footnotesize\color{muted}theorem; \code{KomorebiPaperIV/BranchZetaContinuation.lean}, line 49.\par}\smallskip
{\small\code{Komorebi.PaperIV.branchedReducedZeta_summable}\par}
{\footnotesize\color{muted}theorem; \code{KomorebiPaperIV/BranchZetaContinuation.lean}, line 59.\par}\smallskip
{\small\code{Komorebi.PaperIV.branchedReducedZeta_eq_absolute_add_entire}\par}
{\footnotesize\color{muted}theorem; \code{KomorebiPaperIV/BranchZetaContinuation.lean}, line 79.\par}\smallskip
{\small\code{Komorebi.PaperIV.branchedReducedZetaContinuation}\par}
{\footnotesize\color{muted}definition; \code{KomorebiPaperIV/BranchZetaContinuation.lean}, line 103.\par}\smallskip
{\small\code{Komorebi.PaperIV.branchedReducedZetaContinuation_eq}\par}
{\footnotesize\color{muted}theorem; \code{KomorebiPaperIV/BranchZetaContinuation.lean}, line 107.\par}\smallskip
{\small\code{Komorebi.PaperIV.branchedReducedZetaContinuation_meromorphic}\par}
{\footnotesize\color{muted}theorem; \code{KomorebiPaperIV/BranchZetaContinuation.lean}, line 114.\par}\smallskip
{\small\code{Komorebi.PaperIV.branchedReducedZetaContinuation_holomorphic_off_pole}\par}
{\footnotesize\color{muted}theorem; \code{KomorebiPaperIV/BranchZetaContinuation.lean}, line 120.\par}\smallskip
{\small\code{Komorebi.PaperIV.branchedReducedZetaContinuation_residue}\par}
{\footnotesize\color{muted}theorem; \code{KomorebiPaperIV/BranchZetaContinuation.lean}, line 125.\par}\smallskip
{\small\code{Komorebi.PaperIV.branchedReducedZetaContinuation_simple_pole}\par}
{\footnotesize\color{muted}theorem; \code{KomorebiPaperIV/BranchZetaContinuation.lean}, line 131.\par}\smallskip
{\small\code{Komorebi.PaperIV.branchedReducedZetaContinuation_difference_entire}\par}
{\footnotesize\color{muted}theorem; \code{KomorebiPaperIV/BranchZetaContinuation.lean}, line 141.\par}\smallskip
\Needspace{7\baselineskip}
\subsection{Row 94: 14.4 Exact quartic factorization and zeros}
\label{row:94}
\textbf{Exact row scope.} Proves D4(z)=cos(z/sqrt(2))*cosh(z/sqrt(2)), its exhaustive real/imaginary zero strings, and simplicity of every zero.\par\smallskip
\textbf{Manuscript labels.} \code{cor:exact_zeta_models}\par\smallskip
{\small\code{Komorebi.PaperIV.quartic_root_square}\par}
{\footnotesize\color{muted}theorem; \code{KomorebiPaperIV/QuarticExactZeros.lean}, line 12.\par}\smallskip
{\small\code{Komorebi.PaperIV.canonicalSignedDenominator_four_eq_cos_cosh}\par}
{\footnotesize\color{muted}theorem; \code{KomorebiPaperIV/QuarticExactZeros.lean}, line 19.\par}\smallskip
{\small\code{Komorebi.PaperIV.complex_cos_zero_iff_nat}\par}
{\footnotesize\color{muted}theorem; \code{KomorebiPaperIV/QuarticExactZeros.lean}, line 43.\par}\smallskip
{\small\code{Komorebi.PaperIV.quarticPositiveZero}\par}
{\footnotesize\color{muted}definition; \code{KomorebiPaperIV/QuarticExactZeros.lean}, line 49.\par}\smallskip
{\small\code{Komorebi.PaperIV.quarticPositiveZero_pos}\par}
{\footnotesize\color{muted}theorem; \code{KomorebiPaperIV/QuarticExactZeros.lean}, line 51.\par}\smallskip
{\small\code{Komorebi.PaperIV.quarticPositiveZero_div_sqrt}\par}
{\footnotesize\color{muted}theorem; \code{KomorebiPaperIV/QuarticExactZeros.lean}, line 55.\par}\smallskip
{\small\code{Komorebi.PaperIV.canonicalSignedDenominator_four_zero_iff}\par}
{\footnotesize\color{muted}theorem; \code{KomorebiPaperIV/QuarticExactZeros.lean}, line 63.\par}\smallskip
{\small\code{Komorebi.PaperIV.quarticPositiveZero_simple}\par}
{\footnotesize\color{muted}theorem; \code{KomorebiPaperIV/QuarticExactZeros.lean}, line 112.\par}\smallskip
{\small\code{Komorebi.PaperIV.canonicalSignedDenominator_four_zero_order_one}\par}
{\footnotesize\color{muted}theorem; \code{KomorebiPaperIV/QuarticExactZeros.lean}, line 141.\par}\smallskip
\Needspace{7\baselineskip}
\subsection{Row 95: 14.4 Full exact order-two/four models}
\label{row:95}
\textbf{Exact row scope.} Proves both exact zeta formulas in orders two and four with no omitted exceptional zeros. The identity principle proves that the entire corrections constructed in Theorem 14.3 vanish identically. Fully encodes Corollary 14.4.\par\smallskip
\textbf{Manuscript labels.} \code{cor:exact_zeta_models}\par\smallskip
{\small\code{Komorebi.PaperIV.exactLatticeEnumeration}\par}
{\footnotesize\color{muted}definition; \code{KomorebiPaperIV/ExactZetaModels.lean}, line 13.\par}\smallskip
{\small\code{Komorebi.PaperIV.SignedZeroEnumeration.reducedZero_surjective_of_orbitZero}\par}
{\footnotesize\color{muted}theorem; \code{KomorebiPaperIV/ExactZetaModels.lean}, line 41.\par}\smallskip
{\small\code{Komorebi.PaperIV.SignedZeroEnumeration.exceptions_zero_of_surjective}\par}
{\footnotesize\color{muted}theorem; \code{KomorebiPaperIV/ExactZetaModels.lean}, line 57.\par}\smallskip
{\small\code{Komorebi.PaperIV.zero_zeta_exact_of_lattice}\par}
{\footnotesize\color{muted}theorem; \code{KomorebiPaperIV/ExactZetaModels.lean}, line 66.\par}\smallskip
{\small\code{Komorebi.PaperIV.signedZeroCenter_four}\par}
{\footnotesize\color{muted}theorem; \code{KomorebiPaperIV/ExactZetaModels.lean}, line 100.\par}\smallskip
{\small\code{Komorebi.PaperIV.orderTwo_full_lattice}\par}
{\footnotesize\color{muted}theorem; \code{KomorebiPaperIV/ExactZetaModels.lean}, line 108.\par}\smallskip
{\small\code{Komorebi.PaperIV.orderFour_full_lattice}\par}
{\footnotesize\color{muted}theorem; \code{KomorebiPaperIV/ExactZetaModels.lean}, line 129.\par}\smallskip
{\small\code{Komorebi.PaperIV.orderTwo_exact_zero_zeta_models}\par}
{\footnotesize\color{muted}theorem; \code{KomorebiPaperIV/ExactZetaModels.lean}, line 157.\par}\smallskip
{\small\code{Komorebi.PaperIV.orderFour_exact_zero_zeta_models}\par}
{\footnotesize\color{muted}theorem; \code{KomorebiPaperIV/ExactZetaModels.lean}, line 165.\par}\smallskip
{\small\code{Komorebi.PaperIV.entire_zero_of_halfplane}\par}
{\footnotesize\color{muted}theorem; \code{KomorebiPaperIV/ExactZetaModels.lean}, line 177.\par}\smallskip
{\small\code{Komorebi.PaperIV.zero_zeta_corrections_of_lattice}\par}
{\footnotesize\color{muted}theorem; \code{KomorebiPaperIV/ExactZetaModels.lean}, line 188.\par}\smallskip
{\small\code{Komorebi.PaperIV.orderTwo_zeta_corrections_zero}\par}
{\footnotesize\color{muted}theorem; \code{KomorebiPaperIV/ExactZetaModels.lean}, line 207.\par}\smallskip
{\small\code{Komorebi.PaperIV.orderFour_zeta_corrections_zero}\par}
{\footnotesize\color{muted}theorem; \code{KomorebiPaperIV/ExactZetaModels.lean}, line 213.\par}\smallskip
\Needspace{7\baselineskip}
\subsection{Row 96: 12.1 Complete multiplicity-indexed product}
\label{row:96}
\textbf{Exact row scope.} Constructs the product over every actual reduced zero repeated according to analytic multiplicity, including all exceptional zeros and N=2. Proves countability, reciprocal summability, compact-uniform convergence, entire realization, normalization and exact zero set. Equality with Delta is supplied by row 102.\par\smallskip
\textbf{Manuscript labels.} \code{thm:genus_zero_product}\par\smallskip
{\small\code{Komorebi.PaperIV.ZeroMultiplicityIndex}\par}
{\footnotesize\color{muted}definition; \code{KomorebiPaperIV/FullCanonicalProduct.lean}, line 14.\par}\smallskip
{\small\code{Komorebi.PaperIV.indexedZero}\par}
{\footnotesize\color{muted}definition; \code{KomorebiPaperIV/FullCanonicalProduct.lean}, line 17.\par}\smallskip
{\small\code{Komorebi.PaperIV.fullCanonicalProduct}\par}
{\footnotesize\color{muted}definition; \code{KomorebiPaperIV/FullCanonicalProduct.lean}, line 19.\par}\smallskip
{\small\code{Komorebi.PaperIV.zeroMultiplicity_reciprocal_summable}\par}
{\footnotesize\color{muted}theorem; \code{KomorebiPaperIV/FullCanonicalProduct.lean}, line 22.\par}\smallskip
{\small\code{Komorebi.PaperIV.fullCanonicalProduct_factors_summable}\par}
{\footnotesize\color{muted}theorem; \code{KomorebiPaperIV/FullCanonicalProduct.lean}, line 31.\par}\smallskip
{\small\code{Komorebi.PaperIV.fullCanonicalProduct_multipliable}\par}
{\footnotesize\color{muted}theorem; \code{KomorebiPaperIV/FullCanonicalProduct.lean}, line 35.\par}\smallskip
{\small\code{Komorebi.PaperIV.fullCanonicalProduct_uniform_on_compact}\par}
{\footnotesize\color{muted}theorem; \code{KomorebiPaperIV/FullCanonicalProduct.lean}, line 39.\par}\smallskip
{\small\code{Komorebi.PaperIV.fullCanonicalProduct_locally_uniform}\par}
{\footnotesize\color{muted}theorem; \code{KomorebiPaperIV/FullCanonicalProduct.lean}, line 51.\par}\smallskip
{\small\code{Komorebi.PaperIV.fullCanonicalProduct_entire}\par}
{\footnotesize\color{muted}theorem; \code{KomorebiPaperIV/FullCanonicalProduct.lean}, line 58.\par}\smallskip
{\small\code{Komorebi.PaperIV.fullCanonicalProduct_zero}\par}
{\footnotesize\color{muted}theorem; \code{KomorebiPaperIV/FullCanonicalProduct.lean}, line 68.\par}\smallskip
{\small\code{Komorebi.PaperIV.fullCanonicalProduct_ne_zero}\par}
{\footnotesize\color{muted}theorem; \code{KomorebiPaperIV/FullCanonicalProduct.lean}, line 71.\par}\smallskip
{\small\code{Komorebi.PaperIV.fullCanonicalProduct_zero_iff_index}\par}
{\footnotesize\color{muted}theorem; \code{KomorebiPaperIV/FullCanonicalProduct.lean}, line 89.\par}\smallskip
{\small\code{Komorebi.PaperIV.indexedZero_ne_zero}\par}
{\footnotesize\color{muted}theorem; \code{KomorebiPaperIV/FullCanonicalProduct.lean}, line 108.\par}\smallskip
{\small\code{Komorebi.PaperIV.entire_zero_multiplicity_pos}\par}
{\footnotesize\color{muted}theorem; \code{KomorebiPaperIV/FullCanonicalProduct.lean}, line 116.\par}\smallskip
{\small\code{Komorebi.PaperIV.fullCanonicalProduct_zero_iff}\par}
{\footnotesize\color{muted}theorem; \code{KomorebiPaperIV/FullCanonicalProduct.lean}, line 122.\par}\smallskip
{\small\code{Komorebi.PaperIV.zeroMultiplicityIndex_countable}\par}
{\footnotesize\color{muted}theorem; \code{KomorebiPaperIV/FullCanonicalProduct.lean}, line 132.\par}\smallskip
{\small\code{Komorebi.PaperIV.entireReducedDenominator_full_product_properties}\par}
{\footnotesize\color{muted}theorem; \code{KomorebiPaperIV/FullCanonicalProduct.lean}, line 143.\par}\smallskip
\Needspace{7\baselineskip}
\subsection{Row 97: 12.1 Exact product multiplicities}
\label{row:97}
\textbf{Exact row scope.} Factors out every copy of a selected zero; the remaining product is entire and nonzero there. Proves that the complete product and the reduced denominator have exactly the same analytic order at every point.\par\smallskip
\textbf{Manuscript labels.} \code{thm:genus_zero_product}\par\smallskip
{\small\code{Komorebi.PaperIV.reciprocalProduct}\par}
{\footnotesize\color{muted}definition; \code{KomorebiPaperIV/CanonicalProductMultiplicity.lean}, line 13.\par}\smallskip
{\small\code{Komorebi.PaperIV.reciprocalProduct_locally_uniform}\par}
{\footnotesize\color{muted}theorem; \code{KomorebiPaperIV/CanonicalProductMultiplicity.lean}, line 15.\par}\smallskip
{\small\code{Komorebi.PaperIV.reciprocalProduct_entire}\par}
{\footnotesize\color{muted}theorem; \code{KomorebiPaperIV/CanonicalProductMultiplicity.lean}, line 27.\par}\smallskip
{\small\code{Komorebi.PaperIV.reciprocalProduct_ne_zero}\par}
{\footnotesize\color{muted}theorem; \code{KomorebiPaperIV/CanonicalProductMultiplicity.lean}, line 37.\par}\smallskip
{\small\code{Komorebi.PaperIV.zeroIndexFiberEquiv}\par}
{\footnotesize\color{muted}definition; \code{KomorebiPaperIV/CanonicalProductMultiplicity.lean}, line 47.\par}\smallskip
{\small\code{Komorebi.PaperIV.otherZeroProduct}\par}
{\footnotesize\color{muted}definition; \code{KomorebiPaperIV/CanonicalProductMultiplicity.lean}, line 60.\par}\smallskip
{\small\code{Komorebi.PaperIV.otherZeroProduct_entire}\par}
{\footnotesize\color{muted}theorem; \code{KomorebiPaperIV/CanonicalProductMultiplicity.lean}, line 63.\par}\smallskip
{\small\code{Komorebi.PaperIV.otherZeroProduct_ne_zero_at_zero}\par}
{\footnotesize\color{muted}theorem; \code{KomorebiPaperIV/CanonicalProductMultiplicity.lean}, line 67.\par}\smallskip
{\small\code{Komorebi.PaperIV.fullCanonicalProduct_factor_at_zero}\par}
{\footnotesize\color{muted}theorem; \code{KomorebiPaperIV/CanonicalProductMultiplicity.lean}, line 76.\par}\smallskip
{\small\code{Komorebi.PaperIV.linear_zero_factor_order}\par}
{\footnotesize\color{muted}theorem; \code{KomorebiPaperIV/CanonicalProductMultiplicity.lean}, line 95.\par}\smallskip
{\small\code{Komorebi.PaperIV.fullCanonicalProduct_order_at_zero}\par}
{\footnotesize\color{muted}theorem; \code{KomorebiPaperIV/CanonicalProductMultiplicity.lean}, line 103.\par}\smallskip
{\small\code{Komorebi.PaperIV.fullCanonicalProduct_same_orders}\par}
{\footnotesize\color{muted}theorem; \code{KomorebiPaperIV/CanonicalProductMultiplicity.lean}, line 118.\par}\smallskip
{\small\code{Komorebi.PaperIV.entireReducedDenominator_full_product_same_orders}\par}
{\footnotesize\color{muted}theorem; \code{KomorebiPaperIV/CanonicalProductMultiplicity.lean}, line 128.\par}\smallskip
\Needspace{7\baselineskip}
\subsection{Row 98: 12.1 Entire nonvanishing quotient}
\label{row:98}
\textbf{Exact row scope.} Removes all quotient singularities using the proved order equality. Constructs a globally entire, nowhere-zero quotient with value 1 at zero, and proves Delta equals the full product times this quotient. Constancy is proved in row 102.\par\smallskip
\textbf{Manuscript labels.} \code{thm:genus_zero_product}\par\smallskip
{\small\code{Komorebi.PaperIV.canonicalProductQuotient}\par}
{\footnotesize\color{muted}definition; \code{KomorebiPaperIV/CanonicalProductQuotient.lean}, line 12.\par}\smallskip
{\small\code{Komorebi.PaperIV.rawCanonicalProductQuotient_meromorphic}\par}
{\footnotesize\color{muted}theorem; \code{KomorebiPaperIV/CanonicalProductQuotient.lean}, line 15.\par}\smallskip
{\small\code{Komorebi.PaperIV.rawCanonicalProductQuotient_order_zero}\par}
{\footnotesize\color{muted}theorem; \code{KomorebiPaperIV/CanonicalProductQuotient.lean}, line 20.\par}\smallskip
{\small\code{Komorebi.PaperIV.canonicalProductQuotient_entire}\par}
{\footnotesize\color{muted}theorem; \code{KomorebiPaperIV/CanonicalProductQuotient.lean}, line 32.\par}\smallskip
{\small\code{Komorebi.PaperIV.canonicalProductQuotient_ne_zero}\par}
{\footnotesize\color{muted}theorem; \code{KomorebiPaperIV/CanonicalProductQuotient.lean}, line 40.\par}\smallskip
{\small\code{Komorebi.PaperIV.canonicalProductQuotient_eq_div}\par}
{\footnotesize\color{muted}theorem; \code{KomorebiPaperIV/CanonicalProductQuotient.lean}, line 47.\par}\smallskip
{\small\code{Komorebi.PaperIV.fullCanonicalProduct_times_quotient}\par}
{\footnotesize\color{muted}theorem; \code{KomorebiPaperIV/CanonicalProductQuotient.lean}, line 58.\par}\smallskip
{\small\code{Komorebi.PaperIV.canonicalProductQuotient_zero}\par}
{\footnotesize\color{muted}theorem; \code{KomorebiPaperIV/CanonicalProductQuotient.lean}, line 66.\par}\smallskip
{\small\code{Komorebi.PaperIV.entireReducedDenominator_product_quotient}\par}
{\footnotesize\color{muted}theorem; \code{KomorebiPaperIV/CanonicalProductQuotient.lean}, line 71.\par}\smallskip
\Needspace{7\baselineskip}
\subsection{Row 99: 12.1 Growth reduction for the quotient}
\label{row:99}
\textbf{Exact row scope.} Constructs an entire normalized logarithm of the quotient. Proves a Borel-Caratheodory rigidity criterion from a uniform disk log-norm bound M(R) with M(R)/R tending to zero. The required bound for the actual quotient is derived in rows 100--102, completing Theorem 12.1.\par\smallskip
\textbf{Manuscript labels.} \code{thm:genus_zero_product}\par\smallskip
{\small\code{Komorebi.PaperIV.entire_normalized_exists_log}\par}
{\footnotesize\color{muted}theorem; \code{KomorebiPaperIV/SublinearExponentialRigidity.lean}, line 16.\par}\smallskip
{\small\code{Komorebi.PaperIV.entire_realPart_sublinear_constant}\par}
{\footnotesize\color{muted}theorem; \code{KomorebiPaperIV/SublinearExponentialRigidity.lean}, line 41.\par}\smallskip
{\small\code{Komorebi.PaperIV.entire_zeroFree_sublinear_logNorm_eq_one}\par}
{\footnotesize\color{muted}theorem; \code{KomorebiPaperIV/SublinearExponentialRigidity.lean}, line 64.\par}\smallskip
{\small\code{Komorebi.PaperIV.entireReducedDenominator_product_quotient_has_log}\par}
{\footnotesize\color{muted}theorem; \code{KomorebiPaperIV/SublinearExponentialRigidity.lean}, line 82.\par}\smallskip
\Needspace{7\baselineskip}
\subsection{Row 100: 12.1 Jensen-Poisson quotient estimate}
\label{row:100}
\textbf{Exact row scope.} Proves a normalized entire reciprocal proximity bound by Jensen, quotient proximity by division, and a Poisson real-part estimate. These yield rigidity for a factorization with proved sublinear bounds on both entire factors.\par\smallskip
\textbf{Manuscript labels.} \code{thm:genus_zero_product}\par\smallskip
{\small\code{Komorebi.PaperIV.entire_logCounting_poles_zero}\par}
{\footnotesize\color{muted}theorem; \code{KomorebiPaperIV/EntireQuotientGrowth.lean}, line 14.\par}\smallskip
{\small\code{Komorebi.PaperIV.entire_normalized_proximity_inv_le}\par}
{\footnotesize\color{muted}theorem; \code{KomorebiPaperIV/EntireQuotientGrowth.lean}, line 20.\par}\smallskip
{\small\code{Komorebi.PaperIV.entire_eventually_codiscrete_ne_zero}\par}
{\footnotesize\color{muted}theorem; \code{KomorebiPaperIV/EntireQuotientGrowth.lean}, line 32.\par}\smallskip
{\small\code{Komorebi.PaperIV.entire_quotient_proximity_le}\par}
{\footnotesize\color{muted}theorem; \code{KomorebiPaperIV/EntireQuotientGrowth.lean}, line 45.\par}\smallskip
{\small\code{Komorebi.PaperIV.entire_re_le_three_proximity_exp}\par}
{\footnotesize\color{muted}theorem; \code{KomorebiPaperIV/EntireQuotientGrowth.lean}, line 62.\par}\smallskip
{\small\code{Komorebi.PaperIV.entire_proximity_le_of_growth}\par}
{\footnotesize\color{muted}theorem; \code{KomorebiPaperIV/EntireQuotientGrowth.lean}, line 98.\par}\smallskip
{\small\code{Komorebi.PaperIV.entire_sublinear_factor_rigidity}\par}
{\footnotesize\color{muted}theorem; \code{KomorebiPaperIV/EntireQuotientGrowth.lean}, line 113.\par}\smallskip
\Needspace{7\baselineskip}
\subsection{Row 101: 12.1 Actual denominator and product growth}
\label{row:101}
\textbf{Exact row scope.} Derives the needed growth bounds from actual zero-power summability and the entire denominator estimate. The fixed exponent 3/4 is valid for every N>=2. No quotient-growth certificate is assumed.\par\smallskip
\textbf{Manuscript labels.} \code{thm:genus_zero_product}\par\smallskip
{\small\code{Komorebi.PaperIV.log_one_add_le_rpow_div}\par}
{\footnotesize\color{muted}theorem; \code{KomorebiPaperIV/CanonicalProductGrowth.lean}, line 14.\par}\smallskip
{\small\code{Komorebi.PaperIV.zeroMultiplicity_reciprocal_rpow_summable}\par}
{\footnotesize\color{muted}theorem; \code{KomorebiPaperIV/CanonicalProductGrowth.lean}, line 23.\par}\smallskip
{\small\code{Komorebi.PaperIV.reciprocalPowerMass}\par}
{\footnotesize\color{muted}definition; \code{KomorebiPaperIV/CanonicalProductGrowth.lean}, line 33.\par}\smallskip
{\small\code{Komorebi.PaperIV.reciprocalPowerMass_nonneg}\par}
{\footnotesize\color{muted}theorem; \code{KomorebiPaperIV/CanonicalProductGrowth.lean}, line 36.\par}\smallskip
{\small\code{Komorebi.PaperIV.fullCanonicalProduct_growth}\par}
{\footnotesize\color{muted}theorem; \code{KomorebiPaperIV/CanonicalProductGrowth.lean}, line 40.\par}\smallskip
{\small\code{Komorebi.PaperIV.entireReducedDenominator_norm_le}\par}
{\footnotesize\color{muted}theorem; \code{KomorebiPaperIV/CanonicalProductGrowth.lean}, line 82.\par}\smallskip
{\small\code{Komorebi.PaperIV.entireReducedDenominator_and_product_growth}\par}
{\footnotesize\color{muted}theorem; \code{KomorebiPaperIV/CanonicalProductGrowth.lean}, line 91.\par}\smallskip
\Needspace{7\baselineskip}
\subsection{Row 102: 12.1 Full Hadamard identification}
\label{row:102}
\textbf{Exact row scope.} The actual normalized quotient equals one. Proves Delta equals the complete multiplicity-indexed product, compact-uniform convergence to Delta, and the original-variable product for D\_N. Fully encodes Theorem 12.1.\par\smallskip
\textbf{Manuscript labels.} \code{thm:genus_zero_product}\par\smallskip
{\small\code{Komorebi.PaperIV.entireReducedDenominator_quotient_eq_one}\par}
{\footnotesize\color{muted}theorem; \code{KomorebiPaperIV/HadamardIdentification.lean}, line 15.\par}\smallskip
{\small\code{Komorebi.PaperIV.entireReducedDenominator_eq_fullCanonicalProduct}\par}
{\footnotesize\color{muted}theorem; \code{KomorebiPaperIV/HadamardIdentification.lean}, line 32.\par}\smallskip
{\small\code{Komorebi.PaperIV.entireReducedDenominator_hadamard_uniform}\par}
{\footnotesize\color{muted}theorem; \code{KomorebiPaperIV/HadamardIdentification.lean}, line 41.\par}\smallskip
{\small\code{Komorebi.PaperIV.canonicalSignedDenominator_hadamard}\par}
{\footnotesize\color{muted}theorem; \code{KomorebiPaperIV/HadamardIdentification.lean}, line 50.\par}\smallskip
\Needspace{7\baselineskip}
\subsection{Row 103: 12.7 Convergent Mittag-Leffler expansion}
\label{row:103}
\textbf{Exact row scope.} Proves the logarithmic derivative of the full product and the exact tangent partial fractions, with absolute and locally uniform convergence away from poles. Auxiliary tail-convergence lemmas use Lean total division at poles; the analytic identity is restricted to nonpoles. Fully encodes Corollary 12.7.\par\smallskip
\textbf{Manuscript labels.} \code{cor:mittag_leffler_tangent}\par\smallskip
{\small\code{Komorebi.PaperIV.zeroPartialFraction}\par}
{\footnotesize\color{muted}definition; \code{KomorebiPaperIV/CanonicalPartialFractions.lean}, line 14.\par}\smallskip
{\small\code{Komorebi.PaperIV.zeroPartialFraction_uniform_majorant}\par}
{\footnotesize\color{muted}theorem; \code{KomorebiPaperIV/CanonicalPartialFractions.lean}, line 17.\par}\smallskip
{\small\code{Komorebi.PaperIV.zeroPartialFraction_norm_summable}\par}
{\footnotesize\color{muted}theorem; \code{KomorebiPaperIV/CanonicalPartialFractions.lean}, line 34.\par}\smallskip
{\small\code{Komorebi.PaperIV.zeroPartialFraction_uniform_on_compact}\par}
{\footnotesize\color{muted}theorem; \code{KomorebiPaperIV/CanonicalPartialFractions.lean}, line 43.\par}\smallskip
{\small\code{Komorebi.PaperIV.logDeriv_linear_factor}\par}
{\footnotesize\color{muted}theorem; \code{KomorebiPaperIV/CanonicalPartialFractions.lean}, line 53.\par}\smallskip
{\small\code{Komorebi.PaperIV.fullCanonicalProduct_logDeriv}\par}
{\footnotesize\color{muted}theorem; \code{KomorebiPaperIV/CanonicalPartialFractions.lean}, line 61.\par}\smallskip
{\small\code{Komorebi.PaperIV.entireReducedDenominator_partialFractions}\par}
{\footnotesize\color{muted}theorem; \code{KomorebiPaperIV/CanonicalPartialFractions.lean}, line 84.\par}\smallskip
{\small\code{Komorebi.PaperIV.zeroPartialFraction_locally_uniform}\par}
{\footnotesize\color{muted}theorem; \code{KomorebiPaperIV/CanonicalPartialFractions.lean}, line 97.\par}\smallskip
{\small\code{Komorebi.PaperIV.analyticTransferTangent}\par}
{\footnotesize\color{muted}definition; \code{KomorebiPaperIV/CanonicalPartialFractions.lean}, line 105.\par}\smallskip
{\small\code{Komorebi.PaperIV.analyticTransferTangent_eq_reduced_logDeriv}\par}
{\footnotesize\color{muted}theorem; \code{KomorebiPaperIV/CanonicalPartialFractions.lean}, line 108.\par}\smallskip
{\small\code{Komorebi.PaperIV.analyticTransferTangent_mittagLeffler}\par}
{\footnotesize\color{muted}theorem; \code{KomorebiPaperIV/CanonicalPartialFractions.lean}, line 119.\par}\smallskip
{\small\code{Komorebi.PaperIV.analyticTransferTangent_mittagLeffler_absolute}\par}
{\footnotesize\color{muted}theorem; \code{KomorebiPaperIV/CanonicalPartialFractions.lean}, line 127.\par}\smallskip
{\small\code{Komorebi.PaperIV.analyticTransferTangent_mittagLeffler_locally_uniform}\par}
{\footnotesize\color{muted}theorem; \code{KomorebiPaperIV/CanonicalPartialFractions.lean}, line 135.\par}\smallskip
\Needspace{7\baselineskip}
\subsection{Row 104: 12.6 Actual zero moments and derivatives}
\label{row:104}
\textbf{Exact row scope.} Defines moments over all actual zeros repeated by analytic multiplicity. Proves termwise derivatives of the locally uniform resolvent series and the exact factorial-normalized derivatives at zero.\par\smallskip
\textbf{Manuscript labels.} \code{thm:newton_sums}\par\smallskip
{\small\code{Komorebi.PaperIV.locally_uniform_sum_deriv}\par}
{\footnotesize\color{muted}theorem; \code{KomorebiPaperIV/ZeroMomentDerivatives.lean}, line 13.\par}\smallskip
{\small\code{Komorebi.PaperIV.zeroResolventJet}\par}
{\footnotesize\color{muted}definition; \code{KomorebiPaperIV/ZeroMomentDerivatives.lean}, line 22.\par}\smallskip
{\small\code{Komorebi.PaperIV.zeroResolventJet_hasDerivAt}\par}
{\footnotesize\color{muted}theorem; \code{KomorebiPaperIV/ZeroMomentDerivatives.lean}, line 25.\par}\smallskip
{\small\code{Komorebi.PaperIV.zeroResolventJet_locally_uniform}\par}
{\footnotesize\color{muted}theorem; \code{KomorebiPaperIV/ZeroMomentDerivatives.lean}, line 37.\par}\smallskip
{\small\code{Komorebi.PaperIV.zeroPowerSum}\par}
{\footnotesize\color{muted}definition; \code{KomorebiPaperIV/ZeroMomentDerivatives.lean}, line 65.\par}\smallskip
{\small\code{Komorebi.PaperIV.zeroPowerSum_hasSum_from_resolvent}\par}
{\footnotesize\color{muted}theorem; \code{KomorebiPaperIV/ZeroMomentDerivatives.lean}, line 68.\par}\smallskip
{\small\code{Komorebi.PaperIV.zeroPowerSum_resolvent_derivative}\par}
{\footnotesize\color{muted}theorem; \code{KomorebiPaperIV/ZeroMomentDerivatives.lean}, line 75.\par}\smallskip
\Needspace{7\baselineskip}
\subsection{Row 105: 12.6 Kappa identification and Newton recurrence}
\label{row:105}
\textbf{Exact row scope.} Identifies every positive-index formal kappa coefficient with the actual zero moment and proves the exact Newton recurrence. The logarithm-series clause is supplied by row 109.\par\smallskip
\textbf{Manuscript labels.} \code{thm:newton_sums}\par\smallskip
{\small\code{Komorebi.PaperIV.entireReducedDenominator_iteratedDeriv_zero}\par}
{\footnotesize\color{muted}theorem; \code{KomorebiPaperIV/ZeroMomentCoefficients.lean}, line 14.\par}\smallskip
{\small\code{Komorebi.PaperIV.zeroResolvent_analyticAt_zero}\par}
{\footnotesize\color{muted}theorem; \code{KomorebiPaperIV/ZeroMomentCoefficients.lean}, line 26.\par}\smallskip
{\small\code{Komorebi.PaperIV.reducedZeroPowerSum_convolution}\par}
{\footnotesize\color{muted}theorem; \code{KomorebiPaperIV/ZeroMomentCoefficients.lean}, line 43.\par}\smallskip
{\small\code{Komorebi.PaperIV.zeroMomentSeries}\par}
{\footnotesize\color{muted}definition; \code{KomorebiPaperIV/ZeroMomentCoefficients.lean}, line 90.\par}\smallskip
{\small\code{Komorebi.PaperIV.reducedDenominator_mul_zeroMomentSeries}\par}
{\footnotesize\color{muted}theorem; \code{KomorebiPaperIV/ZeroMomentCoefficients.lean}, line 93.\par}\smallskip
{\small\code{Komorebi.PaperIV.kappaCoeff_eq_zeroPowerSum}\par}
{\footnotesize\color{muted}theorem; \code{KomorebiPaperIV/ZeroMomentCoefficients.lean}, line 103.\par}\smallskip
{\small\code{Komorebi.PaperIV.reducedZeroPowerSum_newton}\par}
{\footnotesize\color{muted}theorem; \code{KomorebiPaperIV/ZeroMomentCoefficients.lean}, line 114.\par}\smallskip
\Needspace{7\baselineskip}
\subsection{Row 106: 12.9 Tangent coefficients from zero sums}
\label{row:106}
\textbf{Exact row scope.} Proves the precise factorial-normalized coefficient formula and its convergent Taylor realization as the actual analytic tangent on every zero-free disk in the reduced variable. Fully encodes Corollary 12.9.\par\smallskip
\textbf{Manuscript labels.} \code{cor:tangent_zero_sums}\par\smallskip
{\small\code{Komorebi.PaperIV.transferTangentCoeff_zero_sum}\par}
{\footnotesize\color{muted}theorem; \code{KomorebiPaperIV/ZeroMomentCoefficients.lean}, line 135.\par}\smallskip
{\small\code{Komorebi.PaperIV.analyticTransferTangent_coefficient_hasSum}\par}
{\footnotesize\color{muted}theorem; \code{KomorebiPaperIV/ZeroMomentLogarithm.lean}, line 110.\par}\smallskip
\Needspace{7\baselineskip}
\subsection{Row 107: 12.10 Reciprocal Bell formula from zero data}
\label{row:107}
\textbf{Exact row scope.} The named reciprocal coefficient family equals the complete exponential Bell polynomial in the actual zero moments, for all n including zero and every complex parameter r. Fully encodes Corollary 12.10; earlier general-domain complex-power bridges retain their own boundaries.\par\smallskip
\textbf{Manuscript labels.} \code{cor:reciprocal_zero_data}\par\smallskip
{\small\code{Komorebi.PaperIV.bellInput_zero_sums}\par}
{\footnotesize\color{muted}theorem; \code{KomorebiPaperIV/ZeroMomentCoefficients.lean}, line 140.\par}\smallskip
{\small\code{Komorebi.PaperIV.reciprocalCoeff_zero_sum_bell}\par}
{\footnotesize\color{muted}theorem; \code{KomorebiPaperIV/ZeroMomentCoefficients.lean}, line 150.\par}\smallskip
\Needspace{7\baselineskip}
\subsection{Row 108: 12.8 Theta Bell formula from zero data}
\label{row:108}
\textbf{Exact row scope.} Uses formal ODE uniqueness and the proved kappa/zero-moment identity to derive the exact theta coefficient formula, including the alternating sign and u\textasciicircum{}(Nm)+u factor. Fully encodes Proposition 12.8 for the named theta coefficient family.\par\smallskip
\textbf{Manuscript labels.} \code{prop:theta_bell}\par\smallskip
{\small\code{Komorebi.PaperIV.powerSeries_linear_ode_unique}\par}
{\footnotesize\color{muted}theorem; \code{KomorebiPaperIV/ThetaZeroBell.lean}, line 10.\par}\smallskip
{\small\code{Komorebi.PaperIV.thetaZeroBellInput}\par}
{\footnotesize\color{muted}definition; \code{KomorebiPaperIV/ThetaZeroBell.lean}, line 29.\par}\smallskip
{\small\code{Komorebi.PaperIV.thetaZeroBellInput_derivative}\par}
{\footnotesize\color{muted}theorem; \code{KomorebiPaperIV/ThetaZeroBell.lean}, line 33.\par}\smallskip
{\small\code{Komorebi.PaperIV.thetaSeries_eq_zero_bell_exp}\par}
{\footnotesize\color{muted}theorem; \code{KomorebiPaperIV/ThetaZeroBell.lean}, line 44.\par}\smallskip
{\small\code{Komorebi.PaperIV.thetaCoeff_zero_sum_bell}\par}
{\footnotesize\color{muted}theorem; \code{KomorebiPaperIV/ThetaZeroBell.lean}, line 60.\par}\smallskip
\Needspace{7\baselineskip}
\subsection{Row 109: 12.6 Normalized analytic logarithm series}
\label{row:109}
\textbf{Exact row scope.} Constructs the logarithm branch normalized at zero on every zero-free disk, proves exp(-L)=Delta, and proves its convergent zero-moment series throughout that disk using Cauchy estimates. Together with rows 104--105 this fully encodes Theorem 12.6. Also proves the convergent resolvent Taylor series.\par\smallskip
\textbf{Manuscript labels.} \code{thm:newton_sums}\par\smallskip
{\small\code{Komorebi.PaperIV.differentiableOn_ball_taylor_hasSum}\par}
{\footnotesize\color{muted}theorem; \code{KomorebiPaperIV/ZeroMomentLogarithm.lean}, line 14.\par}\smallskip
{\small\code{Komorebi.PaperIV.reducedZeroPowerSum_logarithm}\par}
{\footnotesize\color{muted}theorem; \code{KomorebiPaperIV/ZeroMomentLogarithm.lean}, line 34.\par}\smallskip
{\small\code{Komorebi.PaperIV.reducedZeroPowerSum_resolvent_hasSum}\par}
{\footnotesize\color{muted}theorem; \code{KomorebiPaperIV/ZeroMomentLogarithm.lean}, line 87.\par}\smallskip
\Needspace{7\baselineskip}
\subsection{Row 110: 11.2 Dominant exponential on a vertex ray}
\label{row:110}
\textbf{Exact row scope.} Proves finite exponential dominance with an explicit positive projection gap and specializes to the actual canonical signed denominator on a vertex-facing ray. No dominance or lower-growth certificate is assumed.\par\smallskip
\textbf{Manuscript labels.} \code{prop:orders}\par\smallskip
{\small\code{Komorebi.PaperIV.finite_positive_gap}\par}
{\footnotesize\color{muted}theorem; \code{KomorebiPaperIV/DominantExponential.lean}, line 14.\par}\smallskip
{\small\code{Komorebi.PaperIV.finite_exponential_dominant_bound}\par}
{\footnotesize\color{muted}theorem; \code{KomorebiPaperIV/DominantExponential.lean}, line 21.\par}\smallskip
{\small\code{Komorebi.PaperIV.finite_exponential_dominant_tendsto}\par}
{\footnotesize\color{muted}theorem; \code{KomorebiPaperIV/DominantExponential.lean}, line 43.\par}\smallskip
{\small\code{Komorebi.PaperIV.unit_complex_re_lt_one}\par}
{\footnotesize\color{muted}theorem; \code{KomorebiPaperIV/DominantExponential.lean}, line 59.\par}\smallskip
{\small\code{Komorebi.PaperIV.canonicalSignedVertex_injective}\par}
{\footnotesize\color{muted}theorem; \code{KomorebiPaperIV/DominantExponential.lean}, line 71.\par}\smallskip
{\small\code{Komorebi.PaperIV.canonicalSignedVertex_mul_star_self}\par}
{\footnotesize\color{muted}theorem; \code{KomorebiPaperIV/DominantExponential.lean}, line 78.\par}\smallskip
{\small\code{Komorebi.PaperIV.canonicalSignedDenominator_dominant_ray}\par}
{\footnotesize\color{muted}theorem; \code{KomorebiPaperIV/DominantExponential.lean}, line 84.\par}\smallskip
\Needspace{7\baselineskip}
\subsection{Row 111: 11.2 Sharp growth order and exponential type}
\label{row:111}
\textbf{Exact row scope.} Defines the closed-disk maximum modulus and the standard logarithmic order/type limsups. Proves two-sided maximum-modulus bounds, actual limits for both logarithmic growth ratios, order one and type one for D\_N, and order 1/N for Delta\_N. Fully encodes Proposition 11.2; imported standard fact 11.1 is not separately formalized or assumed by this proof.\par\smallskip
\textbf{Manuscript labels.} \code{prop:orders}\par\smallskip
{\small\code{Komorebi.PaperIV.diskMaximumModulus}\par}
{\footnotesize\color{muted}definition; \code{KomorebiPaperIV/SharpGrowth.lean}, line 15.\par}\smallskip
{\small\code{Komorebi.PaperIV.entireGrowthOrder}\par}
{\footnotesize\color{muted}definition; \code{KomorebiPaperIV/SharpGrowth.lean}, line 19.\par}\smallskip
{\small\code{Komorebi.PaperIV.entireExponentialType}\par}
{\footnotesize\color{muted}definition; \code{KomorebiPaperIV/SharpGrowth.lean}, line 23.\par}\smallskip
{\small\code{Komorebi.PaperIV.norm_le_diskMaximumModulus}\par}
{\footnotesize\color{muted}theorem; \code{KomorebiPaperIV/SharpGrowth.lean}, line 26.\par}\smallskip
{\small\code{Komorebi.PaperIV.diskMaximumModulus_le}\par}
{\footnotesize\color{muted}theorem; \code{KomorebiPaperIV/SharpGrowth.lean}, line 31.\par}\smallskip
{\small\code{Komorebi.PaperIV.log_growth_ratio_of_exp_bounds}\par}
{\footnotesize\color{muted}theorem; \code{KomorebiPaperIV/SharpGrowth.lean}, line 39.\par}\smallskip
{\small\code{Komorebi.PaperIV.order_ratio_of_log_growth}\par}
{\footnotesize\color{muted}theorem; \code{KomorebiPaperIV/SharpGrowth.lean}, line 60.\par}\smallskip
{\small\code{Komorebi.PaperIV.canonicalSignedDenominator_ray_lower}\par}
{\footnotesize\color{muted}theorem; \code{KomorebiPaperIV/SharpGrowth.lean}, line 77.\par}\smallskip
{\small\code{Komorebi.PaperIV.canonicalSignedDenominator_maximum_bounds}\par}
{\footnotesize\color{muted}theorem; \code{KomorebiPaperIV/SharpGrowth.lean}, line 100.\par}\smallskip
{\small\code{Komorebi.PaperIV.canonicalSignedDenominator_type_limit}\par}
{\footnotesize\color{muted}theorem; \code{KomorebiPaperIV/SharpGrowth.lean}, line 113.\par}\smallskip
{\small\code{Komorebi.PaperIV.canonicalSignedDenominator_order_limit}\par}
{\footnotesize\color{muted}theorem; \code{KomorebiPaperIV/SharpGrowth.lean}, line 119.\par}\smallskip
{\small\code{Komorebi.PaperIV.canonicalSignedDenominator_order_one}\par}
{\footnotesize\color{muted}theorem; \code{KomorebiPaperIV/SharpGrowth.lean}, line 124.\par}\smallskip
{\small\code{Komorebi.PaperIV.canonicalSignedDenominator_type_one}\par}
{\footnotesize\color{muted}theorem; \code{KomorebiPaperIV/SharpGrowth.lean}, line 128.\par}\smallskip
{\small\code{Komorebi.PaperIV.entireReducedDenominator_maximum_bounds}\par}
{\footnotesize\color{muted}theorem; \code{KomorebiPaperIV/SharpGrowth.lean}, line 133.\par}\smallskip
{\small\code{Komorebi.PaperIV.entireReducedDenominator_growth_limit}\par}
{\footnotesize\color{muted}theorem; \code{KomorebiPaperIV/SharpGrowth.lean}, line 159.\par}\smallskip
{\small\code{Komorebi.PaperIV.entireReducedDenominator_order_limit}\par}
{\footnotesize\color{muted}theorem; \code{KomorebiPaperIV/SharpGrowth.lean}, line 164.\par}\smallskip
{\small\code{Komorebi.PaperIV.entireReducedDenominator_order_exact}\par}
{\footnotesize\color{muted}theorem; \code{KomorebiPaperIV/SharpGrowth.lean}, line 170.\par}\smallskip
\Needspace{7\baselineskip}
\subsection{Row 112: 11.4 Indicator polygon and tie geometry}
\label{row:112}
\textbf{Exact row scope.} Identifies the actual support function with the manuscript maximum of cosines, proves continuity and nonnegativity for N>=2, and derives the exact periodic tie directions and strict maximizing vertex off those directions.\par\smallskip
\textbf{Manuscript labels.} \code{thm:indicator_polygon}\par\smallskip
{\small\code{Komorebi.PaperIV.rayDirection}\par}
{\footnotesize\color{muted}definition; \code{KomorebiPaperIV/IndicatorGeometry.lean}, line 13.\par}\smallskip
{\small\code{Komorebi.PaperIV.norm_rayDirection}\par}
{\footnotesize\color{muted}theorem; \code{KomorebiPaperIV/IndicatorGeometry.lean}, line 15.\par}\smallskip
{\small\code{Komorebi.PaperIV.signedProjection}\par}
{\footnotesize\color{muted}definition; \code{KomorebiPaperIV/IndicatorGeometry.lean}, line 18.\par}\smallskip
{\small\code{Komorebi.PaperIV.signedIndicator}\par}
{\footnotesize\color{muted}definition; \code{KomorebiPaperIV/IndicatorGeometry.lean}, line 21.\par}\smallskip
{\small\code{Komorebi.PaperIV.signedProjection_cos}\par}
{\footnotesize\color{muted}theorem; \code{KomorebiPaperIV/IndicatorGeometry.lean}, line 24.\par}\smallskip
{\small\code{Komorebi.PaperIV.signedIndicator_cos}\par}
{\footnotesize\color{muted}theorem; \code{KomorebiPaperIV/IndicatorGeometry.lean}, line 32.\par}\smallskip
{\small\code{Komorebi.PaperIV.continuous_signedProjection}\par}
{\footnotesize\color{muted}theorem; \code{KomorebiPaperIV/IndicatorGeometry.lean}, line 40.\par}\smallskip
{\small\code{Komorebi.PaperIV.continuous_signedIndicator}\par}
{\footnotesize\color{muted}theorem; \code{KomorebiPaperIV/IndicatorGeometry.lean}, line 44.\par}\smallskip
{\small\code{Komorebi.PaperIV.signedProjection_le_indicator}\par}
{\footnotesize\color{muted}theorem; \code{KomorebiPaperIV/IndicatorGeometry.lean}, line 47.\par}\smallskip
{\small\code{Komorebi.PaperIV.canonicalSignedVertex_sum_zero}\par}
{\footnotesize\color{muted}theorem; \code{KomorebiPaperIV/IndicatorGeometry.lean}, line 51.\par}\smallskip
{\small\code{Komorebi.PaperIV.signedIndicator_nonneg}\par}
{\footnotesize\color{muted}theorem; \code{KomorebiPaperIV/IndicatorGeometry.lean}, line 56.\par}\smallskip
{\small\code{Komorebi.PaperIV.canonicalSignedVertex_pow}\par}
{\footnotesize\color{muted}theorem; \code{KomorebiPaperIV/IndicatorGeometry.lean}, line 65.\par}\smallskip
{\small\code{Komorebi.PaperIV.signed_root_in_vertices}\par}
{\footnotesize\color{muted}theorem; \code{KomorebiPaperIV/IndicatorGeometry.lean}, line 70.\par}\smallskip
{\small\code{Komorebi.PaperIV.unit_complex_same_real}\par}
{\footnotesize\color{muted}theorem; \code{KomorebiPaperIV/IndicatorGeometry.lean}, line 83.\par}\smallskip
{\small\code{Komorebi.PaperIV.signed_maximizer_strict_of_not_tie}\par}
{\footnotesize\color{muted}theorem; \code{KomorebiPaperIV/IndicatorGeometry.lean}, line 95.\par}\smallskip
{\small\code{Komorebi.PaperIV.signedTieDirection}\par}
{\footnotesize\color{muted}definition; \code{KomorebiPaperIV/IndicatorGeometry.lean}, line 131.\par}\smallskip
{\small\code{Komorebi.PaperIV.signedTieDirection_iff}\par}
{\footnotesize\color{muted}theorem; \code{KomorebiPaperIV/IndicatorGeometry.lean}, line 134.\par}\smallskip
\Needspace{7\baselineskip}
\subsection{Row 113: 11.4 Uniform dominance and zero-free sectors}
\label{row:113}
\textbf{Exact row scope.} Derives a uniform positive projection gap on every compact set avoiding the actual tie directions. Proves the normalized dominant-exponential error bound (N-1) exp(-delta r) and eventual zero-freeness there, with one uniform delta and radius. Supplies every dominant-sector clause of Theorem 11.4.\par\smallskip
\textbf{Manuscript labels.} \code{thm:indicator_polygon}\par\smallskip
{\small\code{Komorebi.PaperIV.range_erase_nonempty}\par}
{\footnotesize\color{muted}theorem; \code{KomorebiPaperIV/UniformDominantSectors.lean}, line 12.\par}\smallskip
{\small\code{Komorebi.PaperIV.vertexProjectionGap}\par}
{\footnotesize\color{muted}definition; \code{KomorebiPaperIV/UniformDominantSectors.lean}, line 18.\par}\smallskip
{\small\code{Komorebi.PaperIV.directionSeparation}\par}
{\footnotesize\color{muted}definition; \code{KomorebiPaperIV/UniformDominantSectors.lean}, line 22.\par}\smallskip
{\small\code{Komorebi.PaperIV.continuous_vertexProjectionGap}\par}
{\footnotesize\color{muted}theorem; \code{KomorebiPaperIV/UniformDominantSectors.lean}, line 26.\par}\smallskip
{\small\code{Komorebi.PaperIV.continuous_directionSeparation}\par}
{\footnotesize\color{muted}theorem; \code{KomorebiPaperIV/UniformDominantSectors.lean}, line 30.\par}\smallskip
{\small\code{Komorebi.PaperIV.directionSeparation_pos}\par}
{\footnotesize\color{muted}theorem; \code{KomorebiPaperIV/UniformDominantSectors.lean}, line 34.\par}\smallskip
{\small\code{Komorebi.PaperIV.compact_directions_uniform_gap}\par}
{\footnotesize\color{muted}theorem; \code{KomorebiPaperIV/UniformDominantSectors.lean}, line 49.\par}\smallskip
{\small\code{Komorebi.PaperIV.canonicalSignedDenominator_dominant_error}\par}
{\footnotesize\color{muted}theorem; \code{KomorebiPaperIV/UniformDominantSectors.lean}, line 65.\par}\smallskip
{\small\code{Komorebi.PaperIV.canonicalSignedDenominator_compact_dominant_sectors}\par}
{\footnotesize\color{muted}theorem; \code{KomorebiPaperIV/UniformDominantSectors.lean}, line 90.\par}\smallskip
\Needspace{7\baselineskip}
\subsection{Row 114: 11.4 Aligned radii in tie directions}
\label{row:114}
\textbf{Exact row scope.} Constructs explicit radii tending to infinity on which the leading conjugate phases align; lower-projection terms decay and the normalized finite sum has a positive limit. This proves the indicator lower bound even at tie directions, without a phase-alignment certificate.\par\smallskip
\textbf{Manuscript labels.} \code{thm:indicator_polygon}\par\smallskip
{\small\code{Komorebi.PaperIV.alignedRadii}\par}
{\footnotesize\color{muted}definition; \code{KomorebiPaperIV/AlignedExponentialRadii.lean}, line 14.\par}\smallskip
{\small\code{Komorebi.PaperIV.alignedRadii_tendsto}\par}
{\footnotesize\color{muted}theorem; \code{KomorebiPaperIV/AlignedExponentialRadii.lean}, line 17.\par}\smallskip
{\small\code{Komorebi.PaperIV.alignedRadii_phase}\par}
{\footnotesize\color{muted}theorem; \code{KomorebiPaperIV/AlignedExponentialRadii.lean}, line 26.\par}\smallskip
{\small\code{Komorebi.PaperIV.alignedRadii_equal_projection}\par}
{\footnotesize\color{muted}theorem; \code{KomorebiPaperIV/AlignedExponentialRadii.lean}, line 47.\par}\smallskip
{\small\code{Komorebi.PaperIV.alignedRadii_lower_projection}\par}
{\footnotesize\color{muted}theorem; \code{KomorebiPaperIV/AlignedExponentialRadii.lean}, line 65.\par}\smallskip
{\small\code{Komorebi.PaperIV.finite_unit_exponential_aligned_limit}\par}
{\footnotesize\color{muted}theorem; \code{KomorebiPaperIV/AlignedExponentialRadii.lean}, line 79.\par}\smallskip
{\small\code{Komorebi.PaperIV.log_norm_slope_of_normalized_limit}\par}
{\footnotesize\color{muted}theorem; \code{KomorebiPaperIV/AlignedExponentialRadii.lean}, line 99.\par}\smallskip
\Needspace{7\baselineskip}
\subsection{Row 115: 11.4 Exact indicator in every direction}
\label{row:115}
\textbf{Exact row scope.} Proves the exact indicator limsup for every angle and every N>=2. Includes the classical formulation along radii where the denominator is nonzero, so the theorem does not depend on Lean Real.log zero. Together with rows 112--114 this fully encodes Theorem 11.4, including uniform compact-sector estimates and eventual zero-freeness.\par\smallskip
\textbf{Manuscript labels.} \code{thm:indicator_polygon}\par\smallskip
{\small\code{Komorebi.PaperIV.canonicalSignedDenominator_indicator_upper}\par}
{\footnotesize\color{muted}theorem; \code{KomorebiPaperIV/IndicatorLimit.lean}, line 14.\par}\smallskip
{\small\code{Komorebi.PaperIV.canonicalSignedDenominator_log_indicator_upper}\par}
{\footnotesize\color{muted}theorem; \code{KomorebiPaperIV/IndicatorLimit.lean}, line 38.\par}\smallskip
{\small\code{Komorebi.PaperIV.canonicalSignedDenominator_real_normalization}\par}
{\footnotesize\color{muted}theorem; \code{KomorebiPaperIV/IndicatorLimit.lean}, line 50.\par}\smallskip
{\small\code{Komorebi.PaperIV.canonicalSignedDenominator_indicator_sequence}\par}
{\footnotesize\color{muted}theorem; \code{KomorebiPaperIV/IndicatorLimit.lean}, line 67.\par}\smallskip
{\small\code{Komorebi.PaperIV.limsup_eq_of_upper_and_sequence}\par}
{\footnotesize\color{muted}theorem; \code{KomorebiPaperIV/IndicatorLimit.lean}, line 95.\par}\smallskip
{\small\code{Komorebi.PaperIV.canonicalSignedDenominator_indicator}\par}
{\footnotesize\color{muted}theorem; \code{KomorebiPaperIV/IndicatorLimit.lean}, line 104.\par}\smallskip
{\small\code{Komorebi.PaperIV.canonicalSignedDenominator_indicator_nonzero_radii}\par}
{\footnotesize\color{muted}theorem; \code{KomorebiPaperIV/IndicatorLimit.lean}, line 113.\par}\smallskip
\Needspace{7\baselineskip}
\subsection{Row 116: 12.2 DiagonalNuclear}
\label{row:116}
\textbf{Exact row scope.} Actual diagonal operator on l2, norm-summable rank-one expansion, Hilbert nuclear/trace-class characterization, and exact basis and power actions.\par\smallskip
\textbf{Manuscript labels.} \code{thm:trace_class_model}\par\smallskip
{\small\code{Komorebi.PaperIV.IsHilbertTraceClass}\par}
{\footnotesize\color{muted}definition; \code{KomorebiPaperIV/DiagonalNuclear.lean}, line 15.\par}\smallskip
{\small\code{Komorebi.PaperIV.diagonalHilbertBasis}\par}
{\footnotesize\color{muted}definition; \code{KomorebiPaperIV/DiagonalNuclear.lean}, line 20.\par}\smallskip
{\small\code{Komorebi.PaperIV.diagonalHilbertBasis_inner}\par}
{\footnotesize\color{muted}theorem; \code{KomorebiPaperIV/DiagonalNuclear.lean}, line 23.\par}\smallskip
{\small\code{Komorebi.PaperIV.diagonalHilbertBasis_norm}\par}
{\footnotesize\color{muted}theorem; \code{KomorebiPaperIV/DiagonalNuclear.lean}, line 27.\par}\smallskip
{\small\code{Komorebi.PaperIV.diagonalHilbertBasis_apply}\par}
{\footnotesize\color{muted}theorem; \code{KomorebiPaperIV/DiagonalNuclear.lean}, line 30.\par}\smallskip
{\small\code{Komorebi.PaperIV.diagonalRankOne}\par}
{\footnotesize\color{muted}definition; \code{KomorebiPaperIV/DiagonalNuclear.lean}, line 35.\par}\smallskip
{\small\code{Komorebi.PaperIV.diagonalRankOne_norm}\par}
{\footnotesize\color{muted}theorem; \code{KomorebiPaperIV/DiagonalNuclear.lean}, line 39.\par}\smallskip
{\small\code{Komorebi.PaperIV.diagonalRankOne_summable}\par}
{\footnotesize\color{muted}theorem; \code{KomorebiPaperIV/DiagonalNuclear.lean}, line 42.\par}\smallskip
{\small\code{Komorebi.PaperIV.diagonalOperator}\par}
{\footnotesize\color{muted}definition; \code{KomorebiPaperIV/DiagonalNuclear.lean}, line 47.\par}\smallskip
{\small\code{Komorebi.PaperIV.diagonalOperator_traceClass}\par}
{\footnotesize\color{muted}theorem; \code{KomorebiPaperIV/DiagonalNuclear.lean}, line 50.\par}\smallskip
{\small\code{Komorebi.PaperIV.diagonalOperator_apply}\par}
{\footnotesize\color{muted}theorem; \code{KomorebiPaperIV/DiagonalNuclear.lean}, line 56.\par}\smallskip
{\small\code{Komorebi.PaperIV.diagonalOperator_apply_basis}\par}
{\footnotesize\color{muted}theorem; \code{KomorebiPaperIV/DiagonalNuclear.lean}, line 69.\par}\smallskip
{\small\code{Komorebi.PaperIV.diagonalOperator_pow_apply}\par}
{\footnotesize\color{muted}theorem; \code{KomorebiPaperIV/DiagonalNuclear.lean}, line 76.\par}\smallskip
\Needspace{7\baselineskip}
\subsection{Row 117: 12.2 DiagonalTrace}
\label{row:117}
\textbf{Exact row scope.} Parseval and absolutely convergent double sums prove the operator trace in every orthonormal Hilbert basis, including powers.\par\smallskip
\textbf{Manuscript labels.} \code{thm:trace_class_model}\par\smallskip
{\small\code{Komorebi.PaperIV.hilbertBasis_hasSum_squared_coordinates}\par}
{\footnotesize\color{muted}theorem; \code{KomorebiPaperIV/DiagonalTrace.lean}, line 13.\par}\smallskip
{\small\code{Komorebi.PaperIV.diagonal_trace_double_summable}\par}
{\footnotesize\color{muted}theorem; \code{KomorebiPaperIV/DiagonalTrace.lean}, line 19.\par}\smallskip
{\small\code{Komorebi.PaperIV.hilbertBasisTrace}\par}
{\footnotesize\color{muted}definition; \code{KomorebiPaperIV/DiagonalTrace.lean}, line 36.\par}\smallskip
{\small\code{Komorebi.PaperIV.diagonalOperator_trace_in_any_basis}\par}
{\footnotesize\color{muted}theorem; \code{KomorebiPaperIV/DiagonalTrace.lean}, line 39.\par}\smallskip
{\small\code{Komorebi.PaperIV.diagonalOperator_trace_eq}\par}
{\footnotesize\color{muted}theorem; \code{KomorebiPaperIV/DiagonalTrace.lean}, line 73.\par}\smallskip
{\small\code{Komorebi.PaperIV.diagonalOperator_pow_eq}\par}
{\footnotesize\color{muted}theorem; \code{KomorebiPaperIV/DiagonalTrace.lean}, line 78.\par}\smallskip
{\small\code{Komorebi.PaperIV.diagonalOperator_pow_trace}\par}
{\footnotesize\color{muted}theorem; \code{KomorebiPaperIV/DiagonalTrace.lean}, line 84.\par}\smallskip
\Needspace{7\baselineskip}
\subsection{Row 118: 12.2 DiagonalFredholm}
\label{row:118}
\textbf{Exact row scope.} Actual finite-compression matrix determinants converge locally uniformly to the reduced denominator; the zero-indexed operator is trace class and every positive power has trace kappa. Specialized diagonal Fredholm theory, with no general operator determinant assumed.\par\smallskip
\textbf{Manuscript labels.} \code{thm:trace_class_model}\par\smallskip
{\small\code{Komorebi.PaperIV.diagonalCompression}\par}
{\footnotesize\color{muted}definition; \code{KomorebiPaperIV/DiagonalFredholm.lean}, line 14.\par}\smallskip
{\small\code{Komorebi.PaperIV.diagonalCompression_eq_diagonal}\par}
{\footnotesize\color{muted}theorem; \code{KomorebiPaperIV/DiagonalFredholm.lean}, line 17.\par}\smallskip
{\small\code{Komorebi.PaperIV.diagonalCompression_determinant}\par}
{\footnotesize\color{muted}theorem; \code{KomorebiPaperIV/DiagonalFredholm.lean}, line 28.\par}\smallskip
{\small\code{Komorebi.PaperIV.diagonalFredholmDeterminant}\par}
{\footnotesize\color{muted}definition; \code{KomorebiPaperIV/DiagonalFredholm.lean}, line 40.\par}\smallskip
{\small\code{Komorebi.PaperIV.diagonalFredholmDeterminant_compression_limit}\par}
{\footnotesize\color{muted}theorem; \code{KomorebiPaperIV/DiagonalFredholm.lean}, line 42.\par}\smallskip
{\small\code{Komorebi.PaperIV.reducedZeroEigenvalue}\par}
{\footnotesize\color{muted}definition; \code{KomorebiPaperIV/DiagonalFredholm.lean}, line 50.\par}\smallskip
{\small\code{Komorebi.PaperIV.reducedZeroEigenvalue_norm_summable}\par}
{\footnotesize\color{muted}theorem; \code{KomorebiPaperIV/DiagonalFredholm.lean}, line 53.\par}\smallskip
{\small\code{Komorebi.PaperIV.reducedZeroOperator}\par}
{\footnotesize\color{muted}definition; \code{KomorebiPaperIV/DiagonalFredholm.lean}, line 58.\par}\smallskip
{\small\code{Komorebi.PaperIV.reducedZeroOperator_traceClass}\par}
{\footnotesize\color{muted}theorem; \code{KomorebiPaperIV/DiagonalFredholm.lean}, line 60.\par}\smallskip
{\small\code{Komorebi.PaperIV.reducedZeroOperator_basis}\par}
{\footnotesize\color{muted}theorem; \code{KomorebiPaperIV/DiagonalFredholm.lean}, line 64.\par}\smallskip
{\small\code{Komorebi.PaperIV.reducedZeroOperator_fredholm}\par}
{\footnotesize\color{muted}theorem; \code{KomorebiPaperIV/DiagonalFredholm.lean}, line 70.\par}\smallskip
{\small\code{Komorebi.PaperIV.reducedZeroOperator_trace_power}\par}
{\footnotesize\color{muted}theorem; \code{KomorebiPaperIV/DiagonalFredholm.lean}, line 77.\par}\smallskip
{\small\code{Komorebi.PaperIV.reducedZeroOperator_finite_determinants_uniform}\par}
{\footnotesize\color{muted}theorem; \code{KomorebiPaperIV/DiagonalFredholm.lean}, line 87.\par}\smallskip
\Needspace{7\baselineskip}
\subsection{Row 119: 12.3 CompleteHomogeneous}
\label{row:119}
\textbf{Exact row scope.} Finite complete homogeneous functions are actual monomial sums indexed by weak compositions; their generating function and Newton recurrence are proved independently of reciprocal coefficients.\par\smallskip
\textbf{Manuscript labels.} \code{cor:complete_symmetric_form}\par\smallskip
{\small\code{Komorebi.PaperIV.powerSumLog}\par}
{\footnotesize\color{muted}definition; \code{KomorebiPaperIV/CompleteHomogeneous.lean}, line 11.\par}\smallskip
{\small\code{Komorebi.PaperIV.powerSumLog_constantCoeff}\par}
{\footnotesize\color{muted}theorem; \code{KomorebiPaperIV/CompleteHomogeneous.lean}, line 14.\par}\smallskip
{\small\code{Komorebi.PaperIV.powerSumLog_derivative_coeff}\par}
{\footnotesize\color{muted}theorem; \code{KomorebiPaperIV/CompleteHomogeneous.lean}, line 17.\par}\smallskip
{\small\code{Komorebi.PaperIV.homogeneousSeries}\par}
{\footnotesize\color{muted}definition; \code{KomorebiPaperIV/CompleteHomogeneous.lean}, line 22.\par}\smallskip
{\small\code{Komorebi.PaperIV.homogeneousCoeff}\par}
{\footnotesize\color{muted}definition; \code{KomorebiPaperIV/CompleteHomogeneous.lean}, line 24.\par}\smallskip
{\small\code{Komorebi.PaperIV.homogeneousSeries_constantCoeff}\par}
{\footnotesize\color{muted}theorem; \code{KomorebiPaperIV/CompleteHomogeneous.lean}, line 26.\par}\smallskip
{\small\code{Komorebi.PaperIV.homogeneousSeries_derivative}\par}
{\footnotesize\color{muted}theorem; \code{KomorebiPaperIV/CompleteHomogeneous.lean}, line 29.\par}\smallskip
{\small\code{Komorebi.PaperIV.homogeneousCoeff_recurrence}\par}
{\footnotesize\color{muted}theorem; \code{KomorebiPaperIV/CompleteHomogeneous.lean}, line 34.\par}\smallskip
{\small\code{Komorebi.PaperIV.homogeneousGeometric}\par}
{\footnotesize\color{muted}definition; \code{KomorebiPaperIV/CompleteHomogeneous.lean}, line 42.\par}\smallskip
{\small\code{Komorebi.PaperIV.homogeneousGeometric_coeff}\par}
{\footnotesize\color{muted}theorem; \code{KomorebiPaperIV/CompleteHomogeneous.lean}, line 44.\par}\smallskip
{\small\code{Komorebi.PaperIV.homogeneousGeometric_constantCoeff}\par}
{\footnotesize\color{muted}theorem; \code{KomorebiPaperIV/CompleteHomogeneous.lean}, line 47.\par}\smallskip
{\small\code{Komorebi.PaperIV.homogeneousGeometric_inverse}\par}
{\footnotesize\color{muted}theorem; \code{KomorebiPaperIV/CompleteHomogeneous.lean}, line 50.\par}\smallskip
{\small\code{Komorebi.PaperIV.homogeneousGeometric_derivative}\par}
{\footnotesize\color{muted}theorem; \code{KomorebiPaperIV/CompleteHomogeneous.lean}, line 55.\par}\smallskip
{\small\code{Komorebi.PaperIV.finiteHomogeneousSeries}\par}
{\footnotesize\color{muted}definition; \code{KomorebiPaperIV/CompleteHomogeneous.lean}, line 72.\par}\smallskip
{\small\code{Komorebi.PaperIV.finiteHomogeneousSeries_constantCoeff}\par}
{\footnotesize\color{muted}theorem; \code{KomorebiPaperIV/CompleteHomogeneous.lean}, line 75.\par}\smallskip
{\small\code{Komorebi.PaperIV.finiteHomogeneousSeries_derivative}\par}
{\footnotesize\color{muted}theorem; \code{KomorebiPaperIV/CompleteHomogeneous.lean}, line 78.\par}\smallskip
{\small\code{Komorebi.PaperIV.finiteHomogeneousSeries_eq_homogeneous}\par}
{\footnotesize\color{muted}theorem; \code{KomorebiPaperIV/CompleteHomogeneous.lean}, line 88.\par}\smallskip
{\small\code{Komorebi.PaperIV.finiteHomogeneousSeries_monomials}\par}
{\footnotesize\color{muted}theorem; \code{KomorebiPaperIV/CompleteHomogeneous.lean}, line 100.\par}\smallskip
\Needspace{7\baselineskip}
\subsection{Row 120: 12.3 InfiniteSymmetric}
\label{row:120}
\textbf{Exact row scope.} Norm-summable eigenvalues give coefficientwise limits of the finite monomial sums. Repeating eigenvalues r times gives the exact positive-integer reciprocal coefficient; the complex-parameter Bell continuation is proved.\par\smallskip
\textbf{Manuscript labels.} \code{cor:complete_symmetric_form}\par\smallskip
{\small\code{Komorebi.PaperIV.summable_norm_positive_powers}\par}
{\footnotesize\color{muted}theorem; \code{KomorebiPaperIV/InfiniteSymmetric.lean}, line 13.\par}\smallskip
{\small\code{Komorebi.PaperIV.homogeneousCoeff_zero}\par}
{\footnotesize\color{muted}theorem; \code{KomorebiPaperIV/InfiniteSymmetric.lean}, line 21.\par}\smallskip
{\small\code{Komorebi.PaperIV.homogeneousCoeff_succ}\par}
{\footnotesize\color{muted}theorem; \code{KomorebiPaperIV/InfiniteSymmetric.lean}, line 24.\par}\smallskip
{\small\code{Komorebi.PaperIV.homogeneousCoeff_continuous}\par}
{\footnotesize\color{muted}theorem; \code{KomorebiPaperIV/InfiniteSymmetric.lean}, line 30.\par}\smallskip
{\small\code{Komorebi.PaperIV.homogeneousSeries_congr_positive}\par}
{\footnotesize\color{muted}theorem; \code{KomorebiPaperIV/InfiniteSymmetric.lean}, line 42.\par}\smallskip
{\small\code{Komorebi.PaperIV.completeHomogeneous}\par}
{\footnotesize\color{muted}definition; \code{KomorebiPaperIV/InfiniteSymmetric.lean}, line 52.\par}\smallskip
{\small\code{Komorebi.PaperIV.finiteHomogeneousSeries_coefficient_limit}\par}
{\footnotesize\color{muted}theorem; \code{KomorebiPaperIV/InfiniteSymmetric.lean}, line 55.\par}\smallskip
{\small\code{Komorebi.PaperIV.completeHomogeneous_monomial_limit}\par}
{\footnotesize\color{muted}theorem; \code{KomorebiPaperIV/InfiniteSymmetric.lean}, line 77.\par}\smallskip
{\small\code{Komorebi.PaperIV.repeatedEigenvalue}\par}
{\footnotesize\color{muted}definition; \code{KomorebiPaperIV/InfiniteSymmetric.lean}, line 83.\par}\smallskip
{\small\code{Komorebi.PaperIV.repeatedEigenvalue_norm_summable}\par}
{\footnotesize\color{muted}theorem; \code{KomorebiPaperIV/InfiniteSymmetric.lean}, line 85.\par}\smallskip
{\small\code{Komorebi.PaperIV.repeatedEigenvalue_powerSum}\par}
{\footnotesize\color{muted}theorem; \code{KomorebiPaperIV/InfiniteSymmetric.lean}, line 91.\par}\smallskip
{\small\code{Komorebi.PaperIV.reciprocalSeries_eq_homogeneous_zero_sums}\par}
{\footnotesize\color{muted}theorem; \code{KomorebiPaperIV/InfiniteSymmetric.lean}, line 98.\par}\smallskip
{\small\code{Komorebi.PaperIV.reciprocalCoeff_completeHomogeneous}\par}
{\footnotesize\color{muted}theorem; \code{KomorebiPaperIV/InfiniteSymmetric.lean}, line 111.\par}\smallskip
{\small\code{Komorebi.PaperIV.reciprocalCoeff_symmetric_polynomial_continuation}\par}
{\footnotesize\color{muted}theorem; \code{KomorebiPaperIV/InfiniteSymmetric.lean}, line 131.\par}\smallskip
\Needspace{7\baselineskip}
\subsection{Row 121: 12.4 NormalizedEntire}
\label{row:121}
\textbf{Exact row scope.} The literal normalized factorial series is entire and equals the residue transfer section after multiplication by z\textasciicircum{}j. Actual Taylor derivatives and value one at zero are proved.\par\smallskip
\textbf{Manuscript labels.} \code{prop:normalized_section_product}\par\smallskip
{\small\code{Komorebi.PaperIV.entireNormalizedSection}\par}
{\footnotesize\color{muted}definition; \code{KomorebiPaperIV/NormalizedEntire.lean}, line 13.\par}\smallskip
{\small\code{Komorebi.PaperIV.normalizedSection_coeff_norm}\par}
{\footnotesize\color{muted}theorem; \code{KomorebiPaperIV/NormalizedEntire.lean}, line 16.\par}\smallskip
{\small\code{Komorebi.PaperIV.normalizedSection_entire_radius}\par}
{\footnotesize\color{muted}theorem; \code{KomorebiPaperIV/NormalizedEntire.lean}, line 22.\par}\smallskip
{\small\code{Komorebi.PaperIV.analyticAt_entireNormalizedSection}\par}
{\footnotesize\color{muted}theorem; \code{KomorebiPaperIV/NormalizedEntire.lean}, line 38.\par}\smallskip
{\small\code{Komorebi.PaperIV.hasSum_entireNormalizedSection}\par}
{\footnotesize\color{muted}theorem; \code{KomorebiPaperIV/NormalizedEntire.lean}, line 43.\par}\smallskip
{\small\code{Komorebi.PaperIV.entireNormalizedSection_zero}\par}
{\footnotesize\color{muted}theorem; \code{KomorebiPaperIV/NormalizedEntire.lean}, line 51.\par}\smallskip
{\small\code{Komorebi.PaperIV.entireNormalizedSection_iteratedDeriv_zero}\par}
{\footnotesize\color{muted}theorem; \code{KomorebiPaperIV/NormalizedEntire.lean}, line 54.\par}\smallskip
{\small\code{Komorebi.PaperIV.entireNormalizedSection_transfer}\par}
{\footnotesize\color{muted}theorem; \code{KomorebiPaperIV/NormalizedEntire.lean}, line 67.\par}\smallskip
{\small\code{Komorebi.PaperIV.entireTransfer_unit_norm_le}\par}
{\footnotesize\color{muted}theorem; \code{KomorebiPaperIV/NormalizedEntire.lean}, line 93.\par}\smallskip
\Needspace{7\baselineskip}
\subsection{Row 122: 12.4 JensenZeroGrowth}
\label{row:122}
\textbf{Exact row scope.} Jensen bounds the actual multiplicity-weighted zero counts of a normalized entire function by its disk growth.\par\smallskip
\textbf{Manuscript labels.} \code{prop:normalized_section_product}\par\smallskip
{\small\code{Komorebi.PaperIV.finite_sum_le_finsum_nonneg}\par}
{\footnotesize\color{muted}theorem; \code{KomorebiPaperIV/JensenZeroGrowth.lean}, line 12.\par}\smallskip
{\small\code{Komorebi.PaperIV.entire_divisor_eq_orderNat}\par}
{\footnotesize\color{muted}theorem; \code{KomorebiPaperIV/JensenZeroGrowth.lean}, line 22.\par}\smallskip
{\small\code{Komorebi.PaperIV.entire_zero_finite_weight_growth}\par}
{\footnotesize\color{muted}theorem; \code{KomorebiPaperIV/JensenZeroGrowth.lean}, line 31.\par}\smallskip
{\small\code{Komorebi.PaperIV.entire_zero_subtype_weight_growth}\par}
{\footnotesize\color{muted}theorem; \code{KomorebiPaperIV/JensenZeroGrowth.lean}, line 62.\par}\smallskip
{\small\code{Komorebi.PaperIV.entire_zero_bounded_finite}\par}
{\footnotesize\color{muted}theorem; \code{KomorebiPaperIV/JensenZeroGrowth.lean}, line 74.\par}\smallskip
\Needspace{7\baselineskip}
\subsection{Row 123: 12.4 WeightedGrowthSummability}
\label{row:123}
\textbf{Exact row scope.} Multiplicity-weighted counting growth implies the exact subcritical power-sum summability bounds, by finite shell estimates.\par\smallskip
\textbf{Manuscript labels.} \code{prop:normalized_section_product}\par\smallskip
{\small\code{Komorebi.PaperIV.weighted_growth_summable_large}\par}
{\footnotesize\color{muted}theorem; \code{KomorebiPaperIV/WeightedGrowthSummability.lean}, line 14.\par}\smallskip
{\small\code{Komorebi.PaperIV.weighted_growth_summable}\par}
{\footnotesize\color{muted}theorem; \code{KomorebiPaperIV/WeightedGrowthSummability.lean}, line 61.\par}\smallskip
{\small\code{Komorebi.PaperIV.entire_growth_zeroPowerSummable}\par}
{\footnotesize\color{muted}theorem; \code{KomorebiPaperIV/WeightedGrowthSummability.lean}, line 79.\par}\smallskip
\Needspace{7\baselineskip}
\subsection{Row 124: 12.4 SublinearHadamard}
\label{row:124}
\textbf{Exact row scope.} General normalized entire functions with sublinear power growth equal the complete genus-zero product over their actual zeros with multiplicity. All summability and quotient-rigidity prerequisites are derived.\par\smallskip
\textbf{Manuscript labels.} \code{prop:normalized_section_product}\par\smallskip
{\small\code{Komorebi.PaperIV.entire_sublinear_growth_hadamard}\par}
{\footnotesize\color{muted}theorem; \code{KomorebiPaperIV/SublinearHadamard.lean}, line 15.\par}\smallskip
\Needspace{7\baselineskip}
\subsection{Row 125: 12.4 NormalizedGrowth}
\label{row:125}
\textbf{Exact row scope.} Global factorial/exponential bounds for the actual normalized sections imply every zero-power summability exponent above 1/N.\par\smallskip
\textbf{Manuscript labels.} \code{prop:normalized_section_product}\par\smallskip
{\small\code{Komorebi.PaperIV.entireNormalizedSection_norm_le}\par}
{\footnotesize\color{muted}theorem; \code{KomorebiPaperIV/NormalizedGrowth.lean}, line 13.\par}\smallskip
{\small\code{Komorebi.PaperIV.normalizedSectionGrowthConstant}\par}
{\footnotesize\color{muted}definition; \code{KomorebiPaperIV/NormalizedGrowth.lean}, line 39.\par}\smallskip
{\small\code{Komorebi.PaperIV.normalizedSectionGrowthConstant_pos}\par}
{\footnotesize\color{muted}theorem; \code{KomorebiPaperIV/NormalizedGrowth.lean}, line 41.\par}\smallskip
{\small\code{Komorebi.PaperIV.entireNormalizedSection_disk_growth}\par}
{\footnotesize\color{muted}theorem; \code{KomorebiPaperIV/NormalizedGrowth.lean}, line 47.\par}\smallskip
{\small\code{Komorebi.PaperIV.entireNormalizedSection_zeroPowerSummable}\par}
{\footnotesize\color{muted}theorem; \code{KomorebiPaperIV/NormalizedGrowth.lean}, line 61.\par}\smallskip
\Needspace{7\baselineskip}
\subsection{Row 126: 12.4 NormalizedOrder}
\label{row:126}
\textbf{Exact row scope.} The finite exponential representation proves the dominant positive-real-ray limit for each transfer residue section.\par\smallskip
\textbf{Manuscript labels.} \code{prop:normalized_section_product}\par\smallskip
{\small\code{Komorebi.PaperIV.entireTransfer_positive_ray}\par}
{\footnotesize\color{muted}theorem; \code{KomorebiPaperIV/NormalizedOrder.lean}, line 14.\par}\smallskip
\Needspace{7\baselineskip}
\subsection{Row 127: 12.4 NormalizedExactOrder}
\label{row:127}
\textbf{Exact row scope.} Matching upper and polynomially corrected lower bounds prove the exact order 1/N, including N=1. The residue condition is j<N.\par\smallskip
\textbf{Manuscript labels.} \code{prop:normalized_section_product}\par\smallskip
{\small\code{Komorebi.PaperIV.log_growth_ratio_of_polynomial_exp_bounds}\par}
{\footnotesize\color{muted}theorem; \code{KomorebiPaperIV/NormalizedExactOrder.lean}, line 13.\par}\smallskip
{\small\code{Komorebi.PaperIV.entireTransfer_positive_ray_lower}\par}
{\footnotesize\color{muted}theorem; \code{KomorebiPaperIV/NormalizedExactOrder.lean}, line 49.\par}\smallskip
{\small\code{Komorebi.PaperIV.entireNormalizedSection_maximum_bounds}\par}
{\footnotesize\color{muted}theorem; \code{KomorebiPaperIV/NormalizedExactOrder.lean}, line 66.\par}\smallskip
{\small\code{Komorebi.PaperIV.entireNormalizedSection_growth_limit}\par}
{\footnotesize\color{muted}theorem; \code{KomorebiPaperIV/NormalizedExactOrder.lean}, line 103.\par}\smallskip
{\small\code{Komorebi.PaperIV.entireNormalizedSection_order_exact}\par}
{\footnotesize\color{muted}theorem; \code{KomorebiPaperIV/NormalizedExactOrder.lean}, line 110.\par}\smallskip
\Needspace{7\baselineskip}
\subsection{Row 128: 12.4 NormalizedProduct}
\label{row:128}
\textbf{Exact row scope.} For N>=2, each normalized section equals its full multiplicity-indexed product, uniformly on compact sets. The actual analytic Taylor series is the named formal series.\par\smallskip
\textbf{Manuscript labels.} \code{prop:normalized_section_product}\par\smallskip
{\small\code{Komorebi.PaperIV.entireNormalizedSection_hadamard}\par}
{\footnotesize\color{muted}theorem; \code{KomorebiPaperIV/NormalizedProduct.lean}, line 16.\par}\smallskip
{\small\code{Komorebi.PaperIV.entireNormalizedSection_hadamard_uniform}\par}
{\footnotesize\color{muted}theorem; \code{KomorebiPaperIV/NormalizedProduct.lean}, line 26.\par}\smallskip
{\small\code{Komorebi.PaperIV.analyticTaylorSeries_entireNormalizedSection}\par}
{\footnotesize\color{muted}theorem; \code{KomorebiPaperIV/NormalizedProduct.lean}, line 35.\par}\smallskip
\Needspace{7\baselineskip}
\subsection{Row 129: 12.4 NormalizedZeroBell}
\label{row:129}
\textbf{Exact row scope.} The logarithmic coefficients equal the actual zero power sums, giving the exact (Nn)!/n! Bell normalization for the named residue coefficients and every complex parameter.\par\smallskip
\textbf{Manuscript labels.} \code{prop:normalized_section_product}\par\smallskip
{\small\code{Komorebi.PaperIV.normalized_zeroMomentSeries}\par}
{\footnotesize\color{muted}theorem; \code{KomorebiPaperIV/NormalizedZeroBell.lean}, line 11.\par}\smallskip
{\small\code{Komorebi.PaperIV.normalized_zeroPowerSum}\par}
{\footnotesize\color{muted}theorem; \code{KomorebiPaperIV/NormalizedZeroBell.lean}, line 29.\par}\smallskip
{\small\code{Komorebi.PaperIV.bellInput_normalized_zeros}\par}
{\footnotesize\color{muted}theorem; \code{KomorebiPaperIV/NormalizedZeroBell.lean}, line 35.\par}\smallskip
{\small\code{Komorebi.PaperIV.residuePowerCoeff_zero_sum_bell}\par}
{\footnotesize\color{muted}theorem; \code{KomorebiPaperIV/NormalizedZeroBell.lean}, line 52.\par}\smallskip
\Needspace{7\baselineskip}
\subsection{Row 130: 12.4 NormalizedAnalyticPowers}
\label{row:130}
\textbf{Exact row scope.} Actual principal-logarithm powers near zero have the named residue formal Taylor series and exact zero-moment Bell coefficients for every complex parameter. This is a germ at zero; global branch domains remain an earlier mapping boundary.\par\smallskip
\textbf{Manuscript labels.} \code{prop:normalized_section_product}\par\smallskip
{\small\code{Komorebi.PaperIV.powerSeries_eq_of_derivative_constant}\par}
{\footnotesize\color{muted}theorem; \code{KomorebiPaperIV/NormalizedAnalyticPowers.lean}, line 13.\par}\smallskip
{\small\code{Komorebi.PaperIV.analyticTaylorSeries_log}\par}
{\footnotesize\color{muted}theorem; \code{KomorebiPaperIV/NormalizedAnalyticPowers.lean}, line 23.\par}\smallskip
{\small\code{Komorebi.PaperIV.analyticTaylorSeries_exp}\par}
{\footnotesize\color{muted}theorem; \code{KomorebiPaperIV/NormalizedAnalyticPowers.lean}, line 46.\par}\smallskip
{\small\code{Komorebi.PaperIV.normalizedAnalyticPower}\par}
{\footnotesize\color{muted}definition; \code{KomorebiPaperIV/NormalizedAnalyticPowers.lean}, line 63.\par}\smallskip
{\small\code{Komorebi.PaperIV.normalizedAnalyticPower_taylor}\par}
{\footnotesize\color{muted}theorem; \code{KomorebiPaperIV/NormalizedAnalyticPowers.lean}, line 66.\par}\smallskip
{\small\code{Komorebi.PaperIV.normalizedAnalyticPower_zero_sum_bell}\par}
{\footnotesize\color{muted}theorem; \code{KomorebiPaperIV/NormalizedAnalyticPowers.lean}, line 82.\par}\smallskip
\Needspace{7\baselineskip}
\subsection{Row 131: 13.1 AnalyticTaylorAlgebra}
\label{row:131}
\textbf{Exact row scope.} Actual iterated-derivative Taylor coefficients preserve analytic germ equality, sums, products, natural powers and inverse products.\par\smallskip
\textbf{Manuscript labels.} \code{thm:nearest_circle_expansion}\par\smallskip
{\small\code{Komorebi.PaperIV.analyticTaylorSeries}\par}
{\footnotesize\color{muted}definition; \code{KomorebiPaperIV/AnalyticTaylorAlgebra.lean}, line 14.\par}\smallskip
{\small\code{Komorebi.PaperIV.analyticTaylorSeries_coeff}\par}
{\footnotesize\color{muted}theorem; \code{KomorebiPaperIV/AnalyticTaylorAlgebra.lean}, line 16.\par}\smallskip
{\small\code{Komorebi.PaperIV.analyticTaylorSeries_constantCoeff}\par}
{\footnotesize\color{muted}theorem; \code{KomorebiPaperIV/AnalyticTaylorAlgebra.lean}, line 20.\par}\smallskip
{\small\code{Komorebi.PaperIV.analyticTaylorSeries_congr}\par}
{\footnotesize\color{muted}theorem; \code{KomorebiPaperIV/AnalyticTaylorAlgebra.lean}, line 23.\par}\smallskip
{\small\code{Komorebi.PaperIV.analyticTaylorSeries_mul}\par}
{\footnotesize\color{muted}theorem; \code{KomorebiPaperIV/AnalyticTaylorAlgebra.lean}, line 28.\par}\smallskip
{\small\code{Komorebi.PaperIV.analyticTaylorSeries_add}\par}
{\footnotesize\color{muted}theorem; \code{KomorebiPaperIV/AnalyticTaylorAlgebra.lean}, line 46.\par}\smallskip
{\small\code{Komorebi.PaperIV.analyticTaylorSeries_const}\par}
{\footnotesize\color{muted}theorem; \code{KomorebiPaperIV/AnalyticTaylorAlgebra.lean}, line 51.\par}\smallskip
{\small\code{Komorebi.PaperIV.analyticTaylorSeries_pow}\par}
{\footnotesize\color{muted}theorem; \code{KomorebiPaperIV/AnalyticTaylorAlgebra.lean}, line 58.\par}\smallskip
{\small\code{Komorebi.PaperIV.analyticTaylorSeries_inverse_product}\par}
{\footnotesize\color{muted}theorem; \code{KomorebiPaperIV/AnalyticTaylorAlgebra.lean}, line 64.\par}\smallskip
\Needspace{7\baselineskip}
\subsection{Row 132: 12.4 and 13.1 AnalyticTaylorDifferential}
\label{row:132}
\textbf{Exact row scope.} Actual Taylor coefficients respect differentiation and negation. Analytic Hadamard products yield the formal logarithmic zero-moment identity.\par\smallskip
\textbf{Manuscript labels.} \code{prop:normalized_section_product}, \code{thm:nearest_circle_expansion}\par\smallskip
{\small\code{Komorebi.PaperIV.analyticTaylorSeries_neg}\par}
{\footnotesize\color{muted}theorem; \code{KomorebiPaperIV/AnalyticTaylorDifferential.lean}, line 13.\par}\smallskip
{\small\code{Komorebi.PaperIV.analyticTaylorSeries_deriv}\par}
{\footnotesize\color{muted}theorem; \code{KomorebiPaperIV/AnalyticTaylorDifferential.lean}, line 18.\par}\smallskip
{\small\code{Komorebi.PaperIV.analyticTaylorSeries_zeroResolvent}\par}
{\footnotesize\color{muted}theorem; \code{KomorebiPaperIV/AnalyticTaylorDifferential.lean}, line 29.\par}\smallskip
{\small\code{Komorebi.PaperIV.hadamard_taylor_zeroMoment_identity}\par}
{\footnotesize\color{muted}theorem; \code{KomorebiPaperIV/AnalyticTaylorDifferential.lean}, line 38.\par}\smallskip
\Needspace{7\baselineskip}
\subsection{Row 133: 13.1 PoleCoefficientBounds}
\label{row:133}
\textbf{Exact row scope.} Cauchy's estimate supplies explicit Taylor coefficient bounds, Big-O bounds and exact factorial/block rescaling.\par\smallskip
\textbf{Manuscript labels.} \code{thm:nearest_circle_expansion}\par\smallskip
{\small\code{Komorebi.PaperIV.analyticTaylorCoefficient}\par}
{\footnotesize\color{muted}definition; \code{KomorebiPaperIV/PoleCoefficientBounds.lean}, line 13.\par}\smallskip
{\small\code{Komorebi.PaperIV.analyticTaylorCoefficient_norm_le}\par}
{\footnotesize\color{muted}theorem; \code{KomorebiPaperIV/PoleCoefficientBounds.lean}, line 17.\par}\smallskip
{\small\code{Komorebi.PaperIV.analyticTaylorCoefficient_isBigO}\par}
{\footnotesize\color{muted}theorem; \code{KomorebiPaperIV/PoleCoefficientBounds.lean}, line 28.\par}\smallskip
{\small\code{Komorebi.PaperIV.analyticTaylorCoefficient_block_error}\par}
{\footnotesize\color{muted}theorem; \code{KomorebiPaperIV/PoleCoefficientBounds.lean}, line 45.\par}\smallskip
\Needspace{7\baselineskip}
\subsection{Row 134: 13.1 and 13.2 PrincipalPartCoefficients}
\label{row:134}
\textbf{Exact row scope.} The normalized pole has the exact binomial Taylor coefficients. Lower pole orders satisfy a uniform polynomial coefficient bound.\par\smallskip
\textbf{Manuscript labels.} \code{cor:leading_dominant_zero}, \code{thm:nearest_circle_expansion}\par\smallskip
{\small\code{Komorebi.PaperIV.principalPartSeries}\par}
{\footnotesize\color{muted}definition; \code{KomorebiPaperIV/PrincipalPartCoefficients.lean}, line 13.\par}\smallskip
{\small\code{Komorebi.PaperIV.principalPartSeries_coefficient}\par}
{\footnotesize\color{muted}theorem; \code{KomorebiPaperIV/PrincipalPartCoefficients.lean}, line 16.\par}\smallskip
{\small\code{Komorebi.PaperIV.principalPartSeries_inverse}\par}
{\footnotesize\color{muted}theorem; \code{KomorebiPaperIV/PrincipalPartCoefficients.lean}, line 25.\par}\smallskip
{\small\code{Komorebi.PaperIV.lower_principal_part_choose_bound}\par}
{\footnotesize\color{muted}theorem; \code{KomorebiPaperIV/PrincipalPartCoefficients.lean}, line 33.\par}\smallskip
\Needspace{7\baselineskip}
\subsection{Row 135: 13.1 PoleRemoval}
\label{row:135}
\textbf{Exact row scope.} The manuscript's local O(1) principal-part conditions imply nonnegative meromorphic order. The constructed normal-form remainder is holomorphic throughout the required zero-free disk after all finite principal parts are removed.\par\smallskip
\textbf{Manuscript labels.} \code{thm:nearest_circle_expansion}\par\smallskip
{\small\code{Komorebi.PaperIV.poleTerm}\par}
{\footnotesize\color{muted}definition; \code{KomorebiPaperIV/PoleRemoval.lean}, line 14.\par}\smallskip
{\small\code{Komorebi.PaperIV.polePrincipalPart}\par}
{\footnotesize\color{muted}definition; \code{KomorebiPaperIV/PoleRemoval.lean}, line 16.\par}\smallskip
{\small\code{Komorebi.PaperIV.rawPoleRemainder}\par}
{\footnotesize\color{muted}definition; \code{KomorebiPaperIV/PoleRemoval.lean}, line 19.\par}\smallskip
{\small\code{Komorebi.PaperIV.poleTerm_analyticAt}\par}
{\footnotesize\color{muted}theorem; \code{KomorebiPaperIV/PoleRemoval.lean}, line 22.\par}\smallskip
{\small\code{Komorebi.PaperIV.poleTerm_meromorphicAt}\par}
{\footnotesize\color{muted}theorem; \code{KomorebiPaperIV/PoleRemoval.lean}, line 31.\par}\smallskip
{\small\code{Komorebi.PaperIV.polePrincipalPart_analyticAt}\par}
{\footnotesize\color{muted}theorem; \code{KomorebiPaperIV/PoleRemoval.lean}, line 35.\par}\smallskip
{\small\code{Komorebi.PaperIV.rawPoleRemainder_meromorphic}\par}
{\footnotesize\color{muted}theorem; \code{KomorebiPaperIV/PoleRemoval.lean}, line 39.\par}\smallskip
{\small\code{Komorebi.PaperIV.rawPoleRemainder_analyticAt}\par}
{\footnotesize\color{muted}theorem; \code{KomorebiPaperIV/PoleRemoval.lean}, line 47.\par}\smallskip
{\small\code{Komorebi.PaperIV.meromorphicOrderAt_nonneg_of_bounded}\par}
{\footnotesize\color{muted}theorem; \code{KomorebiPaperIV/PoleRemoval.lean}, line 56.\par}\smallskip
{\small\code{Komorebi.PaperIV.rawPoleRemainder_bounded_at_pole}\par}
{\footnotesize\color{muted}theorem; \code{KomorebiPaperIV/PoleRemoval.lean}, line 66.\par}\smallskip
{\small\code{Komorebi.PaperIV.removedPoleRemainder}\par}
{\footnotesize\color{muted}definition; \code{KomorebiPaperIV/PoleRemoval.lean}, line 87.\par}\smallskip
{\small\code{Komorebi.PaperIV.removedPoleRemainder_analyticOnNhd}\par}
{\footnotesize\color{muted}theorem; \code{KomorebiPaperIV/PoleRemoval.lean}, line 90.\par}\smallskip
{\small\code{Komorebi.PaperIV.removedPoleRemainder_eq_raw}\par}
{\footnotesize\color{muted}theorem; \code{KomorebiPaperIV/PoleRemoval.lean}, line 105.\par}\smallskip
\Needspace{7\baselineskip}
\subsection{Row 136: 13.1 PoleTaylor}
\label{row:136}
\textbf{Exact row scope.} Actual analytic principal parts have the stated binomial coefficients. Integer reciprocal powers of the signed denominator have exactly the manuscript's factorial-scaled P coefficients.\par\smallskip
\textbf{Manuscript labels.} \code{thm:nearest_circle_expansion}\par\smallskip
{\small\code{Komorebi.PaperIV.analyticTaylorSeries_sub}\par}
{\footnotesize\color{muted}theorem; \code{KomorebiPaperIV/PoleTaylor.lean}, line 13.\par}\smallskip
{\small\code{Komorebi.PaperIV.analyticTaylorSeries_finsetSum}\par}
{\footnotesize\color{muted}theorem; \code{KomorebiPaperIV/PoleTaylor.lean}, line 18.\par}\smallskip
{\small\code{Komorebi.PaperIV.analyticTaylorSeries_id}\par}
{\footnotesize\color{muted}theorem; \code{KomorebiPaperIV/PoleTaylor.lean}, line 29.\par}\smallskip
{\small\code{Komorebi.PaperIV.analyticTaylorSeries_poleTerm}\par}
{\footnotesize\color{muted}theorem; \code{KomorebiPaperIV/PoleTaylor.lean}, line 37.\par}\smallskip
{\small\code{Komorebi.PaperIV.analyticTaylorCoefficient_principalPart}\par}
{\footnotesize\color{muted}theorem; \code{KomorebiPaperIV/PoleTaylor.lean}, line 54.\par}\smallskip
{\small\code{Komorebi.PaperIV.analyticTaylorSeries_canonicalSignedDenominator}\par}
{\footnotesize\color{muted}theorem; \code{KomorebiPaperIV/PoleTaylor.lean}, line 70.\par}\smallskip
{\small\code{Komorebi.PaperIV.reciprocalSeries_nat}\par}
{\footnotesize\color{muted}theorem; \code{KomorebiPaperIV/PoleTaylor.lean}, line 79.\par}\smallskip
{\small\code{Komorebi.PaperIV.analyticTaylorSeries_reciprocal_signed}\par}
{\footnotesize\color{muted}theorem; \code{KomorebiPaperIV/PoleTaylor.lean}, line 84.\par}\smallskip
{\small\code{Komorebi.PaperIV.reciprocalCoeff_analyticTaylor}\par}
{\footnotesize\color{muted}theorem; \code{KomorebiPaperIV/PoleTaylor.lean}, line 102.\par}\smallskip
\Needspace{7\baselineskip}
\subsection{Row 137: 13.1 SimplePoleParts}
\label{row:137}
\textbf{Exact row scope.} Analytic Taylor division constructs the full principal part at every simple zero and proves its leading Laurent coefficient and bounded remainder; no principal-part existence certificate is assumed.\par\smallskip
\textbf{Manuscript labels.} \code{thm:nearest_circle_expansion}\par\smallskip
{\small\code{Komorebi.PaperIV.poleNumerator}\par}
{\footnotesize\color{muted}definition; \code{KomorebiPaperIV/SimplePoleParts.lean}, line 13.\par}\smallskip
{\small\code{Komorebi.PaperIV.canonicalPoleCoefficient}\par}
{\footnotesize\color{muted}definition; \code{KomorebiPaperIV/SimplePoleParts.lean}, line 16.\par}\smallskip
{\small\code{Komorebi.PaperIV.poleNumerator_analyticAt}\par}
{\footnotesize\color{muted}theorem; \code{KomorebiPaperIV/SimplePoleParts.lean}, line 19.\par}\smallskip
{\small\code{Komorebi.PaperIV.canonicalPoleCoefficient_leading}\par}
{\footnotesize\color{muted}theorem; \code{KomorebiPaperIV/SimplePoleParts.lean}, line 31.\par}\smallskip
{\small\code{Komorebi.PaperIV.poleNumerator_identity}\par}
{\footnotesize\color{muted}theorem; \code{KomorebiPaperIV/SimplePoleParts.lean}, line 35.\par}\smallskip
{\small\code{Komorebi.PaperIV.polePrincipalPart_scaled}\par}
{\footnotesize\color{muted}theorem; \code{KomorebiPaperIV/SimplePoleParts.lean}, line 50.\par}\smallskip
{\small\code{Komorebi.PaperIV.canonicalPoleCoefficient_local_bound}\par}
{\footnotesize\color{muted}theorem; \code{KomorebiPaperIV/SimplePoleParts.lean}, line 73.\par}\smallskip
\Needspace{7\baselineskip}
\subsection{Row 138: 13.1 PrincipalPartUniqueness}
\label{row:138}
\textbf{Exact row scope.} Every principal part satisfying the local O(1) condition has leading coefficient (-rho D'(rho))\textasciicircum{}(-r), independently of the chosen construction.\par\smallskip
\textbf{Manuscript labels.} \code{thm:nearest_circle_expansion}\par\smallskip
{\small\code{Komorebi.PaperIV.polePrincipalPart_scaled_general}\par}
{\footnotesize\color{muted}theorem; \code{KomorebiPaperIV/PrincipalPartUniqueness.lean}, line 12.\par}\smallskip
{\small\code{Komorebi.PaperIV.polePrincipalPart_scaled_limit}\par}
{\footnotesize\color{muted}theorem; \code{KomorebiPaperIV/PrincipalPartUniqueness.lean}, line 25.\par}\smallskip
{\small\code{Komorebi.PaperIV.principalPart_leading_coefficient}\par}
{\footnotesize\color{muted}theorem; \code{KomorebiPaperIV/PrincipalPartUniqueness.lean}, line 50.\par}\smallskip
\Needspace{7\baselineskip}
\subsection{Row 139: 13.1 NearestPoleExpansion}
\label{row:139}
\textbf{Exact row scope.} The actual reciprocal coefficients satisfy the complete nearest-pole expansion with an explicit Cauchy constant for every coefficient. Its hypotheses state the prescribed local principal parts and the absence of other zeros inside the next modulus.\par\smallskip
\textbf{Manuscript labels.} \code{thm:nearest_circle_expansion}\par\smallskip
{\small\code{Komorebi.PaperIV.removedPoleRemainder_taylor}\par}
{\footnotesize\color{muted}theorem; \code{KomorebiPaperIV/NearestPoleExpansion.lean}, line 12.\par}\smallskip
{\small\code{Komorebi.PaperIV.nearest_pole_taylor_bound}\par}
{\footnotesize\color{muted}theorem; \code{KomorebiPaperIV/NearestPoleExpansion.lean}, line 41.\par}\smallskip
{\small\code{Komorebi.PaperIV.reciprocalCoeff_nearest_pole_bound}\par}
{\footnotesize\color{muted}theorem; \code{KomorebiPaperIV/NearestPoleExpansion.lean}, line 62.\par}\smallskip
\Needspace{7\baselineskip}
\subsection{Row 140: 13.2 LeadingPoleBounds}
\label{row:140}
\textbf{Exact row scope.} The leading term has separately controlled polynomial/radial and Cauchy errors, with actual simple-zero principal parts constructed automatically. For r=1 the first error constant is exactly zero.\par\smallskip
\textbf{Manuscript labels.} \code{cor:leading_dominant_zero}\par\smallskip
{\small\code{Komorebi.PaperIV.principalCoefficientSum}\par}
{\footnotesize\color{muted}definition; \code{KomorebiPaperIV/LeadingPoleBounds.lean}, line 13.\par}\smallskip
{\small\code{Komorebi.PaperIV.leadingPoleCoefficient}\par}
{\footnotesize\color{muted}definition; \code{KomorebiPaperIV/LeadingPoleBounds.lean}, line 16.\par}\smallskip
{\small\code{Komorebi.PaperIV.lowerPrincipalConstant}\par}
{\footnotesize\color{muted}definition; \code{KomorebiPaperIV/LeadingPoleBounds.lean}, line 19.\par}\smallskip
{\small\code{Komorebi.PaperIV.lowerPrincipalConstant_nonneg}\par}
{\footnotesize\color{muted}theorem; \code{KomorebiPaperIV/LeadingPoleBounds.lean}, line 22.\par}\smallskip
{\small\code{Komorebi.PaperIV.lowerPrincipalConstant_one}\par}
{\footnotesize\color{muted}theorem; \code{KomorebiPaperIV/LeadingPoleBounds.lean}, line 25.\par}\smallskip
{\small\code{Komorebi.PaperIV.principalCoefficientSum_split}\par}
{\footnotesize\color{muted}theorem; \code{KomorebiPaperIV/LeadingPoleBounds.lean}, line 28.\par}\smallskip
{\small\code{Komorebi.PaperIV.principalCoefficientSum_leading_error}\par}
{\footnotesize\color{muted}theorem; \code{KomorebiPaperIV/LeadingPoleBounds.lean}, line 41.\par}\smallskip
{\small\code{Komorebi.PaperIV.principalCoefficientSum_one}\par}
{\footnotesize\color{muted}theorem; \code{KomorebiPaperIV/LeadingPoleBounds.lean}, line 67.\par}\smallskip
{\small\code{Komorebi.PaperIV.reciprocalCoeff_leading_bound}\par}
{\footnotesize\color{muted}theorem; \code{KomorebiPaperIV/LeadingPoleBounds.lean}, line 72.\par}\smallskip
\Needspace{7\baselineskip}
\subsection{Row 141: 13.3 OrbitPoleReduction}
\label{row:141}
\textbf{Exact row scope.} Derivative rotation and block exponent invariance give N identical contributions per orbit. An equivalence records the chosen representatives and all distinct rotations; the resulting main term retains both proved error bounds.\par\smallskip
\textbf{Manuscript labels.} \code{cor:orbit_reduced_asymptotic}\par\smallskip
{\small\code{Komorebi.PaperIV.leadingZeroWeight}\par}
{\footnotesize\color{muted}definition; \code{KomorebiPaperIV/OrbitPoleReduction.lean}, line 12.\par}\smallskip
{\small\code{Komorebi.PaperIV.leadingZeroWeight_rotation}\par}
{\footnotesize\color{muted}theorem; \code{KomorebiPaperIV/OrbitPoleReduction.lean}, line 15.\par}\smallskip
{\small\code{Komorebi.PaperIV.leadingPoleCoefficient_eq_zero_weights}\par}
{\footnotesize\color{muted}theorem; \code{KomorebiPaperIV/OrbitPoleReduction.lean}, line 27.\par}\smallskip
{\small\code{Komorebi.PaperIV.leading_zero_sum_orbit_reduction}\par}
{\footnotesize\color{muted}theorem; \code{KomorebiPaperIV/OrbitPoleReduction.lean}, line 34.\par}\smallskip
{\small\code{Komorebi.PaperIV.reciprocalCoeff_orbit_leading_bound}\par}
{\footnotesize\color{muted}theorem; \code{KomorebiPaperIV/OrbitPoleReduction.lean}, line 55.\par}\smallskip
\Needspace{7\baselineskip}
\subsection{Row 142: 13.4 SecantAsymptotic}
\label{row:142}
\textbf{Exact row scope.} The actual cosine zeros and derivatives give the quantitative coefficient error, exponentially vanishing relative error, ratio limit one, and asymptotic equivalence to 2(2n)!(2/pi)\textasciicircum{}(2n+1).\par\smallskip
\textbf{Manuscript labels.} \code{cor:secant_asymptotic}\par\smallskip
{\small\code{Komorebi.PaperIV.cosineNearestZeros}\par}
{\footnotesize\color{muted}definition; \code{KomorebiPaperIV/SecantAsymptotic.lean}, line 13.\par}\smallskip
{\small\code{Komorebi.PaperIV.canonicalSignedDenominator_two_cos}\par}
{\footnotesize\color{muted}theorem; \code{KomorebiPaperIV/SecantAsymptotic.lean}, line 15.\par}\smallskip
{\small\code{Komorebi.PaperIV.cosineNearestZeros_norm}\par}
{\footnotesize\color{muted}theorem; \code{KomorebiPaperIV/SecantAsymptotic.lean}, line 19.\par}\smallskip
{\small\code{Komorebi.PaperIV.cosineNearestZeros_root}\par}
{\footnotesize\color{muted}theorem; \code{KomorebiPaperIV/SecantAsymptotic.lean}, line 23.\par}\smallskip
{\small\code{Komorebi.PaperIV.cosineNearestZeros_simple}\par}
{\footnotesize\color{muted}theorem; \code{KomorebiPaperIV/SecantAsymptotic.lean}, line 28.\par}\smallskip
{\small\code{Komorebi.PaperIV.cosine_nearest_zero_exhaustion}\par}
{\footnotesize\color{muted}theorem; \code{KomorebiPaperIV/SecantAsymptotic.lean}, line 33.\par}\smallskip
{\small\code{Komorebi.PaperIV.secant_leading_coefficient}\par}
{\footnotesize\color{muted}theorem; \code{KomorebiPaperIV/SecantAsymptotic.lean}, line 49.\par}\smallskip
{\small\code{Komorebi.PaperIV.secant_coefficient_error}\par}
{\footnotesize\color{muted}theorem; \code{KomorebiPaperIV/SecantAsymptotic.lean}, line 72.\par}\smallskip
{\small\code{Komorebi.PaperIV.secantMainTerm}\par}
{\footnotesize\color{muted}definition; \code{KomorebiPaperIV/SecantAsymptotic.lean}, line 87.\par}\smallskip
{\small\code{Komorebi.PaperIV.secantMainTerm_norm}\par}
{\footnotesize\color{muted}theorem; \code{KomorebiPaperIV/SecantAsymptotic.lean}, line 90.\par}\smallskip
{\small\code{Komorebi.PaperIV.secantMainTerm_ne_zero}\par}
{\footnotesize\color{muted}theorem; \code{KomorebiPaperIV/SecantAsymptotic.lean}, line 95.\par}\smallskip
{\small\code{Komorebi.PaperIV.secant_relative_error}\par}
{\footnotesize\color{muted}theorem; \code{KomorebiPaperIV/SecantAsymptotic.lean}, line 100.\par}\smallskip
{\small\code{Komorebi.PaperIV.secant_coefficient_ratio_limit}\par}
{\footnotesize\color{muted}theorem; \code{KomorebiPaperIV/SecantAsymptotic.lean}, line 120.\par}\smallskip
{\small\code{Komorebi.PaperIV.secant_coefficient_asymptotic}\par}
{\footnotesize\color{muted}theorem; \code{KomorebiPaperIV/SecantAsymptotic.lean}, line 131.\par}\smallskip
\Needspace{7\baselineskip}
\subsection{Row 143: 3.1, 3.2, 3.3 PointwiseFilters}
\label{row:143}
\textbf{Exact row scope.} Actual pointwise finite averages, coefficient-series HasSum, orthogonal projector resolution, holomorphic range/covariance equivalence on invariant domains, and the complete cyclic product law.\par\smallskip
\textbf{Manuscript labels.} \code{prop:covariance}, \code{thm:cyclic_product}, \code{thm:root_filter}\par\smallskip
{\small\code{Komorebi.PaperIV.pointwiseRootAverage}\par}
{\footnotesize\color{muted}definition; \code{KomorebiPaperIV/PointwiseFilters.lean}, line 14.\par}\smallskip
{\small\code{Komorebi.PaperIV.root_weight_sum}\par}
{\footnotesize\color{muted}theorem; \code{KomorebiPaperIV/PointwiseFilters.lean}, line 17.\par}\smallskip
{\small\code{Komorebi.PaperIV.mapsTo_root_pow}\par}
{\footnotesize\color{muted}theorem; \code{KomorebiPaperIV/PointwiseFilters.lean}, line 43.\par}\smallskip
{\small\code{Komorebi.PaperIV.pointwiseRootAverage_covariant}\par}
{\footnotesize\color{muted}theorem; \code{KomorebiPaperIV/PointwiseFilters.lean}, line 52.\par}\smallskip
{\small\code{Komorebi.PaperIV.pointwiseRootAverage_of_covariant}\par}
{\footnotesize\color{muted}theorem; \code{KomorebiPaperIV/PointwiseFilters.lean}, line 82.\par}\smallskip
{\small\code{Komorebi.PaperIV.pointwiseRootAverage_projector}\par}
{\footnotesize\color{muted}theorem; \code{KomorebiPaperIV/PointwiseFilters.lean}, line 99.\par}\smallskip
{\small\code{Komorebi.PaperIV.pointwiseRootAverage_resolution}\par}
{\footnotesize\color{muted}theorem; \code{KomorebiPaperIV/PointwiseFilters.lean}, line 106.\par}\smallskip
{\small\code{Komorebi.PaperIV.differentiableOn_pointwiseRootAverage}\par}
{\footnotesize\color{muted}theorem; \code{KomorebiPaperIV/PointwiseFilters.lean}, line 128.\par}\smallskip
{\small\code{Komorebi.PaperIV.pointwiseRootAverage_range_iff}\par}
{\footnotesize\color{muted}theorem; \code{KomorebiPaperIV/PointwiseFilters.lean}, line 138.\par}\smallskip
{\small\code{Komorebi.PaperIV.pointwiseRootAverage_sum}\par}
{\footnotesize\color{muted}theorem; \code{KomorebiPaperIV/PointwiseFilters.lean}, line 152.\par}\smallskip
{\small\code{Komorebi.PaperIV.pointwiseRootAverage_product}\par}
{\footnotesize\color{muted}theorem; \code{KomorebiPaperIV/PointwiseFilters.lean}, line 160.\par}\smallskip
{\small\code{Komorebi.PaperIV.pointwiseRootAverage_hasSum}\par}
{\footnotesize\color{muted}theorem; \code{KomorebiPaperIV/PointwiseFilters.lean}, line 185.\par}\smallskip
\Needspace{7\baselineskip}
\subsection{Row 144: 4.4, 5.5 PointwiseAddition}
\label{row:144}
\textbf{Exact row scope.} Actual entire transfer and signed addition formulas for arbitrary complex arguments; the signed carry factor is proved from rho\textasciicircum{}N=-1.\par\smallskip
\textbf{Manuscript labels.} \code{cor:exponential_addition}, \code{prop:signed_addition}\par\smallskip
{\small\code{Komorebi.PaperIV.pointwiseRootAverage_scaled_exp}\par}
{\footnotesize\color{muted}theorem; \code{KomorebiPaperIV/PointwiseAddition.lean}, line 10.\par}\smallskip
{\small\code{Komorebi.PaperIV.entireTransfer_addition}\par}
{\footnotesize\color{muted}theorem; \code{KomorebiPaperIV/PointwiseAddition.lean}, line 20.\par}\smallskip
{\small\code{Komorebi.PaperIV.signed_addition_weight}\par}
{\footnotesize\color{muted}theorem; \code{KomorebiPaperIV/PointwiseAddition.lean}, line 38.\par}\smallskip
{\small\code{Komorebi.PaperIV.entireSigned_addition}\par}
{\footnotesize\color{muted}theorem; \code{KomorebiPaperIV/PointwiseAddition.lean}, line 49.\par}\smallskip
\Needspace{7\baselineskip}
\subsection{Row 145: 7.4, 8.4, 8.5 AnalyticReciprocalLaws}
\label{row:145}
\textbf{Exact row scope.} Parameter addition and reciprocal differentiation on every open domain with a single holomorphic logarithm; the actual differential Bell/Riccati identity at every point where D\_N is nonzero. Formal coefficient clauses are covered by the earlier rows.\par\smallskip
\textbf{Manuscript labels.} \code{cor:reciprocal_derivative}, \code{thm:higher_riccati}, \code{thm:reciprocal_parameter_addition}\par\smallskip
{\small\code{Komorebi.PaperIV.IsLogarithmOn}\par}
{\footnotesize\color{muted}definition; \code{KomorebiPaperIV/AnalyticReciprocalLaws.lean}, line 14.\par}\smallskip
{\small\code{Komorebi.PaperIV.logarithmPower}\par}
{\footnotesize\color{muted}definition; \code{KomorebiPaperIV/AnalyticReciprocalLaws.lean}, line 18.\par}\smallskip
{\small\code{Komorebi.PaperIV.logarithmPower_add}\par}
{\footnotesize\color{muted}theorem; \code{KomorebiPaperIV/AnalyticReciprocalLaws.lean}, line 21.\par}\smallskip
{\small\code{Komorebi.PaperIV.IsLogarithmOn.nonzero}\par}
{\footnotesize\color{muted}theorem; \code{KomorebiPaperIV/AnalyticReciprocalLaws.lean}, line 28.\par}\smallskip
{\small\code{Komorebi.PaperIV.IsLogarithmOn.deriv}\par}
{\footnotesize\color{muted}theorem; \code{KomorebiPaperIV/AnalyticReciprocalLaws.lean}, line 33.\par}\smallskip
{\small\code{Komorebi.PaperIV.hasDerivAt_logarithmPower}\par}
{\footnotesize\color{muted}theorem; \code{KomorebiPaperIV/AnalyticReciprocalLaws.lean}, line 45.\par}\smallskip
{\small\code{Komorebi.PaperIV.reciprocal_fixed_log_derivative}\par}
{\footnotesize\color{muted}theorem; \code{KomorebiPaperIV/AnalyticReciprocalLaws.lean}, line 56.\par}\smallskip
{\small\code{Komorebi.PaperIV.analyticDifferentialBell}\par}
{\footnotesize\color{muted}definition; \code{KomorebiPaperIV/AnalyticReciprocalLaws.lean}, line 65.\par}\smallskip
{\small\code{Komorebi.PaperIV.analyticOnNhd_differentialBell}\par}
{\footnotesize\color{muted}theorem; \code{KomorebiPaperIV/AnalyticReciprocalLaws.lean}, line 69.\par}\smallskip
{\small\code{Komorebi.PaperIV.iteratedDeriv_eq_analyticDifferentialBell}\par}
{\footnotesize\color{muted}theorem; \code{KomorebiPaperIV/AnalyticReciprocalLaws.lean}, line 75.\par}\smallskip
{\small\code{Komorebi.PaperIV.analytic_higher_riccati}\par}
{\footnotesize\color{muted}theorem; \code{KomorebiPaperIV/AnalyticReciprocalLaws.lean}, line 93.\par}\smallskip
\Needspace{7\baselineskip}
\subsection{Row 146: 9.9 OperatorEncodingBridge}
\label{row:146}
\textbf{Exact row scope.} The derivative-truncation and factorial/Hasse encodings are proved equal on every polynomial. The named residue Appell polynomial has the exact forward action and the Phi(partial\_x\textasciicircum{}N)\textasciicircum{}r inverse, for arbitrary parameter and degree.\par\smallskip
\textbf{Manuscript labels.} \code{prop:appell_operator_form}\par\smallskip
{\small\code{Komorebi.PaperIV.truncatedDifferentialOperator_apply}\par}
{\footnotesize\color{muted}theorem; \code{KomorebiPaperIV/OperatorEncodingBridge.lean}, line 12.\par}\smallskip
{\small\code{Komorebi.PaperIV.truncatedDifferentialOperator_egfAction}\par}
{\footnotesize\color{muted}theorem; \code{KomorebiPaperIV/OperatorEncodingBridge.lean}, line 23.\par}\smallskip
{\small\code{Komorebi.PaperIV.residueAppell_operator_forward}\par}
{\footnotesize\color{muted}theorem; \code{KomorebiPaperIV/OperatorEncodingBridge.lean}, line 37.\par}\smallskip
{\small\code{Komorebi.PaperIV.residueAppell_operator_inverse}\par}
{\footnotesize\color{muted}theorem; \code{KomorebiPaperIV/OperatorEncodingBridge.lean}, line 48.\par}\smallskip
\Needspace{7\baselineskip}
\subsection{Row 147: 8.2, 8.3 TangentBoundary}
\label{row:147}
\textbf{Exact row scope.} D\_1=exp(-z), T\_1=1, and all order-one tangent coefficients are identified. Together with the prior N>=2 theorem, the original-variable tangent series converges with the exact factorial normalization for every positive order.\par\smallskip
\textbf{Manuscript labels.} \code{prop:tangent_kappa}, \code{thm:tangent_recurrence}\par\smallskip
{\small\code{Komorebi.PaperIV.canonicalSignedDenominator_one}\par}
{\footnotesize\color{muted}theorem; \code{KomorebiPaperIV/TangentBoundary.lean}, line 11.\par}\smallskip
{\small\code{Komorebi.PaperIV.analyticTransferTangent_one}\par}
{\footnotesize\color{muted}theorem; \code{KomorebiPaperIV/TangentBoundary.lean}, line 19.\par}\smallskip
{\small\code{Komorebi.PaperIV.transferTangentCoeff_one_order}\par}
{\footnotesize\color{muted}theorem; \code{KomorebiPaperIV/TangentBoundary.lean}, line 29.\par}\smallskip
{\small\code{Komorebi.PaperIV.analyticTransferTangent_one_hasSum}\par}
{\footnotesize\color{muted}theorem; \code{KomorebiPaperIV/TangentBoundary.lean}, line 47.\par}\smallskip
{\small\code{Komorebi.PaperIV.analyticTransferTangent_hasSum_all_orders}\par}
{\footnotesize\color{muted}theorem; \code{KomorebiPaperIV/TangentBoundary.lean}, line 58.\par}\smallskip
\Needspace{7\baselineskip}
\subsection{Row 148: 7.7 CompactReciprocal}
\label{row:148}
\textbf{Exact row scope.} The second and third reciprocal polynomials equal the literal compact binomial/factorial formulas in the manuscript, over every characteristic-zero field. The existing first-coefficient theorem completes all three displayed polynomials.\par\smallskip
\textbf{Manuscript labels.} \code{cor:first_reciprocal_polynomials}\par\smallskip
{\small\code{Komorebi.PaperIV.reciprocalCoeff_two_compact}\par}
{\footnotesize\color{muted}theorem; \code{KomorebiPaperIV/CompactReciprocal.lean}, line 32.\par}\smallskip
{\small\code{Komorebi.PaperIV.reciprocalCoeff_three_compact}\par}
{\footnotesize\color{muted}theorem; \code{KomorebiPaperIV/CompactReciprocal.lean}, line 44.\par}\smallskip
\Needspace{7\baselineskip}
\subsection{Row 149: 7.6 NewtonDeterminant}
\label{row:149}
\textbf{Exact row scope.} The explicit lower Hessenberg matrix has the displayed p\_k entries and negative integer superdiagonal. Its determinant equals n! b\_n from Newton's recurrence, including n=0; specialization gives the literal reciprocal determinant formula.\par\smallskip
\textbf{Manuscript labels.} \code{thm:logarithmic_newton_recurrence}\par\smallskip
{\small\code{Komorebi.PaperIV.newtonHessenberg}\par}
{\footnotesize\color{muted}definition; \code{KomorebiPaperIV/NewtonDeterminant.lean}, line 15.\par}\smallskip
{\small\code{Komorebi.PaperIV.newtonHessenberg_det}\par}
{\footnotesize\color{muted}theorem; \code{KomorebiPaperIV/NewtonDeterminant.lean}, line 59.\par}\smallskip
{\small\code{Komorebi.PaperIV.reciprocalOrdCoeff_newton_determinant}\par}
{\footnotesize\color{muted}theorem; \code{KomorebiPaperIV/NewtonDeterminant.lean}, line 112.\par}\smallskip
\Needspace{7\baselineskip}
\subsection{Row 150: 4.5 HypergeometricAnalytic}
\label{row:150}
\textbf{Exact row scope.} The literal N-1 denominator parameters exclude the unique parameter 1 and have positive real parts. Gamma normalization identifies the ordinary hypergeometric coefficients with the exact factorial coefficients. The everywhere convergent Mathlib regularized series yields the pointwise manuscript formula, including N=1.\par\smallskip
\textbf{Manuscript labels.} \code{thm:hypergeometric_form}\par\smallskip
{\small\code{Komorebi.PaperIV.residueHypergeometricParameters}\par}
{\footnotesize\color{muted}definition; \code{KomorebiPaperIV/HypergeometricAnalytic.lean}, line 14.\par}\smallskip
{\small\code{Komorebi.PaperIV.residueHypergeometricParameters_card}\par}
{\footnotesize\color{muted}theorem; \code{KomorebiPaperIV/HypergeometricAnalytic.lean}, line 17.\par}\smallskip
{\small\code{Komorebi.PaperIV.residueHypergeometricParameters_re_pos}\par}
{\footnotesize\color{muted}theorem; \code{KomorebiPaperIV/HypergeometricAnalytic.lean}, line 23.\par}\smallskip
{\small\code{Komorebi.PaperIV.regularized_zeroF_normalization}\par}
{\footnotesize\color{muted}theorem; \code{KomorebiPaperIV/HypergeometricAnalytic.lean}, line 50.\par}\smallskip
{\small\code{Komorebi.PaperIV.residueHypergeometric}\par}
{\footnotesize\color{muted}definition; \code{KomorebiPaperIV/HypergeometricAnalytic.lean}, line 71.\par}\smallskip
{\small\code{Komorebi.PaperIV.residueHypergeometric_hasSum}\par}
{\footnotesize\color{muted}theorem; \code{KomorebiPaperIV/HypergeometricAnalytic.lean}, line 75.\par}\smallskip
{\small\code{Komorebi.PaperIV.entireTransfer_hypergeometric}\par}
{\footnotesize\color{muted}theorem; \code{KomorebiPaperIV/HypergeometricAnalytic.lean}, line 99.\par}\smallskip
\Needspace{7\baselineskip}
\subsection{Row 151: 6.2 QuotientExponential}
\label{row:151}
\textbf{Exact row scope.} The quotient algebras receive their finite-dimensional coefficient norm topology. The ordinary exponential power series converges at every complex argument to the displayed entire-section coordinates in both x\textasciicircum{}N=1 and x\textasciicircum{}N=-1 algebras. The already checked derivative systems and addition laws retain their exact normalizations.\par\smallskip
\textbf{Manuscript labels.} \code{thm:finite_algebra_exponentials}\par\smallskip
{\small\code{Komorebi.PaperIV.finiteAlgebraExponential}\par}
{\footnotesize\color{muted}definition; \code{KomorebiPaperIV/QuotientExponential.lean}, line 13.\par}\smallskip
{\small\code{Komorebi.PaperIV.hasSum_cyclicAlgebra_exponential}\par}
{\footnotesize\color{muted}theorem; \code{KomorebiPaperIV/QuotientExponential.lean}, line 22.\par}\smallskip
{\small\code{Komorebi.PaperIV.hasSum_signedAlgebra_exponential}\par}
{\footnotesize\color{muted}theorem; \code{KomorebiPaperIV/QuotientExponential.lean}, line 45.\par}\smallskip
{\small\code{Komorebi.PaperIV.cyclicAlgebraNormedAddCommGroup}\par}
{\footnotesize\color{muted}definition; \code{KomorebiPaperIV/QuotientExponential.lean}, line 70.\par}\smallskip
{\small\code{Komorebi.PaperIV.cyclicAlgebraNormedSpace}\par}
{\footnotesize\color{muted}definition; \code{KomorebiPaperIV/QuotientExponential.lean}, line 76.\par}\smallskip
{\small\code{Komorebi.PaperIV.signedAlgebraNormedAddCommGroup}\par}
{\footnotesize\color{muted}definition; \code{KomorebiPaperIV/QuotientExponential.lean}, line 80.\par}\smallskip
{\small\code{Komorebi.PaperIV.signedAlgebraNormedSpace}\par}
{\footnotesize\color{muted}definition; \code{KomorebiPaperIV/QuotientExponential.lean}, line 86.\par}\smallskip
{\small\code{Komorebi.PaperIV.cyclicAlgebra_pointwise_exponential}\par}
{\footnotesize\color{muted}theorem; \code{KomorebiPaperIV/QuotientExponential.lean}, line 91.\par}\smallskip
{\small\code{Komorebi.PaperIV.signedAlgebra_pointwise_exponential}\par}
{\footnotesize\color{muted}theorem; \code{KomorebiPaperIV/QuotientExponential.lean}, line 99.\par}\smallskip
